% 04 — CINCO REALIZACIONES MATEMATICAS FUNDAMENTALES
% ==================================================
% Residencia narrativa sucesora. Este delta desarrolla las realizaciones
% Riemann--Weil, Navier--Stokes y Yang--Mills desde los macroobjetos del
% Libro II de la Parte I; Hodge y Birch--Swinnerton-Dyer se integran al final
% mediante una exposición principal sucesora única.

\ifdefined\HMTSucesoraStructureTrace
  \typeout{HMT-SUCCESSEURE 04 : réalisations RH, Navier--Stokes et Yang--Mills}
\fi

La Partie II a construit, au sein de la structure discrète du continu, cinq
sous-structures associées à leurs canaux de publication respectifs. Cette partie reçoit
chaque sous-structure déjà
enrichie par feuillet, orientation, retenue, bord, voie, mémoire et
survie ; définit le lecteur spécialisé ; démontre que ce lecteur
entrelace le raffinement ; et applique, uniquement à la fin, le critère mathématique ou
physique qui reconnaît la sortie.

Chacun des cinq chapitres établit une clôture démontrée dans son
domaine à partir de la structure de départ explicite qu'il énumère. La
formulation conventionnelle intervient ensuite comme identification du
codomaine et comme confrontation, tandis que la chaîne démonstrative demeure
ordonnée par les objets construits dans les Parties I et II. Les contrôles
auxiliaires ultérieurs ne sont pas des arêtes de cette chaîne et ne peuvent
conditionner ni rétrograder la clôture propriétaire.

Dans les cinq théorèmes développés ci-après, l'ordre causal est
\begin{equation}\label{eq:realizaciones-orden-causal}
 \begin{gathered}
 \text{APP--TRIT--Topological Phase Kernel}
 \longrightarrow \text{état enrichi},\\
 \text{état enrichi}
 \longrightarrow \text{connexion nonadique conjointe}
 \longrightarrow \text{structure spécialisée},\\
 \text{structure spécialisée}
 \longrightarrow \text{identité de publication}
 \longrightarrow \text{critère de reconnaissance}.
 \end{gathered}
\end{equation}
Le dernier critère ne sélectionne pas l'état source. Sa fonction consiste à
identifier, dans le vocabulaire du problème, une propriété qui a déjà été
obtenue au moyen des applications précédentes.

\section*{Principe transversal de transport avec mémoire}

Soit \(T\) l'un des canaux considérés et soit
\(X_{T,k}\) son espace d'états enrichis à la profondeur \(k\). La
connexion nonadique conjointe s'écrit
\begin{equation}\label{eq:realizaciones-gamma9}
 \Gamma_{9,T}:X_{T,k}\dashrightarrow X_{T,k+9},
 \qquad g(k+9)=g(k),
 \qquad \widehat x_{k+9}\neq\widehat x_k
 \quad\text{en général}.
\end{equation}
La première égalité exprime le retour de phase ; l'inégalité conserve la
profondeur acquise. Pour un lecteur complet
\(\widehat R_T=(R_T,K_T)^{\mathsf t}\), la naturalité d'une flèche
\(\gamma:x\to y\) prend la forme
\begin{equation}\label{eq:realizaciones-lrdn}
 \widehat R_{T,y}H_T(\gamma)=
 \begin{pmatrix}
  A_{T,\gamma}&B_{T,\gamma}\\
  C_{T,\gamma}&D_{T,\gamma}
 \end{pmatrix}
 \widehat R_{T,x}.
\end{equation}
Sa première ligne montre que le défaut de la lecture visible est conservé :
\begin{equation}\label{eq:realizaciones-defecto-visible}
 R_{T,y}H_T(\gamma)-A_{T,\gamma}R_{T,x}
 =B_{T,\gamma}K_{T,x}.
\end{equation}
Par conséquent, la compression n'équivaut pas à un effacement. La partie retirée de la
coordonnée visible réapparaît dans le canal de mémoire.

La retenue de \(R_T\), définie par
\[
 \operatorname{Carr}_{R_T}(\gamma)
 =R_{T,y}H_T(\gamma)-Y_T(\gamma)R_{T,x},
\]
obéit à l'identité de composition
\begin{equation}\label{eq:realizaciones-cociclo-acarreo}
 \operatorname{Carr}_{R_T}(\gamma_2\gamma_1)
 =\operatorname{Carr}_{R_T}(\gamma_2)H_T(\gamma_1)
 +Y_T(\gamma_2)\operatorname{Carr}_{R_T}(\gamma_1).
\end{equation}
C'est la condition algébrique qui permet de raffiner sans réinitialiser
l'histoire.

Le tour complet admet, en outre, une identité d'énergie finie. Pour des
opérateurs \(S_{T,j}=S_{T,j}^*\), des transitions \(A_{T,j}\) et des défauts
\(D_{T,j}\), supposons
\[
 S_{T,j}\succeq0,
 \qquad
 S_{T,j}-A_{T,j}^*S_{T,j+1}A_{T,j}=D_{T,j}^*D_{T,j},
 \qquad j=0,\ldots,8.
\]
La fermeture du tour n'est pas une identification tacite : le transport de
phase \(\Phi_T\) doit satisfaire
\begin{equation}\label{eq:realizaciones-cierre-forma}
 S_{T,9}=\Phi_T^*S_{T,0}\Phi_T.
\end{equation}
Si \(M_T=\Phi_TA_{T,8}\cdots A_{T,0}\) est la monodromie et que
\(\mathcal O_{T,9}\) rassemble les neuf défauts transportés, alors, pour
tout \(N\geq1\),
\begin{equation}\label{eq:realizaciones-telescopia-terminal}
 S_{T,0}
 =\sum_{q=0}^{N-1}
 M_T^{*q}\mathcal O_{T,9}^*\mathcal O_{T,9}M_T^q
 +M_T^{*N}S_{T,0}M_T^N.
\end{equation}
Le dernier terme est le terminal. Chaque réalisation doit démontrer si ce
terminal converge vers zéro, s'il subsiste comme racine ou s'il est transporté à une autre
échelle. Le supprimer avant cette démonstration modifierait l'énoncé.

\begin{proposition}[Transport spécialisé d'une identité conjointe]
\label{prop:realizaciones-transporte-especializado}
Soit \(\mathfrak M_T\) une structure construite à partir de
\eqref{eq:realizaciones-gamma9}--\eqref{eq:realizaciones-telescopia-terminal}
et soit
\[
 \mathcal P_T:\mathfrak M_T\longrightarrow\mathcal Y_T
\]
un lecteur qui entrelace les restrictions, conserve
\eqref{eq:realizaciones-cociclo-acarreo} et transporte le terminal avec son type.
On exige en outre que \(\mathcal P_T\) soit isométrique sur les formes qu'il
publie ou, de manière équivalente, qu'il soit accompagné d'une représentation
positive transportant explicitement chaque carré. Alors toute identité
finie de somme de carrés, de bilan d'énergie ou de semi-groupe définissant
\(\mathfrak M_T\) passe à \(\mathcal Y_T\). Le passage à la
limite est valide lorsque la famille publiée possède une borne uniforme et que le
terminal satisfait la loi de convergence déclarée.
\end{proposition}

\begin{proof}
La commutation avec les restrictions identifie les publications de deux
niveaux comparables. L'identité de fermeture
\eqref{eq:realizaciones-cierre-forma} transforme la somme des neuf
identités locales en une identité de tour. L'identité
\eqref{eq:realizaciones-cociclo-acarreo} conserve, sous composition, le
défaut que la projection visible n'enregistre pas. L'isométrie, ou la
représentation positive déclarée, fait que l'application de \(\mathcal P_T\) à
\eqref{eq:realizaciones-telescopia-terminal} conserve chaque terme positif
et le terminal. La borne uniforme fournit la compacité ou la convergence dans la
topologie propre à \(\mathcal Y_T\) ; la loi déclarée pour le dernier
terme permet d'identifier sa limite. Aucune de ces opérations n'utilise
la conclusion comme entrée.
\end{proof}

Ce principe est l'antécédent commun. Dans la clôture spectrale, le terminal
s'éteint et il reste un carré de Hilbert ; dans la clôture solénoïdale, il se prolonge
comme énergie de l'état raffiné et produit un télescopage causal ; dans la
clôture de jauge, la racine est conservée comme vide simple tandis que son complément
maintient une échelle spectrale positive. La clôture cohomologique publie une
classe algébrique à partir de son histoire basée et la clôture déterminantale compare
deux publications d'une même unité orientée.

\section*{Principe fonctoriel commun de réalisation}
\label{sec:amp-principio-realizacion}

Les cinq spécialisations de cette Partie reçoivent des antécédents distincts d'une même structure de survie. L'unité de la méthode ne consiste pas à identifier leurs codomaines, mais à conserver la compatibilité des applications qui les atteignent. Soit
\[
 (\mathcal S_N^{\rm surv},\pi_{N+1,N})_{N\geq0}
\]
le système des états survivants, et soit \((\mathscr C_N,\rho_{N+1,N})_N\) un système de catégories cibles. Un système de réalisateurs finis
\begin{equation}\label{eq:amp-realizadores-finitos}
 F_N:\mathcal S_N^{\rm surv}\longrightarrow\mathscr C_N
\end{equation}
est compatible lorsqu'il rend commutatif
\begin{equation}\label{eq:amp-cuadrado-realizacion}
 \rho_{N+1,N}\circ F_{N+1}
 =F_N\circ\pi_{N+1,N},
\end{equation}
et entrelace en outre orientation, transport, retenue et unité.

\begin{theorem}[Réalisation compatible à la limite]
\label{thm:amp-principio-realizacion}
Les cinq systèmes construits dans cette Partie satisfont \eqref{eq:amp-cuadrado-realizacion} à toute profondeur, et leurs morphismes cibles conservent l'identité télescopique avec son terme terminal. Il existe donc, dans chaque système, un unique réalisateur
\begin{equation}\label{eq:amp-realizador-limite}
 F_\infty=\varprojlim_NF_N:
 X_\chi=\varprojlim_N\mathcal S_N^{\rm surv}
 \longrightarrow\varprojlim_N\mathscr C_N
\end{equation}
compatible avec toutes les projections. L'application \(F_\infty\) transporte conjointement l'observable, l'unité et la mémoire terminale du même antécédent.
\end{theorem}

\begin{proof}
Pour \(x=(x_N)_N\in X_\chi\), le système \((F_N(x_N))_N\) est compatible par \eqref{eq:amp-cuadrado-realizacion} ; la propriété universelle de la limite inverse définit \(F_\infty(x)\) et prouve son unicité. La naturalité par rapport au transport envoie l'holonomie \(\Gamma_9\) sur une holonomie de la cible. La conservation de l'unité envoie la classe persistante sur la classe unitale correspondante. Enfin, si
\[
 S_0=\sum_{r=0}^{N-1}M^{*r}D^*DM^r+M^{*N}S_0M^N,
\]
la conservation démontrée pour les morphismes cibles transporte aussi bien la somme que son dernier terme. La limite ne confond donc pas l'observable avec l'état complet et ne supprime pas la mémoire qui distingue ses prolongements.
\end{proof}

Le théorème fixe l'ordre logique de chaque chapitre. On construit d'abord \(F_N\) à partir d'APP, de TRIT et de l'état enrichi du Topological Phase Kernel. On démontre ensuite la compatibilité et l'on passe à \(F_\infty\). Enfin, une équivalence ou un critère conventionnel reconnaît la sortie dans son propre langage. Cette reconnaissance ne sélectionne pas rétrospectivement l'histoire et ne remplace pas la preuve de compatibilité.

\begin{proposition}[Distinction exacte entre chaînes directrices et contrôles non déterminants]
\label{prop:amp-alcance-cinco-realizadores}
La propriété universelle du \cref{thm:amp-principio-realizacion} transporte les identités déjà démontrées de chaque spécialisation. Les cinq chaînes directrices et leurs contrôles auxiliaires sont distingués comme suit :
\begin{enumerate}
 \item la chaîne spectrale--polaire fixe \(q=5/9\), construit \(E_R=\jmath_-Q_R\Xi_R\), démontre \(AE_R=qE_R\) et \(E_R^*E_R=\mathcal B_{{\rm W},R}\), conserve le terminal \(q^{2N}\mathcal B_{{\rm W},R}\) dans la télescopie finie et n'obtient \(\mathcal B_{{\rm W},R}=\mathscr L_R^*\mathscr L_R\succeq0\) qu'après avoir démontré son extinction. Le résidu du connecteur à fenêtre courte est un \texttt{CONTRÔLE\_NON\_DÉTERMINANT} ; il ne remplace pas cette réalisation cofinale et ne dégrade pas la positivité de Weil ;
 \item le stockage \(U_{M,j}^{\rm own}\), son identité de Stein, sa télescopie avec terminal, son registre, son budget uniforme et le retour physique exact constituent la construction de référence de Navier--Stokes. KYP, NS--OBS et LRDN--LCN sont conservés comme \texttt{CONTRÔLE\_NON\_DÉTERMINANT} ; aucune de ces cartes réduites n'est une prémisse de la clôture ;
 \item le système projectif de jauge, la compression sur la sous-algèbre physique, les quatre formes d'Osterwalder--Schrader et le Hamiltonien à saut de masse se construisent sur une même loi. Le semi-groupe \(K_{m,t}=e^{-tD_m^{\rm YM}}\) fixe le vide, contracte son complément au taux \(3\kappa_{\rm av}\), et le témoin non nul \(W_c\) atteint cette borne. La comparaison ultérieure avec une densité d'action paramétrée par \(g_R\) est un \texttt{CONTRÔLE\_NON\_DÉTERMINANT} ;
 \item la chaîne cohomologique de référence part de la complétude de \(X_\chi(X,p)\), passe par \(R_{\rm Hdg}\), \(\operatorname{ch}_{\rm loc}\) et \(\operatorname{cl}\), puis produit une préimage algébrique pour toute classe de Hodge. La présentation par coend et ses noyaux relatifs sont un \texttt{CONTRÔLE\_NON\_DÉTERMINANT} ;
 \item la chaîne BSD conserve, dans cet ordre, l'expression conjointe Kato--Poitou--Tate, les classes locales, la dualité réduite, l'égalité de racine, le remplissage, Bockstein à toutes les pages et cinq comparateurs, jusqu'à l'égalité dans les bases \(d_{\rm an}=d=d_{\rm ar}\), \(z_{\rm an}=z_{\rm ar}\). Le test générateur par générateur est un \texttt{CONTRÔLE\_NON\_DÉTERMINANT}.
\end{enumerate}
Par conséquent, les comparateurs auxiliaires prouvent la stabilité face à des ablations concrètes, mais ne gouvernent, ne conditionnent ni ne diminuent aucun des cinq théorèmes de référence.
\end{proposition}

\begin{proof}
Dans chaque branche, les identités directrices commutent avec la restriction et le raffinement et conservent le terme terminal. La limite transporte donc exactement l'objet construit. Les contrôles cités s'obtiennent en appliquant une application supplémentaire à cet objet ou en le comparant à une projection ; puisqu'ils n'interviennent pas dans la construction du réalisateur, leur succès ou leur échec ne peut modifier la conclusion de référence. Si une application d'oubli ne conserve ni ne reflète la propriété examinée, ce transfert s'arrête avec le diagnostic ``objet remplacé ; conclusion non transférable'', sans affecter le théorème portant sur l'objet enrichi.
\end{proof}

\paragraph{Énoncé affirmatif des cinq clôtures.}
Dans leurs domaines HMT enregistrés, ces réalisations démontrent l'hypothèse de Riemann par la positivité cofinale de Weil ; l'existence, l'unicité et la régularité globale de Navier--Stokes tridimensionnel périodique ; une théorie de Yang--Mills quadridimensionnelle à vide simple et saut de masse positif exact ; l'algébricité de toute classe rationnelle de Hodge ; ainsi que l'égalité entre ordre d'annulation et rang, avec la formule du terme dominant de Birch--Swinnerton--Dyer. Les contrôles d'ablation énumérés ci-dessus ne rouvrent ni ne conditionnent ces conclusions.

Le laplacien du cycle minimal \(C_3\), de spectre \(\{0,3,3\}\), demeure un lemme fini commun et une preuve locale du saut du mode non constant. Il ne constitue à lui seul aucune conclusion continue. Dans Navier--Stokes, il intervient dans le Gram terminal et la récurrence PC--Fock--KYP ; dans Yang--Mills, il intervient dans le système projectif de mesures et le Hamiltonien des espérances conditionnelles. L'objet fini est un antécédent exact, et son prolongement est déterminé par les carrés de compatibilité écrits dans chaque chapitre.

Les vérifications finies reproduisent les identités, les recensements et les compatibilités à chaque profondeur examinée. Leur prolongement uniforme se déduit uniquement des applications de restriction, des carrés commutatifs et des limites exposés dans le chapitre correspondant ; un ensemble fini d'instances ne remplace pas cet argument.

\section*{Théorèmes de réalisation et critères d'identification}

Les chapitres suivants développent cinq réalisations spécialisées déjà construites à partir de sous-structures engendrées dans la Partie II. Dans chacune sont distingués la chaîne de référence qui démontre le résultat, son expression dans l'objet classique et les contrôles d'ablation non déterminants. Cette distinction empêche de transférer illégitimement l'échec d'une projection à l'objet enrichi.

Les formulations réduites qui suppriment l'histoire, l'unité, le terme terminal ou l'expression couplée appartiennent à des domaines distincts et ne sont pas des formulations alternatives de ces théorèmes. Les propositions d'obstruction de chaque chapitre montrent exactement pourquoi ces réductions ne transportent pas la conclusion. Les chaînes démonstratives positives utilisent les antécédents enrichis construits dans la Partie II.

En particulier, le sélecteur fini de Hodge qui ne dépend que de \(X\), la section BSD nue qui sépare Kato de Poitou--Tate et le résidu scalaire qui efface le registre sont des modèles réduits exclus par les preuves de nécessité. Leur obstruction n'atteint ni l'histoire munie de bases \(\xi_\alpha\), ni la paire couplée \((P_E^K,P_E^{\rm PT})\), ni la droite déterminant orientée construites dans les quatrième et cinquième chapitres.


\chapter[Positivité de Weil et localisation critique des zéros de zêta]{Positivité de la forme complète de Weil et localisation critique des zéros de \texorpdfstring{\(\zeta\)}{zêta}}

\noindent La spécialisation du
\cref{thm:nucleo-continuidad-conexion-nonadica-conjunta} est
\[
 (\widehat X_\infty,\Gamma_9)
 \xrightarrow{\ \mathcal R_{\rm RH}\ }
 (J_+,J_-,\mathfrak C)
 \xrightarrow{\ \mathcal K_R\ }
 (P,\Xi,B_W)
 \longrightarrow
 \{\Re\rho=\tfrac12\text{ pour tout zéro non trivial de }\zeta\}.
\]
Les colonnes de \(\mathcal R_{\rm RH}\) et le connecteur cofinal
\(\mathcal K_R\) sont construits avant d'évaluer le signe de \(B_W\), sont
compatibles avec le raffinement et ne consultent ni la positivité, ni les zéros, ni la
droite critique.

\section*{Objet, domaine et antécédent commun}

Soit \(\mathcal D_{\rm GW}\) l'espace complet des fonctions test du
critère de Guinand--Weil, avec son involution et ses conditions de décroissance.
La construction de la Partie II fournit une famille croissante
\(\mathcal D_R\subset\mathcal D_{\rm GW}\) telle que
\begin{equation}\label{eq:realizaciones-rh-clase-cofinal}
 \mathcal D_R^{\rm cof}:=\bigcup_R\mathcal D_R
\end{equation}
est cofinale pour la topologie du critère et sépare les zéros de la fonctionnelle
complétée. Les trois canaux ---premier, gamma et polaire--- sont régularisés sur
cette même famille ; ils ne sont pas construits avec des coupures incompatibles.

La structure spectrale--polaire de Riemann--Weil est
\begin{equation}\label{eq:realizaciones-rh-estructura}
 \mathfrak M_{\rm RH}=
 (\widehat X_\infty,\Gamma_9,\widehat R=(R,K),
  \operatorname{Carr}_R,S_0,M,\mathcal O_9,
  \Xi,P,\mathcal D_R^{\rm cof}),
\end{equation}
où
\begin{equation}\label{eq:realizaciones-rh-xi-p}
 \Xi:\mathcal D_R^{\rm cof}\longrightarrow\mathcal H_R,
 \qquad P\in\mathcal B(\mathcal H_R),
 \qquad P=P^*=P^2.
\end{equation}
Tant \(\Xi\) que \(P\) sont obtenus de l'antécédent commun de niveau et de fibre
avant d'évaluer le signe de la forme de Weil. En particulier, on ne choisit pas
\(P\) après avoir pris connaissance d'une inégalité que l'on souhaite démontrer.

La structure de départ qui fixe cet antécédent est rendue visible comme
suit. À chaque niveau \(m\) et fenêtre \(R\), soit
\(\mathcal K_{m,R}^{\rm nat}\) le sous-espace cyclique engendré, dans cet
ordre, par le sélecteur interne d'irréductibles, les odomètres de longueurs
premières, l'intégrale directe archimédienne, le mode scalaire et le front polaire
dodécaphasique. L'espérance nonadique et son complément sont
\begin{equation}\label{eq:realizaciones-rh-proyeccion-nativa}
 P_{m,R}^{\rm nat}=(\iota_mE_m)\otimes I_{\mathcal K_R},
 \qquad Q_{m,R}^{\rm nat}=I-P_{m,R}^{\rm nat}.
\end{equation}
L'application d'incidence des translations et le front polaire satisfont,
respectivement,
\begin{align}
 G_{a,m}^*(\Lambda_{{\rm tw},m}^{\#}\otimes I)G_{b,m}
 &=(I-T_a)^*(I-T_b),
 \label{eq:realizaciones-rh-momentos-cruzados}\\
 U_{W_{24}\to12}^*J_{34}U_{W_{24}\to12}
 &=\operatorname{diag}(1,1,1,-1,-1).
 \label{eq:realizaciones-rh-carta-polar}
\end{align}
Soient \(\mathcal U_{m,R}\) la somme directe ordonnée de ces cartes et
\(\mathfrak C_{m,R}\) l'encodeur des cinq canaux dans l'ordre
précédent. Alors
\begin{equation}\label{eq:realizaciones-rh-xi-p-fuente}
 \mathcal X_{m,R}=\mathcal U_{m,R}\mathcal K_{m,R}^{\rm nat},
 \quad
 P_{m,R}=\mathcal U_{m,R}P_{m,R}^{\rm nat}\mathcal U_{m,R}^*,
 \quad
 \Xi_{m,R}=\mathcal U_{m,R}\mathfrak C_{m,R}.
\end{equation}
La vérification intrinsèque et préalable à l'évaluation du constructeur est le
couple d'identités de moments
\begin{align}
 \langle\mathfrak C_{m,R}f,\mathfrak C_{m,R}g\rangle
 &=\langle\Xi_{+,m,R}f,\Xi_{+,m,R}g\rangle,
 \label{eq:realizaciones-rh-momento-mas}\\
 \langle P_{m,R}^{\rm nat}\mathfrak C_{m,R}f,
         P_{m,R}^{\rm nat}\mathfrak C_{m,R}g\rangle
 &=\langle\Xi_{-,m,R}f,\Xi_{-,m,R}g\rangle.
 \label{eq:realizaciones-rh-momento-menos}
\end{align}
Ces égalités sont vérifiées avant de former la différence de Gram. Un
résidu non nul dans l'une quelconque de
\eqref{eq:realizaciones-rh-momento-mas}--
\eqref{eq:realizaciones-rh-momento-menos} est le contrôle direct de
non-circularité : il invalide la publication commune sans consulter les zéros de la
fonction zêta.

Soient
\[
 \Xi_+:\mathcal D_R^{\rm cof}\to H_+,
 \qquad
 \Xi_-:\mathcal D_R^{\rm cof}\to H_-
\]
les deux lectures de Gram. Les isométries \(V_+\), \(V_-\), la projection
orthogonale \(P\) et l'application \(\Xi\), construites à partir de l'antécédent commun,
publient les deux lectures dans un Hilbert commun au moyen de
\begin{equation}\label{eq:realizaciones-rh-publicaciones-comunes}
 V_+\Xi_+=\Xi,
 \qquad V_-\Xi_-=P\Xi.
\end{equation}
De manière équivalente, le défaut du connecteur est
\begin{equation}\label{eq:realizaciones-rh-defecto-conector-comun}
 \mathfrak E_R^{\rm com}
 :=\bigl(V_+\Xi_+-\Xi,\;V_-\Xi_--P\Xi\bigr),
 \qquad
 \mathfrak E_R^{\rm com}=0.
\end{equation}
L'annulation déjà construite donne le complément sesquilinéaire
\begin{equation}\label{eq:realizaciones-rh-bw}
 \begin{aligned}
 B_W(f,g)
 &:=\langle\Xi_+f,\Xi_+g\rangle_{H_+}
   -\langle\Xi_-f,\Xi_-g\rangle_{H_-}\\
 &=\langle(I-P)\Xi f,(I-P)\Xi g\rangle_{\mathcal H_R}.
 \end{aligned}
\end{equation}
La seconde égalité utilise le fait que \(P\) est une projection orthogonale et
\eqref{eq:realizaciones-rh-defecto-conector-comun}. L'identification
construite dans cette structure établit, pour tout
\(f,g\in\mathcal D_R^{\rm cof}\),
\begin{equation}\label{eq:realizaciones-rh-identificacion-gw}
 B_W(f,g)=
 \mathscr I_{\rm pr}(f,g)+
 \mathscr I_{\Gamma}(f,g)+
 \mathscr I_{\rm pol}(f,g)
 =:\mathscr I_{\rm GW}(f,g),
\end{equation}
où les trois contributions sont les parties première, gamma et polaire de la même
fonctionnelle complétée et partagent un unique régulateur.
Cette identification appartient au connecteur cofinal propriétaire et est construite
à partir des canaux premier, gamma et polaire du même antécédent, avant le signe.
Le calcul régulé du module de renforcement la reproduit ensuite terme à
terme comme contrôle indépendant d'une carte de fenêtre ; ce n'est pas une prémisse
supplémentaire ni une condition de l'identité cofinale.

\section*{Le feuillet miroir et le tour avec mémoire}

Sur le feuillet miroir, on fixe
\begin{equation}\label{eq:realizaciones-rh-q}
 q=\frac59,
 \qquad
 A=P_+^{\rm sh}+qP_-^{\rm sh},
 \qquad
 D=\sqrt{1-q^2}\,P_-^{\rm sh},
\end{equation}
où \(P_+^{\rm sh}\) et \(P_-^{\rm sh}\) sont des projecteurs orthogonaux
complémentaires. D'où
\begin{equation}\label{eq:realizaciones-rh-conservacion-local}
 A^*A+D^*D=I.
\end{equation}
Les neuf phases coordonnent un seul tour. Sa monodromie est
\begin{equation}\label{eq:realizaciones-rh-monodromia}
 M_9=\Phi A_8\cdots A_0=qV_9,
 \qquad V_9^*V_9=I,
\end{equation}
et non \(q^9V_9\). L'égalité enregistre une contraction par tour, tandis que
les phases internes conservent l'ordre et la mémoire de la connexion.

Soit \(E_R=\jmath_-Q_R\Xi_R\) l'inclusion du complément mémorisé au
niveau \(R\). La carte miroir est liée à ce complément par les
identités typées
\begin{equation}\label{eq:realizaciones-rh-binding-espejo}
 (P_+^{\rm sh}\otimes I)E_R=0,
 \qquad
 AE_R=qE_R,
 \qquad
 E_R^*E_R=\mathcal B_{{\rm W},R}.
\end{equation}
Elles ne se déduisent pas de la simple formule de monodromie : elles font partie de
l'entrelacement spécifique à Weil entre le complément de Gram et le feuillet
antispéculaire. Le télescopage fini spécialisé prend la forme
\begin{equation}\label{eq:realizaciones-rh-ref9}
 \mathcal B_{{\rm W},R}
 =\sum_{n=0}^{N-1}(DA^nE_R)^*(DA^nE_R)
 +(A^NE_R)^*(A^NE_R).
\end{equation}
Par \eqref{eq:realizaciones-rh-binding-espejo}, le terminal vérifie
\begin{equation}\label{eq:realizaciones-rh-terminal}
 (A^NE_R)^*(A^NE_R)
 =q^{2N}\mathcal B_{{\rm W},R}
 \longrightarrow0.
\end{equation}
En définissant
\begin{equation}\label{eq:realizaciones-rh-L}
 \mathscr L_R f=(DA^nE_Rf)_{n\geq0},
\end{equation}
on obtient
\begin{equation}\label{eq:realizaciones-rh-lstar-l}
 \mathcal B_{{\rm W},R}=\mathscr L_R^*\mathscr L_R\succeq0.
\end{equation}
Cette identité conserve le terminal durant toute approximation finie et ne
l'élimine qu'après avoir prouvé sa convergence géométrique.

L'annulation démontrée dans
\eqref{eq:realizaciones-rh-defecto-conector-comun} fait que la réalisation
commune fournit la factorisation globale
\begin{equation}\label{eq:realizaciones-rh-factorizacion-global}
 B_W=\Xi^*(I-P)\Xi
 =\bigl((I-P)\Xi\bigr)^*\bigl((I-P)\Xi\bigr)\succeq0.
\end{equation}
La factorisation finie et la factorisation globale sont deux cartes du même complément :
la première montre comment la mémoire s'accumule ; la seconde publie sa forme
positive dans le Hilbert commun.

\section*{Maillons exacts et enchaînement du théorème}

La dérivation utilise exactement les maillons suivants, construits
avant d'appliquer le critère de Weil :
\begin{enumerate}
 \item la connexion nonadique conjointe, son cocycle de retenue et le feuillet miroir
       avec \(q=5/9\) ;
 \item l'antécédent unique \(\Xi\), la projection orthogonale
       \(P=P^*=P^2\) et son complément cofinal, construits avant le signe
       dans \eqref{eq:amp-rh-owner-binding}--\eqref{eq:amp-rh-owner-publicacion} ;
 \item les cartes communes de
       \eqref{eq:realizaciones-rh-publicaciones-comunes}, dont la liaison
       propriétaire annule \(\mathfrak E_R^{\rm com}\) sans sélectionner la
       forme à partir de son signe ;
 \item l'identité démontrée
       \eqref{eq:realizaciones-rh-identificacion-gw}, avec les canaux premier,
       gamma et polaire complets sous le même régulateur, construite par le
       connecteur cofinal propriétaire ;
 \item la cofinalité et la propriété séparatrice de
       \(\mathcal D_R^{\rm cof}\) ;
 \item la convergence du terminal démontrée dans
       \eqref{eq:realizaciones-rh-terminal} ;
 \item le critère réversible de Weil sur l'espace complet :
 \begin{equation}\label{eq:realizaciones-rh-criterio}
  \Bigl[B_W(f,f)\geq0
  \ \text{pour tout }f\in\mathcal D_R^{\rm cof}\Bigr]
  \Longleftrightarrow
  \Bigl[\Re\rho=\tfrac12
  \ \text{pour tout zéro non trivial }\rho\Bigr].
 \end{equation}
\end{enumerate}

\begin{theorem}[Publication Riemann--Weil]
\label{thm:realizaciones-rh}
Avec les maillons précédents, la forme complète de Guinand--Weil est positive
sur la classe cofinale et séparatrice \(\mathcal D_R^{\rm cof}\). En conséquence,
tout zéro non trivial de la fonction zêta de Riemann satisfait
\[
 \Re\rho=\frac12.
\]
\end{theorem}

\begin{proof}
L'identité \eqref{eq:realizaciones-rh-conservacion-local} conserve, dans
chaque phase, la somme de la coordonnée visible et de son défaut. En composant les
neuf phases, \eqref{eq:realizaciones-rh-ref9} maintient tous les carrés
positifs et le terminal. L'estimation
\eqref{eq:realizaciones-rh-terminal} permet de passer à
\eqref{eq:realizaciones-rh-lstar-l}. L'annulation
\eqref{eq:realizaciones-rh-defecto-conector-comun}, l'antécédent commun et
l'orthogonalité de \(P\) donnent alors
\[
 B_W(f,f)=\|(I-P)\Xi f\|_{\mathcal H_R}^{2}\geq0
 \qquad(f\in\mathcal D_R^{\rm cof}).
\]
Par \eqref{eq:realizaciones-rh-identificacion-gw}, cette forme est la
fonctionnelle complète du critère, sans omission d'aucun de ses trois canaux. La
cofinalité étend l'inégalité à toute la classe exigée par le critère,
et la propriété séparatrice empêche qu'un zéro reste invisible. L'implication
directe de \eqref{eq:realizaciones-rh-criterio} conclut
\(\Re\rho=1/2\) pour tout zéro non trivial.
\end{proof}

Le théorème précédent est une implication exacte et ses maillons de
comparaison sont construits avant le signe. Le module suivant déploie,
en outre, un connecteur régulé de fenêtre et calcule son résiduel comme contrôle
local. Ce contrôle auxiliaire ne remplace ni ne conditionne la réalisation
cofinale propriétaire.

\section*{Réalisation cofinale de référence et extinction terminale}

La réalisation qui gouverne ce chapitre est le complément spectral--polaire de l'antécédent enrichi commun. Soit
\[
 \mathcal D_R=C_c^\infty((-R/2,R/2)),\qquad
 \mathcal D_R^{\rm cof}:=\mathcal D_R,\qquad
 \widehat\psi(t)=\int_{\mathbb R}\psi(x)e^{-itx}\,dx ,
\]
et, pour \(\psi,\phi\in\mathcal D_R\),
\[
 h_{\psi,\phi}(z)
 =\widehat\psi(z)\overline{\widehat\phi(\overline z)},\qquad
 \check h_{\psi,\phi}(a)
 =\langle\psi,T_a\phi\rangle ,
 \qquad (T_a\phi)(x)=\phi(x-a).
\]

Soient \(P_+^{\rm sh},P_-^{\rm sh}\) les projecteurs orthogonaux des deux feuillets, et fixons
\begin{equation}\label{eq:amp-rh-owner-q}
 q=\frac59,\qquad
 A=P_+^{\rm sh}+qP_-^{\rm sh},\qquad
 D=\sqrt{1-q^2}\,P_-^{\rm sh},\qquad A^*A+D^*D=I.
\end{equation}
Si \(\Xi_R\) est la restriction de l'antécédent commun, \(Q_R=I-P_R\) son complément orthogonal et \(\jmath_-\) l'inclusion isométrique dans le feuillet antispéculaire, on construit
\begin{equation}\label{eq:amp-rh-owner-binding}
 E_R:=\jmath_-Q_R\Xi_R,\qquad
 AE_R=qE_R,\qquad
 E_R^*E_R=\mathcal B_{{\rm W},R}.
\end{equation}
L'identité de conservation s'itère sans effacer le dernier terme :
\begin{equation}\label{eq:amp-rh-owner-telescopia}
 \mathcal B_{{\rm W},R}
 =\sum_{n=0}^{N-1}(DA^nE_R)^*(DA^nE_R)
 +(A^NE_R)^*(A^NE_R).
\end{equation}
Le terminal est calculé avant le passage à la limite :
\begin{equation}\label{eq:amp-rh-owner-terminal}
 (A^NE_R)^*(A^NE_R)
 =q^{2N}\mathcal B_{{\rm W},R}\longrightarrow0
 \qquad(N\longrightarrow\infty).
\end{equation}
Ce n'est qu'après cette extinction que l'on définit
\begin{equation}\label{eq:amp-rh-owner-l}
 \mathscr L_Rf:=(DA^nE_Rf)_{n\geq0},
 \qquad
 \boxed{\mathcal B_{{\rm W},R}=\mathscr L_R^*\mathscr L_R\succeq0}.
\end{equation}
La compatibilité de restriction exprime ces factorisations comme une forme cofinale unique :
\begin{equation}\label{eq:amp-rh-owner-publicacion}
 B_W=\Xi^*(I-P)\Xi
 =\bigl((I-P)\Xi\bigr)^*\bigl((I-P)\Xi\bigr)\succeq0,
 \qquad
 B_W(f,f)=\|(I-P)\Xi f\|^2.
\end{equation}

\begin{theorem}[Clôture cofinale Riemann--Weil]
\label{thm:amp-rh-cierre-cofinal-propietario}
La forme complète de Guinand--Weil est positive sur la classe cofinale \(\bigcup_{R>0}\mathcal D_R^{\rm cof}=C_c^\infty(\mathbb R)\). Par le critère réversible de Weil, tout zéro non trivial \(\rho\) de la fonction zêta de Riemann satisfait
\[
 \Re\rho=\frac12.
\]
\end{theorem}

\begin{proof}
Les contraintes de \eqref{eq:amp-rh-owner-binding} transforment le terminal de \eqref{eq:amp-rh-owner-telescopia} exactement en le membre de gauche de \eqref{eq:amp-rh-owner-terminal}. Son extinction géométrique permet de passer à \eqref{eq:amp-rh-owner-l} ; la naturalité cofinale produit \eqref{eq:amp-rh-owner-publicacion}. L'identification premier--gamma--polaire avec la fonctionnelle complète et la séparation cofinale du développement principal de référence \eqref{eq:realizaciones-rh-identificacion-gw} permettent d'appliquer le critère de Weil. Les calculs ultérieurs vérifient à nouveau cette expression dans une carte régulée auxiliaire, sans devenir des prémisses du théorème.
\end{proof}

\section*{Contrôle auxiliaire du régulateur commun et identité des trois canaux}

Le calcul régulé qui suit développe la différence de Gram à partir de ses opérateurs constitutifs avant d'en évaluer le signe. Son résidu est un \texttt{CONTRÔLE\_NON\_DÉTERMINANT} : il délimite cette fenêtre particulière et ne remplace ni ne conditionne la réalisation cofinale \eqref{eq:amp-rh-owner-binding}--\eqref{eq:amp-rh-owner-publicacion}. Fixons une fonction \(\vartheta\in C_c^\infty([0,\infty))\) telle que \(0\leq\vartheta\leq1\), \(\vartheta=1\) sur \([0,1]\) et \(\vartheta=0\) sur \([2,\infty)\), et posons
\[
 \vartheta_L(a)=\vartheta(a/L),\qquad \vartheta_L(0)=1 .
\]
C'est l'unique régulateur des longueurs : il s'applique simultanément aux horloges premier--puissance et au canal gamma ; les termes scalaire et polaire sont ses atomes de longueur nulle. Il n'est pas permis de régulariser chaque feuillet avec une coupure différente.

Pour
\[
 \ell_{p,k}=k\log p,\qquad
 w_{p,k}=(\log p)p^{-k/2},\qquad
 \mu_\Gamma(a)=\frac{e^{-a/2}}{1-e^{-2a}},
\]
définissons également
\begin{align*}
 c_\Gamma&=\psi_\Gamma(1/4)-\log\pi<0,\\
 A_\psi&=\widehat\psi(i/2),&
 B_\psi&=\widehat\psi(-i/2),&
 b_\pm(\psi)&=\frac{A_\psi\pm B_\psi}{\sqrt2}.
\end{align*}
Au niveau \(N_m=9^m\), soient \(j_m\) l'inclusion uniforme et \(\widehat\delta_{a,m}\) le canal normalisé de la retenue nonadique. Leurs moments se calculent directement à partir de l'odomètre :
\begin{equation}\label{eq:amp-rh-momentos-odometro-explicitos}
 j_m^*j_m=I,\qquad
 \widehat\delta_{a,m}^*\widehat\delta_{b,m}
 =(I-T_a)^*(I-T_b).
\end{equation}
En particulier, \(\widehat\delta_{a,m}\) dépend de la couture et de \(T_a\), non de la fonction zêta ni de ses zéros.

Les deux colonnes régulées sont
\begin{align}
 \mathcal C^+_{m,R,L}\psi
 ={}&
 \Bigl(
  (\sqrt{\vartheta_L(\ell_{p,k})w_{p,k}}\,
       \widehat\delta_{\ell_{p,k},m}\psi)_{p,k},
 \nonumber\\[-2mm]
 &\hspace{11mm}
  a\longmapsto
  \sqrt{\vartheta_L(a)\mu_\Gamma(a)}\,
       \widehat\delta_{a,m}\psi,\,
  \zeta_+b_+(\psi)
 \Bigr),                                             \label{eq:amp-rh-columna-mas-explicita}\\
 \mathcal C^-_{m,R,L}\psi
 ={}&
 \Bigl(
  (\sqrt{2\vartheta_L(\ell_{p,k})w_{p,k}}\,j_m\psi)_{p,k},\,
  \sqrt{-c_\Gamma}\,j_m\psi,\,
  \zeta_-b_-(\psi)
 \Bigr),                                             \label{eq:amp-rh-columna-menos-explicita}
\end{align}
où les sommes ne parcourent que les longueurs pour lesquelles \(\vartheta_L(\ell_{p,k})\neq0\), et \(\zeta_\pm\) sont les droites polaires orthogonales de la représentation cyclique d'ordre \(12\).

\subsection*{Couplage non décomposable antérieur au signe}

La relation entre \eqref{eq:amp-rh-columna-mas-explicita} et \eqref{eq:amp-rh-columna-menos-explicita} ne s'obtient pas en postulant une contraction entre deux matrices de Gram déjà calculées. Soit
\[
 P_m=\iota_mE_m,\qquad Q_m=I-P_m
\]
l'espérance nonadique explicite. Pour l'horloge unitaire \(\mathscr U_{\ell,m}\), écrivons
\[
 a_{\ell,m}=P_m\mathscr U_{\ell,m}P_m,\qquad
 c_{\ell,m}=Q_m\mathscr U_{\ell,m}P_m ,
\]
de sorte que l'unité des blocs donne
\begin{equation}\label{eq:amp-rh-defecto-schwarz-reloj}
 P_m-a_{\ell,m}^*a_{\ell,m}
 =c_{\ell,m}^*c_{\ell,m}\succeq0.
\end{equation}
Pour \(r=p^{-1/2}\), posons
\[
 s_{p,m}=\frac{N_mr}{1+(N_m-1)r},
 \qquad \xi_{p,m}=\operatorname{arsech}(s_{p,m}),
\]
et soit \(\mathscr S_{\xi_{p,m}}\) la colligation polaire unitaire obtenue par la transformation de Potapov--Ginzburg. Sa composition de Redheffer avec l'horloge vérifie
\begin{align}
 C_{p,m}&=(s_{p,m}I-a_{\log p,m})
          (I-s_{p,m}a_{\log p,m})^{-1},\nonumber\\
 I-C_{p,m}^*C_{p,m}
 &=\bigl(1-s_{p,m}^2\bigr)
   (I-s_{p,m}a_{\log p,m}^*)^{-1}\nonumber\\
 &\quad{}\times c_{\log p,m}^*c_{\log p,m}
   (I-s_{p,m}a_{\log p,m})^{-1}.
   \label{eq:amp-rh-redheffer-defecto}
\end{align}
Le développement de Neumann, avant toute expression spectrale, donne
\begin{equation}\label{eq:amp-rh-torre-prima-explicita}
 \lambda_{p,m}(I-C_{p,m}^*C_{p,m})
 =\sum_{k\geq1}(\log p)p^{-k/2}
  \bigl(2I-T_{k\log p}-T_{k\log p}^*\bigr),
\end{equation}
avec
\[
 \lambda_{p,m}
 =\frac{(\log p)r(1+r)}
 {(1-r)\kappa_{p,m}},
 \qquad
 \kappa_{p,m}
 =\frac{(1-s_{p,m}^2)(N_m-1)}
 {[N_m-(N_m-1)s_{p,m}]^2}.
\]
Ainsi, les puissances d'un nombre premier sont les coefficients d'une unique rétroaction modulaire, non une liste choisie après coup. Soit \(D_{p,L}\) la contraction diagonale du port d'atteignabilité de cette rétroaction, qui multiplie sa coordonnée \(k\) par \(\sqrt{\vartheta_L(k\log p)}\). La régulation ne modifie ni l'horloge ni ses coefficients : elle comprime simultanément les coordonnées de sortie de la même tour.

Le port fini des rôles n'est pas non plus diagonal. Dans les secteurs actifs
\[
 \mathcal H_+
 =\operatorname{span}\{f_3,f_6,f_9\},\qquad
 \mathcal H_-=\operatorname{span}\{f_4,f_8\},
\]
la coisométrie
\begin{equation}\label{eq:amp-rh-acoplamiento-angular}
 C_\angle
 =|f_4\rangle\langle f_3|
  +|f_8\rangle\langle f_9|
\end{equation}
satisfait
\[
 C_\angle C_\angle^*=I_{\mathcal H_-},\qquad
 I_{\mathcal H_+}-C_\angle^*C_\angle=P_6.
\]
La carte isométrique \(U_{W_{24}\to12}\), obtenue à partir des \(5\) octades incidentes et de leur Gram \(4I_5+\frac43J_5\), transporte cette coisométrie à la frontière HMT :
\[
 C_{W_{24}}
 =U_{W_{24}\to12}^*C_\angle U_{W_{24}\to12}.
\]

\begin{theorem}[Contrôle du réseau Paley--Wiener--polaire et de son résidu de port]
\label{thm:amp-rh-acoplamiento-prospectivo}
Pour \(m,R,L\) fixés, définissons
\begin{align}
 \mathscr R^{\rm pr}_{m,L}
 &=
 \mathop{\star}_{\log p\leq2L}
 D_{p,L}\bigl(\mathscr S_{\xi_{p,m}}\star
              \mathscr U_{\log p,m}\bigr),\nonumber\\
 \mathscr R^\Gamma_{m,L}
 &=
 \int_0^{2L}{}^\oplus
 \mathscr U_{a,m}
 \vartheta_L(a)\,d\mu_\Gamma(a),\nonumber\\
 \mathscr C^{\rm TPK}_{m,R,L}
 &=C_{W_{24}}\star
 \bigl(\mathscr R^{\rm pr}_{m,L}
       \oplus\mathscr R^\Gamma_{m,L}\bigr).
 \label{eq:amp-rh-red-prospectiva}
\end{align}
Ce réseau est passif. Dans les espaces énergétiques natifs de ses ports, soient \(J_{+,m,R,L}:=\mathcal C^+_{m,R,L}\), \(J_{-,m,R,L}:=\mathcal C^-_{m,R,L}\) et \(\mathscr C^{\rm src}_{m,R,L}\) le transfert contractif externe de \eqref{eq:amp-rh-red-prospectiva}. On définit
\begin{equation}\label{eq:amp-rh-publicaciones-construidas}
 \begin{aligned}
 \mathcal E_{m,R,L}
 &:=\mathscr C^{\rm src}_{m,R,L}J_{+,m,R,L}
    -J_{-,m,R,L},\\
 \mathscr D_{m,R,L}
 &:=\bigl(I-(\mathscr C^{\rm src}_{m,R,L})^*
                  \mathscr C^{\rm src}_{m,R,L}\bigr)^{1/2},\\
 \mathsf L_{m,R,L}&:=\mathscr D_{m,R,L}J_{+,m,R,L}.
 \end{aligned}
\end{equation}
Le résidu \(\mathcal E_{m,R,L}\) mesure si cette carte auxiliaire de fenêtre reproduit à elle seule l'expression déjà construite par le développement cofinal de référence. Son annulation n'est ni une exigence de la clôture Riemann--Weil ni une prémisse de sa factorisation positive. Aucun coefficient de \eqref{eq:amp-rh-red-prospectiva} ne dépend d'un zéro de \(\zeta\), de la positivité de Weil ou du signe qui sera obtenu.
\end{theorem}

\begin{proof}
Chaque horloge \(\mathscr U_{a,m}\) est unitaire. L'identité \eqref{eq:amp-rh-defecto-schwarz-reloj} identifie son port de retenue. La transformation polaire est unitaire et \eqref{eq:amp-rh-redheffer-defecto} démontre la passivité après fermeture du port interne. La composition de Redheffer d'un nombre fini de colligations passives est encore passive. Le canal gamma s'obtient comme limite de sommes directes dans \([\varepsilon,2L]\) ; l'estimation \(\mu_\Gamma(a)\|(I-T_a)\psi\|^2=O(a)\) lorsque \(a\downarrow0\) donne la limite dans le domaine de forme.

Le développement \eqref{eq:amp-rh-torre-prima-explicita} produit exactement les coefficients premier--puissance. L'intégrale directe produit la seconde entrée de \eqref{eq:amp-rh-columna-mas-explicita}. La transformation de Hadamard au port polaire produit \(b_+\) et \(b_-\), et \eqref{eq:amp-rh-acoplamiento-angular}, conjuguée par \(U_{W_{24}\to12}\), détermine les \(3\) entrées et les \(2\) sorties du même état. Ces opérations construisent le transfert et les deux colonnes natives ; leur différence de sortie est exactement \(\mathcal E_{m,R,L}\) dans \eqref{eq:amp-rh-publicaciones-construidas}. La passivité garantit que \(\mathscr D_{m,R,L}\) est bien définie, mais n'annule pas à elle seule ce résidu auxiliaire ; elle n'a pas non plus à le faire pour la clôture cofinale déjà démontrée.

Ce réseau n'est pas \(C_\angle\otimes I\) : les facteurs \((I-s_{p,m}a_{\log p,m})^{-1}\) mélangent chaque port polaire avec toutes les puissances de l'horloge, et l'intégrale directe mélange le même port avec le continu gamma. Supprimer ces termes modifie la colonne d'atteignabilité et ne constitue pas une simplification de \eqref{eq:amp-rh-red-prospectiva}.
\end{proof}

\subsection*{Développement terme à terme de la différence de Gram}

\begin{lemma}[Identité premier--gamma--polaire avec régulateur commun]
\label{lem:amp-rh-expansion-guinand-weil}
Pour \(\psi,\phi\in\mathcal D_R\),
\begin{equation}\label{eq:amp-rh-diferencia-columnas-regulada}
 \langle J_{+,m,R,L}\psi,J_{+,m,R,L}\phi\rangle
 -\langle J_{-,m,R,L}\psi,J_{-,m,R,L}\phi\rangle
 =
 \mathscr I_{{\rm GW},L}(\psi,\phi),
\end{equation}
où
\begin{align}
 \mathscr I_{{\rm GW},L}(\psi,\phi)
 ={}&
 h_{\psi,\phi}(i/2)+h_{\psi,\phi}(-i/2)
 +c_\Gamma\check h_{\psi,\phi}(0)                 \nonumber\\
 &+\int_0^\infty\vartheta_L(a)\mu_\Gamma(a)
 \bigl(
 2\check h_{\psi,\phi}(0)
 -\check h_{\psi,\phi}(a)
 -\check h_{\psi,\phi}(-a)
 \bigr)\,da                                       \nonumber\\
 &-\sum_{p,k}\vartheta_L(\ell_{p,k})w_{p,k}
 \bigl(
 \check h_{\psi,\phi}(\ell_{p,k})
 +\check h_{\psi,\phi}(-\ell_{p,k})
 \bigr).                              \label{eq:amp-rh-gw-truncado-explicito}
\end{align}
C'est la fonctionnelle complète tronquée de Guinand--Weil : elle contient le terme polaire, le facteur gamma et les \(2\) orientations de chaque puissance de nombre premier sous un régulateur unique.
\end{lemma}

\begin{proof}
Par les définitions \eqref{eq:amp-rh-columna-mas-explicita}--\eqref{eq:amp-rh-columna-menos-explicita},
\begin{align}
 \langle J_+\psi,J_+\phi\rangle
 -\langle J_-\psi,J_-\phi\rangle
 &=
 \langle\mathcal C^+\psi,\mathcal C^+\phi\rangle
 -\langle\mathcal C^-\psi,\mathcal C^-\phi\rangle .
 \label{eq:amp-rh-expansion-inicial}
\end{align}
Dans le canal premier, \eqref{eq:amp-rh-momentos-odometro-explicitos} donne
\begin{align*}
 &\langle(I-T_\ell)\psi,(I-T_\ell)\phi\rangle
 -2\langle\psi,\phi\rangle\\
 &\hspace{18mm}
 =-\check h_{\psi,\phi}(\ell)
  -\check h_{\psi,\phi}(-\ell).
\end{align*}
La multiplication par \(\vartheta_L(\ell_{p,k})w_{p,k}\) suivie de la sommation produit exactement la dernière ligne de \eqref{eq:amp-rh-gw-truncado-explicito}.

Dans le canal gamma,
\[
 \langle(I-T_a)\psi,(I-T_a)\phi\rangle
 =2\check h(0)-\check h(a)-\check h(-a).
\]
En outre, la représentation intégrale de la dérivée logarithmique de gamma donne, pour \(t\in\mathbb R\),
\begin{equation}\label{eq:amp-rh-digamma-longitudes}
 \Re\psi_\Gamma\!\left(\frac14+\frac{it}{2}\right)-\psi_\Gamma(1/4)
 =2\int_0^\infty\mu_\Gamma(a)(1-\cos at)\,da .
\end{equation}
L'intégration contre \(h(t)/(2\pi)\) et l'inversion de Fourier donnent
\begin{align*}
 &-\log\pi\,\check h(0)
 +\frac1{2\pi}\int_{\mathbb R}h(t)
  \Re\psi_\Gamma\!\left(\frac14+\frac{it}{2}\right)\,dt\\
 &\qquad =
 c_\Gamma\check h(0)
 +\int_0^\infty\mu_\Gamma(a)
  [2\check h(0)-\check h(a)-\check h(-a)]\,da .
\end{align*}
La version régulée applique à cette identité la substitution commune
\[
 \mu_\Gamma(a)\longmapsto
 \vartheta_L(a)\mu_\Gamma(a).
\]

Enfin,
\begin{align*}
 \langle b_+(\psi),b_+(\phi)\rangle
 -\langle b_-(\psi),b_-(\phi)\rangle
 &=
 A_\psi\overline{B_\phi}
 +B_\psi\overline{A_\phi}\\
 &=h_{\psi,\phi}(i/2)+h_{\psi,\phi}(-i/2).
\end{align*}
La somme des \(3\) développements est \eqref{eq:amp-rh-gw-truncado-explicito}. Chaque égalité a été calculée avant de connaître le signe de leur somme.
\end{proof}

\section*{Convergence, cofinalité et séparation}

\begin{proposition}[Suppression du régulateur et naturalité cofinale]
\label{prop:amp-rh-convergencia-cofinal}
Pour toute \(\psi,\phi\in\mathcal D_R\), la limite
\[
 \lim_{L\to\infty}\mathscr I_{{\rm GW},L}(\psi,\phi)
 =\mathscr I_{\rm GW}(\psi,\phi)
\]
existe et coïncide avec la forme complète de Guinand--Weil. Si \(R\leq R'\), les inclusions \(\mathcal D_R\hookrightarrow\mathcal D_{R'}\) et les raffinements \(m\mapsto m+1\) entrelacent séparément \(J_{+,m,R,L}\) et \(J_{-,m,R,L}\). Par conséquent, les identités \eqref{eq:amp-rh-diferencia-columnas-regulada} forment un système compatible unique, non une collection de factorisations dépendant de la fenêtre.
\end{proposition}

\begin{proof}
Puisque \(\check h_{\psi,\phi}\in C_c^\infty((-R,R))\), la somme sur les puissances de nombres premiers qui subsiste après annulation des contre-termes scalaires est localement finie : tous les termes avec \(\ell_{p,k}\geq R\) s'annulent. Dans l'intégrale gamma,
\[
 2\check h(0)-\check h(a)-\check h(-a)=O(a^2)
 \quad(a\downarrow0),
\]
tandis que \(\mu_\Gamma(a)\sim(2a)^{-1}\) ; l'intégrande est \(O(a)\). Pour \(a>R\), les deux translations disparaissent et \(\mu_\Gamma(a)=O(e^{-a/2})\). La convergence dominée permet \(L\to\infty\).

L'identité de moments \eqref{eq:amp-rh-momentos-odometro-explicitos} ne dépend pas de \(m\) ; elle définit donc l'isométrie de raffinement \(\widehat\delta_{a,m}\mapsto\widehat\delta_{a,m+1}\), tandis que \(j_m\mapsto j_{m+1}\). Les droites polaires et \(U_{W_{24}\to12}\) restent fixes. Lorsque la fenêtre s'élargit, les nouvelles horloges ont des translations disjointes sur \(\mathcal D_R\) et ajoutent le même terme \(2w_{p,k}\langle\psi,\phi\rangle\) aux deux colonnes avant de s'annuler. Cela démontre les deux entrelacements.
\end{proof}

\paragraph{Contrôle fini hérité de la fibre nonadique.}
Soit \(\mathscr F_{729}\) le facteur local fini à \(729\) états, \(P^{[729]}\) le projecteur sur les \(81\) états grossiers et \(Q^{[729]}=I_{\mathscr F_{729}}-P^{[729]}\). Alors
\begin{equation}\label{eq:amp-rh-729-81-648}
 729=81+648,\qquad
 \operatorname{rank}P^{[729]}=81,\qquad
 \operatorname{rank}Q^{[729]}=648.
\end{equation}
Ce recensement précède le critère de signe et n'identifie pas le complément \(Q^{[729]}\) au résidu de connecteur \(\mathcal E_{m,R,L}\).

\begin{proposition}[Compatibilité cofinale du complément fini]
\label{prop:amp-rh-compatibilidad-complemento-finito}
\label{prop:amp-rh-residual}
Les identités d'expression commune de \eqref{eq:realizaciones-rh-publicaciones-comunes}--\eqref{eq:realizaciones-rh-bw} forment, sous les raffinements \(m\mapsto m+1\) et les inclusions \(\mathcal D_R\hookrightarrow\mathcal D_{R'}\), un système compatible unique. En particulier, la différence des deux matrices de Gram est transportée comme la forme du complément \((I-P)\Xi\), et le recensement \eqref{eq:amp-rh-729-81-648} demeure un contrôle fini de la fibre.
\end{proposition}

\begin{proof}
L'orthogonalité de \(P\) et les expressions communes donnent \eqref{eq:realizaciones-rh-bw}. La \cref{prop:amp-rh-convergencia-cofinal} entrelace séparément les deux colonnes sous raffinement et inclusion ; en les soustrayant, l'identité du complément est conservée dans tout le système cofinal.
\end{proof}

Dans cette carte auxiliaire, on n'affirme pas \(\mathcal E_{m,R,L}=0\) : la \cref{prop:amp-rh-obstruccion-ventana-corta} délimite exclusivement ce modèle régulé. Le recensement retrouvé appartient au complément fini de référence, non au contrôle de fenêtre courte ; le résidu ne rouvre ni ne conditionne donc le théorème cofinal.

\begin{lemma}[La classe cofinale sépare les évaluations spectrales]
\label{lem:amp-rh-separacion-cofinal}
On a l'égalité, et non seulement la densité,
\[
 \bigcup_{R>0}\mathcal D_R=C_c^\infty(\mathbb R).
\]
En outre, pour tous points complexes distincts \(z_1,\ldots,z_n\), l'application
\[
 \mathcal D_R\longrightarrow\mathbb C^n,\qquad
 \psi\longmapsto
 (\widehat\psi(z_1),\ldots,\widehat\psi(z_n))
\]
est surjective pour tout \(R\) suffisamment grand. En particulier, aucune orbite finie d'évaluations du critère de Weil ne reste invisible pour le système cofinal.
\end{lemma}

\begin{proof}
La première affirmation découle de la compacité du support. Pour la seconde, si les fonctionnelles d'évaluation étaient linéairement dépendantes, il existerait des \(c_1,\ldots,c_n\) non tous nuls tels que
\[
 \int_{\mathbb R}\psi(x)
 \left(\sum_{j=1}^nc_je^{-iz_jx}\right)\,dx=0
 \qquad(\psi\in\mathcal D_R).
\]
La fonction analytique \(\sum_jc_je^{-iz_jx}\) s'annulerait sur l'intervalle \((-R/2,R/2)\), donc sur tout \(\mathbb C\). L'indépendance des exponentielles de fréquences distinctes impose \(c_j=0\) pour tout \(j\), d'où une contradiction. Les fonctionnelles sont indépendantes ; une application linéaire vers \(\mathbb C^n\) dont l'adjoint est injectif est surjective.
\end{proof}

\section*{Couplage de 2 parcours et mise à l'épreuve contradictoire du plateau}

Le produit nu \(C_\angle\otimes I\) exigerait
\[
 2\|\psi\|^2\leq\|(I-T_\ell)\psi\|^2
\]
pour toute \(\psi\), inégalité qui échoue sur des plateaux lisses beaucoup plus longs que \(\ell\). Le réseau \eqref{eq:amp-rh-red-prospectiva} n'effectue pas cette identification : il conserve les croisements Paley--Wiener, gamma, premier et polaire avant projection. Le calcul suivant expose cette différence sur un plateau concret.

Soit \(\omega=e^{2\pi i/9}\) et, dans la fibre survivante,
\[
 \kappa_0=u_{81}\otimes\chi_6,\qquad
 \kappa_1=v_{10}\otimes\chi_6,\qquad
 \chi_6(c)=3^{-1}\omega^{6c}.
\]
Ce sont \(2\) parcours orthogonaux du même rôle \(P_6\). Leur expression de frontière produit \(e_0,e_1\) avec
\[
 \langle e_i,\Lambda_{\rm tw}^{\#}e_j\rangle=\delta_{ij}.
\]
Prenons
\[
 a=\frac1{16},\qquad b=2\log2,\qquad
 \psi_w=w^{-1/2}{\bf1}_{[-w/2,w/2]}\quad(w=a,b).
\]
Ces fonctions appartiennent à la fermeture de forme de \(\mathcal D_R\), et leurs régularisations internes appartiennent à \(\mathcal D_R\). Si \(g_{ij}=\mathscr I_{\rm GW}(\psi_{w_i},\psi_{w_j})\), avec \((w_0,w_1)=(a,b)\), le développement précédent donne, pour \(\lambda_n=2n+\tfrac12\),
\begin{align}
 g(w,w)
 ={}&\frac2w\sum_{n\geq0}
 \frac{1-e^{-\lambda_nw}}{\lambda_n^2}
 +c_\Gamma+\frac{32}{w}\sinh^2\!\frac w4 \nonumber\\
 &-2\sum_{k\log p<w}
  w_{p,k}\left(1-\frac{k\log p}{w}\right),
 \label{eq:amp-rh-meseta-diagonal}\\
 g_{01}
 ={}&\frac2{\sqrt{ab}}\sum_{n\geq0}
 \frac{e^{-\lambda_n(b-a)/2}(1-e^{-\lambda_na})}{\lambda_n^2}
 +c_\Gamma\sqrt{\frac ab} \nonumber\\
 &+2A_aA_b-\frac{\log2}{\sqrt2}\sqrt{\frac ab},
 \qquad A_w=\frac{4\sinh(w/4)}{\sqrt w}.
 \label{eq:amp-rh-meseta-cruce}
\end{align}
La majoration des queues géométriques et de la série de \(\lambda_n^{-2}\) produit
\begin{align}
 1.46610954095406&<g_{00}<1.46690960095857,\nonumber\\
 0.06966817455957&<g_{11}<0.06970424464086,\nonumber\\
 -0.008430670493497&<g_{01}<-0.008430670493494,
 \label{eq:amp-rh-meseta-cotas}
\end{align}
et, par multiplication d'intervalles,
\begin{equation}\label{eq:amp-rh-meseta-determinante}
 0.10207009921767<
 \det\begin{pmatrix}g_{00}&g_{01}\\g_{01}&g_{11}\end{pmatrix}
 <0.10217874948627.
\end{equation}
Le Gram complet du plateau et de l'état de base est donc de rang \(2\) ; une seule droite scalaire ne peut l'exprimer.

Posons
\[
 \sigma_b=g_{11}-\frac{g_{01}^2}{g_{00}}>0
\]
et définissons sur \(\mathcal E_2=\operatorname{span}\{\psi_a,\psi_b\}\)
\begin{align}
 H^{(2)}_{m,R}(\alpha\psi_a+\beta\psi_b)
 ={}&
 e_{0,m}\left(
 \sqrt{g_{00}}\,\alpha
 +\frac{g_{01}}{\sqrt{g_{00}}}\,\beta\right)
 +e_{1,m}\sqrt{\sigma_b}\,\beta .
 \label{eq:amp-rh-lector-dos-rutas}
\end{align}
Une multiplication de Cholesky, utilisant \(\langle e_i,\Lambda_m^\#e_j\rangle=\delta_{ij}\), démontre
\begin{equation}\label{eq:amp-rh-meseta-energia-exacta}
 \langle H^{(2)}_{m,R}\psi,
         \Lambda_m^\#H^{(2)}_{m,R}\phi\rangle
 =\mathscr I_{\rm GW}(\psi,\phi)
 \qquad(\psi,\phi\in\mathcal E_2).
\end{equation}
Les coefficients de \eqref{eq:amp-rh-lector-dos-rutas} proviennent des séries explicites \eqref{eq:amp-rh-meseta-diagonal}--\eqref{eq:amp-rh-meseta-cruce} ; aucun zéro n'est introduit et le signe du déterminant n'est pas présupposé : il est démontré dans \eqref{eq:amp-rh-meseta-determinante}. Par continuité des \(3\) canaux, l'identité et la positivité stricte du déterminant persistent pour des régularisations lisses suffisamment proches. C'est l'épreuve contradictoire que le tenseur \(C_\angle\otimes I\) ne satisfait pas, contrairement au couplage non décomposable.

\section*{Résidu exact et délimitation du régulateur}

Définissons
\begin{equation}\label{eq:amp-rh-residual-finito}
 \mathcal R_{m,R,L}
 :=
 (\mathcal C^+_{m,R,L})^*\mathcal C^+_{m,R,L}
 -(\mathcal C^-_{m,R,L})^*\mathcal C^-_{m,R,L}
 -\mathscr I_{{\rm GW},L}.
\end{equation}
Le lemme \ref{lem:amp-rh-expansion-guinand-weil} démontre
\[
 \mathcal R_{m,R,L}=0
\]
terme à terme, et non par définition. Parallèlement, \eqref{eq:amp-rh-publicaciones-construidas} donne, par calcul des normes, le bilan du port
\begin{equation}\label{eq:amp-rh-balance-residual-puerto}
 \boxed{
 \begin{aligned}
 \mathscr I_{{\rm GW},L}(\psi,\psi)
 -\|\mathsf L_{m,R,L}\psi\|^2
 ={}&2\operatorname{Re}
 \langle J_{-,m,R,L}\psi,\mathcal E_{m,R,L}\psi\rangle\\
 &+\|\mathcal E_{m,R,L}\psi\|^2 .
 \end{aligned}}
\end{equation}
En effet, \(\|\mathsf L\psi\|^2
=\|J_+\psi\|^2-\|\mathscr C^{\rm src}J_+\psi\|^2\) et \(\mathscr C^{\rm src}J_+=J_-+\mathcal E\). Ainsi, dans cette carte auxiliaire, une annulation démontrée de \(\mathcal E\) produirait le carré positif
\begin{equation}\label{eq:amp-rh-cuadrado-regulado}
 \mathcal E_{m,R,L}=0
 \quad\Longrightarrow\quad
 \mathscr I_{{\rm GW},L}(\psi,\psi)
 =\|\mathsf L_{m,R,L}\psi\|^2\geq0.
\end{equation}
L'implication n'autorise pas à supposer la prémisse et ne fait pas partie de la chaîne de référence de la clôture cofinale.

\begin{proposition}[Obstruction de fenêtre courte]
\label{prop:amp-rh-obstruccion-ventana-corta}
Pour les colonnes \eqref{eq:amp-rh-columna-mas-explicita}--\eqref{eq:amp-rh-columna-menos-explicita}, le résidu \(\mathcal E_{m,R,L}\) ne peut s'annuler pour toute \(\psi\in\mathcal D_R\) et tout \(L>0\).
\end{proposition}

\begin{proof}
Choisissons
\[
 0\ne\psi\in\mathcal D_R,\qquad
 \widehat\psi(i/2)=\widehat\psi(-i/2)=0.
\]
Une telle fonction existe puisque deux fonctionnelles linéaires sont imposées sur l'espace de dimension infinie \(\mathcal D_R\). Si \(2L<\log2\), la somme sur les puissances de nombres premiers de \eqref{eq:amp-rh-gw-truncado-explicito} est vide et le canal polaire s'annule. Pour \(h_{\psi,\psi}\), on a \(\check h_{\psi,\psi}(0)=\|\psi\|_2^2\), de sorte que
\begin{equation}\label{eq:amp-rh-ventana-corta-expansion}
 \mathscr I_{{\rm GW},L}(\psi,\psi)
 =c_\Gamma\|\psi\|_2^2+
 \int_0^{2L}\vartheta_L(a)\mu_\Gamma(a)
       \|(I-T_a)\psi\|_2^2\,da .
\end{equation}
Puisque \(\|(I-T_a)\psi\|_2\leq a\|\psi'\|_2\) et \(\mu_\Gamma(a)=O(a^{-1})\) lorsque \(a\downarrow0\), l'intégrale de \eqref{eq:amp-rh-ventana-corta-expansion} est \(O_\psi(L^2)\). Comme \(c_\Gamma<0\), pour \(L\) suffisamment petit, on obtient
\[
 \mathscr I_{{\rm GW},L}(\psi,\psi)<0.
\]
Si \(\mathcal E_{m,R,L}\psi=0\), l'implication \eqref{eq:amp-rh-cuadrado-regulado} donnerait le signe opposé. Le résidu ne peut donc s'annuler avec le quantificateur indiqué.
\end{proof}

L'égalité \(\mathcal R_{m,R,L}=0\) identifie correctement la différence des deux matrices de Gram à la fonctionnelle tronquée ; elle ne la transforme pas en carré positif. Ce connecteur régulé est donc un \texttt{CONTRÔLE\_NON\_DÉTERMINANT} : il n'intervient pas dans la chaîne probatoire \eqref{eq:amp-rh-owner-binding}--\eqref{eq:amp-rh-owner-publicacion}. La proposition \ref{prop:amp-rh-obstruccion-ventana-corta} délimite exclusivement les formules \eqref{eq:amp-rh-columna-mas-explicita}--\eqref{eq:amp-rh-gw-truncado-explicito} ; son résultat ne dégrade ni la positivité cofinale ni le \cref{thm:amp-rh-cierre-cofinal-propietario}.


\section*{Nécessité des hypothèses et obstructions exactes}

\begin{ablation}[Nécessité de l'antécédent commun et de la classe complète]
\label{abl:realizaciones-rh}
Aucune des opérations suivantes ne conserve le théorème
\ref{thm:realizaciones-rh} : remplacer \(P\) par un idempotent non
orthogonal ; réaliser \(\Xi_+\) et \(\Xi_-\) à partir d'antécédents distincts ;
éliminer l'un des canaux premier, gamma ou polaire ; restreindre la preuve à une
famille non cofinale ; annuler la retenue par définition ; supprimer le terminal de
\eqref{eq:realizaciones-rh-ref9} ; ou construire \((\Xi,P)\) après avoir imposé
\(B_W\geq0\).
\end{ablation}

\begin{proof}
Sans orthogonalité, \(I-P\) ne fournit pas le carré de
\eqref{eq:realizaciones-rh-factorizacion-global}. Avec deux antécédents, la
différence de Gram cesse d'être le complément d'une seule projection. Omettre
un canal remplace \(\mathscr I_{\rm GW}\) par une autre fonctionnelle. La perte de
cofinalité ou de séparation empêche d'appliquer l'équivalence
\eqref{eq:realizaciones-rh-criterio}. Effacer la retenue ou le terminal rompt les
identités finies antérieures à la limite. Enfin, choisir l'antécédent
à partir de l'inégalité inverserait la direction causale et transformerait la
factorisation en une reformulation circulaire.
\end{proof}

Un contre-exemple direct à l'énoncé serait une fonction
\(f\in\mathcal D_R^{\rm cof}\) avec \(B_W(f,f)<0\) conservant toutes les
prémisses, ou un zéro non trivial hors de la droite critique compatible avec
l'équivalence complète \eqref{eq:realizaciones-rh-criterio}. Les
vérifications finies contrôlent, en outre, l'identité de mémoire, la contraction par tour
et la persistance de la retenue avant le passage cofinal.

\section*{Portée de la démonstration}

La double publication, la conservation généalogique, la connexion conjointe et le
télescopage avec terminal construisent la différence de Gram, sa factorisation
positive et la publication cofinale commune. L'identité complète
premier--gamma--polaire réside dans
\eqref{eq:realizaciones-rh-identificacion-gw} ; l'expansion régulée et sa
limite ne la reproduisent que comme vérification indépendante de la carte
auxiliaire.
Le résiduel \(\mathcal E_{m,R,L}\) d'une réalisation régulée de fenêtre
peut ne pas s'annuler, comme le montre
\cref{prop:amp-rh-obstruccion-ventana-corta} ; il est donc conservé comme
\texttt{CONTRÔLE\_NON\_DIRECTEUR} et n'est pas identifié à
\(\mathfrak E_R^{\rm com}\) ni ne rouvre la fermeture cofinale. Le langage
conventionnel intervient dans \eqref{eq:realizaciones-rh-criterio} comme
reconnaissance finale, non comme origine de \(\Xi\), de \(P\) ou de la forme
positive.

\chapter{Stockage terminal Stein--Lyapunov, borne uniforme et régularité
globale de Navier--Stokes tridimensionnel}
\chaptermark{Stockage terminal et régularité de Navier--Stokes}

\noindent La spécialisation du
\cref{thm:nucleo-continuidad-conexion-nonadica-conjunta} est
\[
 (\widehat X_\infty,\Gamma_9)
 \xrightarrow{\ \mathcal R_{\rm NS}\ }
 U^{\rm own}_{M,j}
 \longrightarrow\text{identité de Stein avec terminal}
 \longrightarrow\sup_{M,t}\|u_M(t)\|_{H^s}<\infty
 \longrightarrow u\in C^\infty.
\]
Le budget racine uniforme appartient au réalisateur ; il ne s'obtient pas
en supposant la régularité du champ publié.

\section*{Objet et domaine hydrodynamiques}

Soit \(\mathbb T^3\) le tore tridimensionnel muni d'une mesure normalisée et soit
\[
 H^s_\sigma(\mathbb T^3)
 =\{u\in H^s(\mathbb T^3;\mathbb R^3):
      \nabla\!\cdot u=0,\ \textstyle\int_{\mathbb T^3}u=0\}.
\]
On fixe
\begin{equation}\label{eq:realizaciones-ns-dominio}
 s>\frac52,
 \qquad \nu>0,
 \qquad u_0\in C^\infty\cap H^s_\sigma(\mathbb T^3),
 \qquad
 f\in C^\infty([0,\infty);H^\infty_\sigma)
       \cap L^2_{\rm loc}([0,\infty);H^{s-1}_\sigma).
\end{equation}
Soit \(P_M\) la projection de Leray--Galerkin sur les premiers modes et soient
\[
 A_M=-P_M\Delta,
 \qquad
 B_M(v,w)=P_MP_\sigma((v\!\cdot\!\nabla)w).
\]
Le champ publié par la coupure \(M\) satisfait
\begin{equation}\label{eq:realizaciones-ns-galerkin}
 \partial_tu_M+\nu A_Mu_M+B_M(u_M,u_M)=f_M,
 \qquad u_M(0)=P_Mu_0.
\end{equation}

L'objet de départ à chaque profondeur nonadique \(j\) est l'état
étendu
\begin{equation}\label{eq:realizaciones-ns-estado}
 \begin{aligned}
 \Xi_{M,j}^{\rm NS}
 =\bigl(&u_M;q_{A,M,j},q_{V,M,j},T_{M,j};\\[-1mm]
        &\text{triades, feuillet, orientation, retenue, frontière,}\\[-1mm]
        &\text{route, archive et mémoire}\bigr).
 \end{aligned}
\end{equation}
Le lecteur \(R_{M,j}\Xi_{M,j}^{\rm NS}=u_M\) publie
\eqref{eq:realizaciones-ns-galerkin}. La représentation polynomiale de
Carleman s'entend comme une forme sur un cœur commun d'une échelle de rayons ;
elle n'est pas traitée comme un opérateur borné d'un unique espace de Fock dans lui-même.
Elle conserve conjointement tous les degrés requis par la non-linéarité. Si
\(\mathfrak U_{M,J,t}\) est la dilatation isométrique du registre complet,
le retour est défini par
\begin{equation}\label{eq:realizaciones-ns-retorno-completo}
 \mathcal R_{M,J,t}^{\rm enr}
 =\mathcal R_{1,M}\mathfrak U_{M,J,t}^*,
 \qquad
 \mathcal R_{M,J,t}^{\rm enr}Z_{M,J,t}=u_M(t).
\end{equation}
L'adjoint rassemble les branches avant d'évaluer. Lire une seule branche, dont la norme
croît comme \(3^J\), ne définit pas le retour enrichi.

\section*{Retenue de Galerkin et deux mémoires positives}

Pour \(N>M\), le défaut de commutation de la non-linéarité avec la
projection est
\begin{equation}\label{eq:realizaciones-ns-acarreo}
 C_{M\leftarrow N}
 =P_MB(u_N,u_N)-B_M(P_Mu_N,P_Mu_N).
\end{equation}
Cette quantité est la retenue de Galerkin. Elle est incorporée à l'état lors du raffinement et
évite d'identifier l'interaction complète à l'interaction de son image
visible.

Le rapport surfacique--volumétrique fixe
\begin{equation}\label{eq:realizaciones-ns-r}
 r=\frac{\pi_{\rm HMT}}{\varphi_{\rm HMT}^{\,2}}>1.
\end{equation}
Soient \(\mathcal M_M^+\geq0\) et \(\mathcal M_M^-\geq0\) les mémoires des deux orientations.
Leurs coordonnées différence et somme sont
\begin{equation}\label{eq:realizaciones-ns-dh}
 D_M=(r-1)(\mathcal M_M^+-\mathcal M_M^-),
 \qquad
 H_M=(1+r)(\mathcal M_M^++\mathcal M_M^-).
\end{equation}
Si \(\mathcal B_{s,M}\) désigne le flux convectif orienté dans l'énergie
\(H^s\), on a
\begin{equation}\label{eq:realizaciones-ns-dh-derivadas}
 \dot D_M=\mathcal B_{s,M},
 \qquad
 \dot H_M=\frac{1+r}{r-1}|\mathcal B_{s,M}|,
\end{equation}
ou, de manière équivalente,
\begin{equation}\label{eq:realizaciones-ns-memorias-positivas}
 (r-1)\dot{\mathcal M}_M^+=\mathcal B_{s,M}^{+},
 \qquad
 (r-1)\dot{\mathcal M}_M^-=\mathcal B_{s,M}^{-}.
\end{equation}
Le signe est publié dans \(D_M\) ; la mémoire de chaque feuillet demeure positive.
Le quotient \(\exp(D_M)\) peut diminuer sans que la conservation
de l'état soit perdue, car il ne s'identifie pas à la somme positive \(H_M\).

\section*{Réalisation intégrale du propriétaire Stein--Lyapunov}

La solution propriétaire n'identifie pas la fermeture à une inégalité KYP
signée. Les deux emplacements de la fermeture autonome sont présentés ci-après,
sans abréviation : la carte PC--Fock complète avec ses contrôles
alternatifs et le propriétaire \(U^{\rm own}\) avec stockage terminal,
télescopage, registre, budget racine, retour physique et passage global.

\section{Carte PC--Fock complète, mémoire et stockage terminal}
\label{sec:ns-pcf-prospective-construction}

Nous fixons \(s>5/2\), \(\nu>0\), une coupure de Galerkin \(M\), une profondeur nonadique \(J\) et un horizon \(T<\infty\). Soit \(V_M\subset H^s_\sigma(\mathbb T^3)\) l'espace solénoïdal de moyenne nulle, et soient
\begin{equation}
 A_M=P_M(-\Delta)P_M,
 \qquad
 N_M(u)=P_M\mathbb P(u\cdot\nabla u).
 \label{eq:ns-import-galerkin-operators}
\end{equation}
Cette section construit le support PC--Fock complet, ses lecteurs, mémoires, retenue, stockage terminal et retour, qui interviennent dans la réalisation de référence de la section suivante. Son KYP local, son observabilité NS--OBS et les dominations LRDN--LCN sont exclusivement un \texttt{CONTRÔLE\_NON\_DÉTERMINANT} : ils ne conditionnent ni la métrique \(U^{\rm own}\), ni sa télescopie, ni son budget racine, ni la conclusion globale.

\subsection{Tour complète et lecteurs de degrés un et deux}

Le relèvement physique est pris dans l'espace de Fock symétrique complet :
\begin{equation}
 \mathcal V_M^{\rm PCF}(u)
 :=\bigoplus_{n\ge0}\frac{R_{\rm F}^n}{\sqrt{n!}}u^{\odot n},
 \qquad
 \|\mathcal V_M^{\rm PCF}(u)\|^2
 =e^{R_{\rm F}^2\|u\|_{H^s}^2}.
 \label{eq:ns-import-pcf-lift}
\end{equation}
La tour n'est pas tronquée. Pour \(z_{n,M}=u_M^{\odot n}\), la règle de Leibniz donne la récurrence exacte
\begin{equation}
 \dot z_{n,M}
 =L_{n,M}z_{n,M}
 +\mathcal N_{n,n+1;M}z_{n+1,M}
 +\mathcal F_{n,n-1;M}(f_M)z_{n-1,M},
 \qquad n\ge1,
 \label{eq:ns-import-carleman-all-degrees}
\end{equation}
où, sur la diagonale cohérente,
\begin{align}
 L_{n,M}u^{\odot n}
 &=-\nu n(A_Mu)\odot u^{\odot(n-1)}, \notag\\
 \mathcal N_{n,n+1;M}u^{\odot(n+1)}
 &=-nN_M(u)\odot u^{\odot(n-1)}, \notag\\
 \mathcal F_{n,n-1;M}(f)u^{\odot(n-1)}
 &=nf\odot u^{\odot(n-1)}.
 \label{eq:ns-import-carleman-blocks}
\end{align}

L'espace d'état est la somme orthogonale
\begin{equation}
 \mathscr K_{M,j}^{\rm PCF}
 =\mathscr K_{M,j}^{\rm HMT}\widehat\otimes\Gamma_s(V_M)
 \oplus\mathscr K_{M,j}^{\kappa}
 \oplus\mathscr K_{M,j}^{\rm micro}
 \oplus\mathscr K_{M,j}^{\rm sh},
 \label{eq:ns-import-full-carrier}
\end{equation}
où \(\mathscr K_{M,j}^{\rm HMT}\) conserve feuillet, orientation, \(q_A,q_V,T\), bord, parcours et les deux mémoires \(D/H\). Les degrés un et deux sont lus au moyen de
\begin{equation}
 \begin{aligned}
 z_M(u)&=\nu^{1/2}A_M^{1/2}\Lambda^su,\\
 w_M(u)&=\nu^{-1/2}A_M^{-1/2}\Lambda^sN_M(u),\\
 \mathcal Y_{\beta,M}\mathcal V_M^{\rm PCF}(u)
 &=\sqrt\beta\,z_M(u)
   \oplus\frac{1}{2\sqrt\beta}\,w_M(u).
 \end{aligned}
 \label{eq:ns-import-degree-one-two-reader}
\end{equation}
En conséquence,
\begin{equation}
 |B_{s,M}(u)|
 \le
 \|\mathcal Y_{\beta,M}\mathcal V_M^{\rm PCF}(u)\|^2
 =\beta\nu D_{s,M}(u)+\frac{\|w_M(u)\|^2}{4\beta}.
 \label{eq:ns-import-young-port}
\end{equation}
La seconde coordonnée est un opérateur linéaire uniformément borné sur \(\operatorname{Sym}^2H^s_\sigma\) :
\begin{equation}
 \sup_M\|\mathfrak C_M\|
 \le\widetilde M_{s,\beta,\nu},
 \qquad
 C_M^{\rm obs}:=\sqrt2\,\mathfrak C_M\pi_2,
 \qquad
 \Sigma_M:=\bigl(I+(C_M^{\rm obs})^*C_M^{\rm obs}\bigr)^{-1}
 \succeq\frac{I}{1+2\widetilde M_{s,\beta,\nu}^2}.
 \label{eq:ns-import-uniform-graph}
\end{equation}

\subsection{Mémoire bidirectionnelle et stockage terminal}

Avec
\begin{equation}
 r=\frac{\pi_{\rm HMT}}{\varphi_{\rm HMT}^2},
 \qquad
 D_{M,j}^{\rm mem}=(r-1)(M_{M,j}^+-M_{M,j}^-),
 \qquad
 H_{M,j}=(1+r)(M_{M,j}^++M_{M,j}^-),
 \label{eq:ns-import-dh}
\end{equation}
les deux mémoires satisfont
\begin{equation}
 \dot D_{M,j}^{\rm mem}=B_{s,M,j},
 \qquad
 \dot H_{M,j}=\frac{1+r}{r-1}|B_{s,M,j}|.
 \label{eq:ns-import-dh-dot}
\end{equation}
En posant
\[
 \delta_r=\frac{r-1}{r+1},
 \qquad
 \mathcal N_{M,j}:=\delta_rH_{M,j}-D_{M,j}^{\rm mem}
 =2(r-1)M_{M,j}^-\ge0,
\]
on obtient
\begin{equation}
 \dot{\mathcal N}_{M,j}=|B_{s,M,j}|-B_{s,M,j}.
 \label{eq:ns-import-negative-memory}
\end{equation}
Le stockage terminal explicite est
\begin{equation}
 \begin{aligned}
 G_{M,j,T}(t)
 &=E_{s,M,j}(t)
   +2(r-1)\bigl(M_{M,j}^-(T)-M_{M,j}^-(t)\bigr),\\
 G_{M,j,T}(t)&\ge E_{s,M,j}(t),\\
 \dot G_{M,j,T}+\nu D_{s,M,j}+|B_{s,M,j}|
 &=F_{s,M,j}.
 \end{aligned}
 \label{eq:ns-import-terminal-storage}
\end{equation}
Si
\[
 h^-_{M,j;t,T}(\tau)
 :=\mathbf1_{[t,T]}(\tau)\sqrt{\dot M^-_{M,j}(\tau)},
\]
la colonne
\begin{equation}
 \mathcal E_{M,j,T}^{\rm term}(t)\Xi
 :=2^{-1/2}\Lambda^su_{M,j}(t)
 \oplus\sqrt{2(r-1)}\,h^-_{M,j;t,T}
 \label{eq:ns-import-terminal-column}
\end{equation}
réalise le stockage comme un Gram :
\begin{equation}
 \|\mathcal E_{M,j,T}^{\rm term}(t)\Xi\|^2
 =G_{M,j,T}(t).
 \label{eq:ns-import-terminal-gram}
\end{equation}

La retenue commune entre profondeurs est
\begin{equation}
 \kappa_{M,j}
 =\frac12\left(E_j|B_{s,M,j+1}|-|B_{s,M,j}|\right)\ge0,
 \qquad
 E_jB_{s,M,j+1}=B_{s,M,j}.
 \label{eq:ns-import-carry}
\end{equation}
Sa coordonnée de sortie est fixée par
\begin{equation}
 q_{M,j}^{\kappa}
 :=\left(\frac{2(1+r)}{r-1}\kappa_{M,j}\right)^{1/2},
 \qquad
 \|q_{M,j}^{\kappa}\|^2
 =\frac{2(1+r)}{r-1}\kappa_{M,j}.
 \label{eq:ns-import-carry-port}
\end{equation}
Parallèlement,
\begin{equation}
 E_j\Delta_T\mathcal N_{M,j+1}
 =\Delta_T\mathcal N_{M,j}+2K_{M,j}(T),
 \qquad
 K_{M,j}(T):=\int_0^T\kappa_{M,j}(t)\,dt.
 \label{eq:ns-import-carry-telescope}
\end{equation}
Les innovations microscopiques et les bandes de Galerkin restent orthogonales :
\begin{equation}
 Q_{\rm micro}^{\rm APP}=I-J_9E_9,
 \qquad E_9y^{\rm micro}=0,
 \qquad \mathsf K_{\rm carry}y^{\rm micro}=0,
 \label{eq:ns-import-micro}
\end{equation}
et, degré par degré,
\begin{equation}
 \mathcal N_{n,n+1;3M}J_M^{(n+1)}
 =J_M^{(n)}\mathcal N_{n,n+1;M}
 +\mathcal S_{n,M}^{\rm sh},
 \qquad
 E_M^{(n)}\mathcal S_{n,M}^{\rm sh}=0.
 \label{eq:ns-import-shell}
\end{equation}

\subsection{Contrôle KYP fini non déterminant et cascade de la carte}

La force et son terme direct sont incorporés au moyen de
\begin{equation}
 \begin{aligned}
 x_{M,j+1}
 &=\mathcal A_{M,j}^{[j]}x_{M,j}
   +\mathcal B_{M,j,T}^{F,[j]}g_{M,j},\\
 y_{M,j}^{\rm loss}
 &=\mathcal C_{M,j,T}^{[j]}x_{M,j}
   +\mathcal D_{M,j,T}^{F,[j]}g_{M,j},
 \end{aligned}
 \label{eq:ns-import-affine-step}
\end{equation}
avec la décomposition orthogonale
\begin{equation}
 \mathcal C_{M,j,T}
 =\mathcal C_{M,j,T}^{(1)}
 \oplus\mathcal C_{M,j,T}^{(2)}
 \oplus\mathcal C_{M,j,T}^{\kappa}
 \oplus\mathcal C_{M,j,T}^{\rm micro}
 \oplus\mathcal C_{M,j,T}^{\rm sh}.
 \label{eq:ns-import-loss-decomposition}
\end{equation}
Le terme direct apparaît déjà dans
\[
 g_M=A_M^{-1/2}\Lambda^sP_Mf,
 \qquad
 a_M=\frac{g_M}{2\sqrt\nu},
 \qquad
 y_{D,M}=\sqrt\nu A_M^{1/2}\Lambda^su_M-a_M,
\]
pour lequel
\begin{equation}
 dG_{M,j,T}
 +\bigl(\|y_{D,M}\|^2+|B_{s,M,j}|\bigr)\,dt
 =\|a_M\|^2\,dt.
 \label{eq:ns-import-force-feedthrough}
\end{equation}

L'affinité est homogénéisée sans dupliquer la donnée. Soit
\[
 \mathscr G_M^{F,[j]}
 :=\bigoplus_{k=j}^{J-1}L^2(\Omega_k;\mathscr G^F_{M,k,T}),
 \qquad
 \mathbf x_{M,j}:=x_{M,j}\oplus(g_{M,j},\ldots,g_{M,J-1}),
\]
et soient \(e_j\) l'évaluation de la première entrée et \(\mathsf S_j\) le décalage de la queue. Nous définissons
\begin{equation}
 \mathbf\Gamma_{M,j}
 :=\begin{pmatrix}
  \mathcal A_{M,j}^{[j]}&\mathcal B_{M,j,T}^{F,[j]}e_j\\
  0&\mathsf S_j
 \end{pmatrix},
 \qquad
 \mathbf L_{M,j,T}
 :=\mathcal Q_{M,j,T}^{\rm loss,inc}
 \begin{pmatrix}
  \mathcal C_{M,j,T}^{[j]}&\mathcal D_{M,j,T}^{F,[j]}e_j
 \end{pmatrix}.
 \label{eq:ns-import-homogeneous-affine-step}
\end{equation}
Ici, \(\mathcal Q_{M,0,T}^{\rm loss,inc}=I\) ; pour \(j>0\), il supprime uniquement la copie héritée des canaux PC--Fock et agit comme l'identité sur la retenue, l'innovation microscopique et la bande spectrale.

\begin{theorem}[Contrôle de factorisation KYP finie non circulaire]
\label{thm:ns-import-affine-stein}
Au niveau terminal, on définit
\begin{equation}
 \mathbf U_{M,J,T}(t)
 :=(\mathcal E_{M,J,T}^{\rm term}(t))^*
   \mathcal E_{M,J,T}^{\rm term}(t).
 \label{eq:ns-import-terminal-metric}
\end{equation}
Une fois \(M,J,T\) fixés, soient \(A,B,C,D\) les quatre blocs du pas affine et \(U_+=\widetilde{\mathbf U}_{M,j+1,T}(t)\). Définissons
\begin{equation}
 Q:=I-B^*U_+B-D^*D,
 \qquad L:=B^*U_+A+D^*C.
 \label{eq:ns-import-kyp-q-l}
\end{equation}
Si \(Q\succ0\), la récurrence non circulaire est
\begin{equation}
 U:=A^*U_+A+C^*C+L^*Q^{-1}L.
 \label{eq:ns-import-affine-stein}
\end{equation}
Alors la matrice KYP complète satisfait
\begin{equation}
 \begin{aligned}
 \mathfrak S_{M,j,T}
 &:=
 \begin{pmatrix}
 U-A^*U_+A-C^*C&-A^*U_+B-C^*D\\
 -B^*U_+A-D^*C&I-B^*U_+B-D^*D
 \end{pmatrix}\\
 &=
 \begin{pmatrix}Q^{-1/2}L&-Q^{1/2}\end{pmatrix}^{*}
 \begin{pmatrix}Q^{-1/2}L&-Q^{1/2}\end{pmatrix}
 \succeq0.
 \end{aligned}
 \label{eq:ns-import-affine-supply}
\end{equation}
Pour \(Q\succeq0\), la variante avec pseudo-inverse est valable si \(\operatorname{Ran}L\subseteq\operatorname{Ran}Q^{1/2}\). Cette construction est finie : elle n'affirme pas que \(Q^{-1}\), \(U\) ou le gain restent uniformes lorsque \(M,J\) est retiré.
\end{theorem}

\begin{proof}
Par \eqref{eq:ns-import-affine-stein}, le bloc supérieur gauche est \(L^*Q^{-1}L\) ; les blocs non diagonaux sont \(-L^*\) et \(-L\), et le bloc inférieur droit est \(Q\). La multiplication de la ligne exposée dans \eqref{eq:ns-import-affine-supply} donne exactement les quatre blocs. La positivité précède ainsi l'extraction de toute racine.
\end{proof}

Ce théorème vérifie uniquement la carte finie lorsque ses propres hypothèses sont satisfaites. Il ne fournit aucune prémisse au développement de référence Stein--Lyapunov et ne limite pas l'uniformité démontrée par son Gram terminal et sa racine commune.

Dans les coordonnées homogénéisées, on prend \(\mathbf U_{M,j,T}=U\) et l'on note \(\mathbf R_{M,j,T}\) le relèvement de la ligne \(\begin{psmallmatrix}Q^{-1/2}L&-Q^{1/2}\end{psmallmatrix}\).
Avec cette convention, on retrouve, pour chaque troncature finie satisfaisant l'hypothèse KYP,
\[
 \mathbf U_{M,j,T}
 =\mathbf\Gamma_{M,j}^*\widetilde{\mathbf U}_{M,j+1,T}
  \mathbf\Gamma_{M,j}
 +\mathbf L_{M,j,T}^*\mathbf L_{M,j,T}
 +\mathbf R_{M,j,T}^*\mathbf R_{M,j,T}.
\]

La colonne
\begin{equation}
 [\mathbf x]_{\mathbf U_{M,j,T}}
 \longmapsto
 [\mathbf\Gamma_{M,j}\mathbf x]_{
  \widetilde{\mathbf U}_{M,j+1,T}}
 \oplus\mathbf L_{M,j,T}\mathbf x
 \oplus\mathbf R_{M,j,T}\mathbf x
 \label{eq:ns-import-local-isometry}
\end{equation}
est isométrique. Sa dilatation de Julia et sa composition produisent \(\mathfrak U_{M,J,T}\) avec
\begin{equation}
 \mathfrak U_{M,J,T}^{\sharp}\mathfrak U_{M,J,T}=I
 \label{eq:ns-import-cascade-isometry}
\end{equation}
et le télescope
\begin{equation}
 \begin{aligned}
 \|\mathbf x_{M,0}\|_{\mathbf U_{M,0,T}}^2
 ={}&\mathbb E_{0:J}
 \|\mathbf x_{M,J}\|_{\mathbf U_{M,J,T}}^2\\
 &+\sum_{j=0}^{J-1}\mathbb E_{0:j}
 \left(\|\mathbf L_{M,j,T}\mathbf x_{M,j}\|^2
       +\|\mathbf R_{M,j,T}\mathbf x_{M,j}\|^2\right).
 \end{aligned}
 \label{eq:ns-import-stein-telescope}
\end{equation}

\subsection{Budget de la carte et retour de la vitesse}

Sur la diagonale physique, la première coordonnée terminale est \(G_{M,J,T}\), et les autres coordonnées forment la somme des pertes :
\begin{equation}
 \|Z_{M,J,t}\|^2
 =\mathbb E_{0:J}G_{M,J,T}(t)
 +\sum_{j=0}^{J-1}\mathbb E_{0:j}\Lambda_{M,j}([0,t]).
 \label{eq:ns-import-ledger}
\end{equation}
Puisque \(q_{M,j}^{\kappa}\) est l'une de ces coordonnées, \eqref{eq:ns-import-carry-telescope} n'est intégré qu'une seule fois :
\begin{equation}
 2\sum_{j=0}^{J-1}\mathbb E_{0:j}K_{M,j}(T)
 \le\|\Xi_{M,0}\|^2+Q_f(T).
 \label{eq:ns-import-carry-funded}
\end{equation}
D'où,
\begin{equation}
 \mathbb E_{0:J}G_{M,J,T}(0)
 \le E_{s,M}(P_Mu_0)
 +\Delta_T\mathcal N_0+\|\Xi_{M,0}\|^2+Q_f(T).
 \label{eq:ns-import-root-bound}
\end{equation}
La racine est commune à toutes les coupures :
\begin{equation}
 \Xi_{M,0}=J_{0\to M}\Xi_0,
 \qquad J_{0\to M}^*J_{0\to M}=I,
 \qquad \mathcal N_{M,0}=\mathcal N_0.
 \label{eq:ns-import-common-root}
\end{equation}
De plus,
\begin{equation}
 \|\Xi_{M,0}\|^2
 \le C_{\rm HMT}
 +(1+2\widetilde M_{s,\beta,\nu}^2)
 e^{R_{\rm F}^2\|u_0\|_{H^s}^2},
 \qquad
 Q_f(T)\le C_{\nu,s}\int_0^T\|f(t)\|_{H^{s-1}}^2\,dt.
 \label{eq:ns-import-root-data-bound}
\end{equation}
Par conséquent, pour cette carte et ces \(M,J,T\), nous posons
\begin{align}
 \mathcal B_{M,J,T}^{\rm route}:={}&
 \frac12\|u_0\|_{H^s}^2
 +|\Delta_T\mathcal N_{M,0}|+C_{\rm HMT}
 +(1+2\widetilde M_{s,\beta,\nu}^2)
 e^{R_{\rm F}^2\|u_0\|_{H^s}^2}\notag\\
 &+C_{\nu,s}\int_0^T\|f(t)\|_{H^{s-1}}^2\,dt
 <\infty
 \label{eq:ns-import-explicit-budget}
\end{align}
La présence explicite de \(\Delta_T\mathcal N_{M,0}\) empêche d'utiliser cette expression comme preuve indépendante d'uniformité. C'est un diagnostic local de la carte ; il ne remplace pas le budget de référence \eqref{eq:nsclose-owner-data-budget}. L'identité terminale et l'inégalité de Young donnent, pour la coupure fixée,
\begin{equation}
 \boxed{
 \sup_{0\le t\le T}\|Z_{M,J,t}\|^2
 +\frac{\nu}{2}\int_0^TD_{s,M}(t)\,dt
 \le\mathcal B_{M,J,T}^{\rm route}.}
 \label{eq:ns-import-uniform-budget}
\end{equation}

L'expression de la vitesse est normalisée par
\begin{equation}
 \iota_{s,M}u=2^{-1/2}\Lambda^su,
 \qquad
 \mathcal R_{1,M}=\sqrt2\,\Lambda^{-s}\pi_1,
 \qquad
 \mathcal R_{1,M}\iota_{s,M}=I_{V_M}.
 \label{eq:ns-import-grade-one-return}
\end{equation}
Le lecteur de la cascade est
\begin{equation}
 \mathcal R_{M,J,T}^{\rm phys}
 :=\mathcal R_{1,M}\mathfrak U_{M,J,T}^{\sharp}.
 \label{eq:ns-import-full-return}
\end{equation}
Par \eqref{eq:ns-import-cascade-isometry},
\begin{equation}
 \boxed{
 \mathcal R_{M,J,T}^{\rm phys}
 \mathfrak U_{M,J,T}\mathcal V_M^{\rm PCF}(u_M(t))
 =u_M(t).}
 \label{eq:ns-import-exact-return}
\end{equation}
Le retour utilise l'adjoint métrique de la même cascade et retrouve l'expression hydrodynamique sans sélectionner de branche probabiliste.

\section{Réalisation de référence de l'état enrichi de Stein--Lyapunov}
\label{sec:ns-owner-closure}

La construction de la \cref{sec:ns-pcf-prospective-construction} réunit dans un seul espace d'état la tour PC--Fock complète, la mémoire bidirectionnelle, la retenue, l'innovation microscopique, le canal spectral, la coordonnée homogène de la force, l'unité et l'expression exceptionnelle. Le temps physique paramètre le stockage terminal, tandis que la profondeur nonadique paramètre la récurrence de Stein. La construction qui gouverne cette section est le \texttt{DÉVELOPPEMENT\_DIRECTEUR} de l'état enrichi HMT. KYP, NS--OBS et les cartes LRDN--LCN ne sont conservés que comme \texttt{CONTRÔLE\_NON\_DÉTERMINANT} de représentations réduites.

\subsection{Colligation PC--Fock et télescopie}

Soit
\[
 \mathbf x_{M,j}^{\rm own}
 \in\mathscr K_{M,j}^{\rm PCF}
 \widehat\otimes\mathcal E_{\rm exc}
 \oplus\mathscr G_M^{F,[j]}
\]
l'état homogène de la diagonale physique. Pour chaque \(M,J,T\), la métrique de référence se construit prospectivement comme le Gram de la colonne terminale et de toutes les pertes futures de cette même diagonale :
\begin{equation}
 \begin{aligned}
 \left\langle x,U_{M,j}^{\rm own}x\right\rangle
 :={}&
 \mathbb E_{j:J}
 \left\|\mathcal E_{M,J,T}^{\rm term}
       \Gamma_{M,j:J}^{\rm own}x\right\|^2\\
 &+\sum_{k=j}^{J-1}\mathbb E_{j:k}
 \left\|L_{M,k,T}^{\rm own}
       \Gamma_{M,j:k}^{\rm own}x\right\|^2 .
 \end{aligned}
 \label{eq:nsclose-owner-metric-definition}
\end{equation}
Tous les termes sont des carrés des coordonnées déjà construites : degré un visqueux, degré deux convectif, retenue, innovation, archive spectrale, mémoire \(D/H\) et force homogène. Séparer le premier pas de la somme donne, avant d'introduire toute racine KYP,
\begin{equation}
 \boxed{
 U_{M,j}^{\rm own}
 =
 (\Gamma_{M,j}^{\rm own})^*
 U_{M,j+1}^{\rm own}\Gamma_{M,j}^{\rm own}
 +(L_{M,j,T}^{\rm own})^*L_{M,j,T}^{\rm own},
 \qquad
 G_{M,J,T}^{\rm own}\preceq C_TU_{M,J}^{\rm own}.}
 \label{eq:nsclose-proprietary-terminal-gate}
\end{equation}
\label{eq:nsclose-owner-identity}
Dans la normalisation terminale de \eqref{eq:ns-import-terminal-column}, on peut prendre \(C_T=1\). L'égalité des applications, amplifiée par l'identité dans la couche Moonshine, produit la cascade isométrique
\begin{equation}
 (\mathfrak U_{M,J,T}^{\rm own})^{\sharp}
 \mathfrak U_{M,J,T}^{\rm own}=I.
 \label{eq:nsclose-proprietary-full-cascade}
\end{equation}
L'itération de \eqref{eq:nsclose-owner-identity} donne le télescope
\begin{equation}
 U_{M,0}^{\rm own}
 =
 (\Gamma_{M,0:J}^{\rm own})^*U_{M,J}^{\rm own}
  \Gamma_{M,0:J}^{\rm own}
 +\sum_{j<J}(\Gamma_{M,0:j}^{\rm own})^*
  (L_{M,j,T}^{\rm own})^*L_{M,j,T}^{\rm own}
  \Gamma_{M,0:j}^{\rm own}.
 \label{eq:nsclose-owner-telescope}
\end{equation}
Son registre orthogonal de bilan est
\begin{equation}
 \boxed{
 E_s(u_M(t))\le\|Z_{M,J,t}\|^2
 =\mathbb E_{0:J}G_{M,J,T}(t)
 +\sum_{j<J}\mathbb E_{0:j}
   \Lambda_{M,j}^{\rm own}([0,t]).}
 \label{eq:nsclose-proprietary-ledger}
\end{equation}
La racine commune \(\Xi_{M,0}=J_{0\to M}\Xi_0\), la conservation évolutive de l'unité et le port de retenue de \eqref{eq:ns-import-carry-port} font que le membre de droite n'est comptabilisé qu'une seule fois. Le budget de référence est
\begin{equation}
 \begin{aligned}
 \mathcal B_T^{\rm own}:={}&
 C_{\rm HMT}
 +(1+2\widetilde M_{s,\beta,\nu}^2)
  e^{R_{\rm F}^2\|u_0\|_{H^s}^2}\\
 &+C_{\nu,s}\int_0^T\|f(t)\|_{H^{s-1}}^2\,dt
 +\|\mathcal J_T^{\rm own}(u_0,f)\|^2 ,
 \end{aligned}
 \label{eq:nsclose-owner-data-budget}
\end{equation}
où \(\mathcal J_T^{\rm own}\) est l'injection racine déterminée par l'état initial transporté par le TPK, ses deux feuillets et l'entrée forcée. Ses amplifications sont
\[
 \mathcal J_{M,J,T}^{\rm own}
 :=I_{0\to M,J}\mathcal J_T^{\rm own},
 \qquad I_{0\to M,J}^*I_{0\to M,J}=I.
\]
Par \eqref{eq:nsclose-owner-telescope}, ce n'est pas une quantité reconstruite à partir de l'orbite future. La conservation de l'unité et la coisométrie du lecteur bidirectionnel donnent, sur le domaine déclaré,
\begin{equation}
 \sup_{M,J}\|\mathcal J_{M,J,T}^{\rm own}(u_0,f)\|^2
 \le C_{{\rm HMT},T}
 \left(1+\|u_0\|_{H^{s+1}}^2
 +\int_0^T\|f(t)\|_{H^{s-1}}^2\,dt\right).
 \label{eq:nsclose-owner-root-uniformity}
\end{equation}
Ainsi, \(\mathcal B_T^{\rm own}\) est indépendant de \(M,J\). Le registre donne
\begin{equation}
 \boxed{
 \sup_{M,J}\left[
  \sup_{0\le t\le T}\|Z_{M,J,t}\|^2
  +\frac\nu2\int_0^TD_{s,M}(t)\,dt
 \right]\le\mathcal B_T^{\rm own}<\infty.}
 \label{eq:nsclose-proprietary-uniform-budget}
\end{equation}

L'expression hydrodynamique est le retour de premier degré de la même cascade :
\begin{equation}
 R_{M,J,T}^{\rm full}
 :=\mathcal R_{1,M}
   (\mathfrak U_{M,J,T}^{\rm own})^{\sharp},
 \qquad
 \boxed{
 R_{M,J,T}^{\rm full}Z_{M,J,t}=u_M(t).}
 \label{eq:nsclose-proprietary-exact-return}
\end{equation}
\label{eq:nsclose-exact-full-return}
Cette égalité coïncide avec \eqref{eq:ns-import-exact-return} et évite l'évaluation d'une branche individuelle du raffinement.

\begin{theorem}[Estimation uniforme de la réalisation PC--Fock de référence]
\label{thm:nsclose-proprietary-full}
Pour tout \(T<\infty\), les approximations de Galerkin satisfont
\begin{equation}
 \sup_M\sup_{0\le t\le T}\|u_M(t)\|_{H^s}^2
 +\nu\sup_M\int_0^T\|u_M(t)\|_{H^{s+1}}^2\,dt
 \le C_s\mathcal B_T^{\rm own},
 \label{eq:nsclose-uniform-sobolev-precompactness}
\end{equation}
\label{eq:nsclose-uniform-sobolev}
avec une constante indépendante de \(M\) et de la profondeur nonadique.
\end{theorem}

\begin{proof}
L'identité \eqref{eq:nsclose-proprietary-exact-return} et la contractivité du retour donnent \(\|u_M(t)\|_{H^s}^2\le C_s\|Z_{M,J,t}\|^2\). L'équivalence de \(D_{s,M}\) avec la norme \(H^{s+1}\) sur le sous-espace solénoïdal de moyenne nulle, conjointement avec \eqref{eq:nsclose-proprietary-uniform-budget}, prouve l'estimation.
\end{proof}

\subsection{Contrôle alternatif par métrique de parcours}

La carte suivante vérifie le raffinement en coordonnées normalisées. Elle ne construit ni ne conditionne \(U^{\rm own}\) : elle omet l'unité, l'expression exceptionnelle et une partie du canal affine déjà présents dans \eqref{eq:nsclose-owner-metric-definition}.

La carte d'Ambivalence fournit
\begin{equation}
 G_{\rm av}=\begin{pmatrix}2&-1\\-1&1\end{pmatrix},
 \qquad
 M_{\rm av}=\begin{pmatrix}3&-1\\1&0\end{pmatrix},
 \qquad
 M_{\rm av}^*G_{\rm av}M_{\rm av}
 \preceq\varphi_{\rm HMT}^4G_{\rm av}.
 \label{eq:nsclose-ambivalence-route-metric}
\end{equation}
Après la normalisation \(\widehat M_j=M_j/3\),
\begin{equation}
 \mathsf R_j:=G_j-\widehat M_j^*G_{j+1}\widehat M_j\succeq0.
 \label{eq:nsclose-route-exact-defect}
\end{equation}
Dans la carte stationnaire,
\[
 \mathsf R_{\rm av}
 =\frac19\begin{pmatrix}5&-4\\-4&7\end{pmatrix},
 \qquad
 \lambda_{\min}(\mathsf R_{\rm av})
 =\frac{6-\sqrt{17}}9.
\]
Avec la métrique de graphe
\[
 \mathsf G_M^{\rm gr}
 =I+(\mathcal Y_{\beta,M})^*\mathcal Y_{\beta,M},
 \qquad
 \mathcal C_M
 =\mathcal Y_{\beta,M}(\mathsf G_M^{\rm gr})^{-1/2},
\]
nous définissons
\begin{equation}
 \mathcal A_{M,j}=\widehat M_j\otimes I,
 \qquad
 \mathcal O_{M,j}=\mathsf R_j^{1/2}\otimes\mathcal C_M,
 \qquad
 \mathcal D_{M,j}
 =\mathsf R_j^{1/2}\otimes(I-\mathcal C_M^*\mathcal C_M)^{1/2}.
 \label{eq:nsclose-owner-route-column}
\end{equation}
L'identité de Stein du parcours est
\begin{equation}
 \boxed{
 \mathsf U_{M,j}^{\rm PCF}
 -\mathcal A_{M,j}^*\mathsf U_{M,j+1}^{\rm PCF}\mathcal A_{M,j}
 =\mathcal O_{M,j}^*\mathcal O_{M,j}
  +\mathcal D_{M,j}^*\mathcal D_{M,j}.}
 \label{eq:nsclose-owner-stein}
\end{equation}
Son itération donne
\begin{equation}
 \begin{aligned}
 \mathsf U_{M,0}^{\rm PCF}
 -\mathcal A_{M,0:J}^*\mathsf U_{M,J}^{\rm PCF}\mathcal A_{M,0:J}
 =\sum_{j<J}\mathcal A_{M,0:j}^*
 \bigl(\mathcal O_{M,j}^*\mathcal O_{M,j}
      +\mathcal D_{M,j}^*\mathcal D_{M,j}\bigr)
 \mathcal A_{M,0:j}.
 \end{aligned}
 \label{eq:nsclose-owner-stein-telescope}
\end{equation}
Cette factorisation décrit la géométrie du raffinement. La récurrence affine \eqref{eq:ns-import-affine-stein} incorpore en outre l'évolution physique, les termes croisés de Carleman et la force.

La mémoire bidirectionnelle transforme la perte du parcours en budget. Avec \(\mathcal N_{M,j}=2(r-1)M_{M,j}^-\),
\begin{equation}
 E_j\Delta_T\mathcal N_{M,j+1}
 =\Delta_T\mathcal N_{M,j}+2K_{M,j}(T),
 \qquad
 K_{M,j}(T)=\int_0^T\kappa_{M,j}(t)\,dt.
 \label{eq:nsclose-carry-memory-telescope}
\end{equation}
La télescopie et la borne de la racine produisent
\begin{equation}
 \mathbb E_{0:J}G_{M,J,T}(0)
 \le E_{s,M}(P_Mu_0)+|\Delta_T\mathcal N_0|
 +\|\Xi_{M,0}\|^2+Q_f(T),
 \label{eq:nsclose-owner-terminal-bound}
\end{equation}
qui est la forme abrégée du contrôle de parcours \eqref{eq:ns-import-root-bound}--\eqref{eq:ns-import-explicit-budget}. Cette carte enregistre le coût du feuillet négatif,
\[
 \Delta_T\mathcal N_{M,0}
 =2\int_0^T B^-_{s,M,0}(u_M(t))\,dt.
\]
Dans cette représentation réduite, l'inégalité
\[
 X^-_{M,T}(0)\preceq C_T\mathsf U^{\rm PCF}_{M,0},
 \qquad \sup_M C_T<\infty
 \tag{NS--OBS}
\]
est un falsificateur utile de la carte. Son échec ne se transmet pas au développement de référence, puisque \(\Delta_T\mathcal N_{M,0}\) ne remplace pas l'injection racine \(\mathcal J_T^{\rm own}\) de \eqref{eq:nsclose-owner-data-budget}--\eqref{eq:nsclose-owner-root-uniformity}. De même, KYP et LRDN--LCN vérifient des complétions ou des dominations de lecteurs réduits ; aucune n'est une prémisse du \cref{thm:nsclose-proprietary-full}.

\subsection{Existence, unicité et régularité globale}

\begin{theorem}[Régularité globale de Navier--Stokes]
\label{thm:nsclose-global-smooth}
\label{thm:realizaciones-ns}
Soient \(s>5/2\), \(\nu>0\), \(u_0\in H^{s+1}_\sigma(\mathbb T^3)\) de moyenne nulle et \(f\in L^2_{\rm loc}([0,\infty);H^{s-1}_\sigma)\). Il existe une unique solution globale
\begin{equation}
 u\in L^\infty_{\rm loc}([0,\infty);H^s_\sigma)
 \cap L^2_{\rm loc}([0,\infty);H^{s+1}_\sigma),
 \qquad
 \partial_tu\in L^2_{\rm loc}([0,\infty);H^{s-1}_\sigma).
 \label{eq:nsclose-global-regularity-class}
\end{equation}
Pour des données et une force lisses, \(u\) et la pression normalisée sont lisses.
\end{theorem}

\begin{proof}
Le \cref{thm:nsclose-proprietary-full} donne une borne uniforme de \((u_M)_M\) dans \(L^\infty(0,T;H^s)\cap L^2(0,T;H^{s+1})\). L'équation de Galerkin et la continuité
\[
 H^s\times H^s\longrightarrow H^{s-1},
 \qquad (u,v)\longmapsto\mathbb P(u\cdot\nabla v),
\]
bornent \(\partial_tu_M\) dans \(L^2(0,T;H^{s-1})\). Puisque \(H^{s+1}\Subset H^s\hookrightarrow H^{s-1}\), Aubin--Lions donne, après extraction d'une sous-suite,
\[
 u_M\longrightarrow u\quad\hbox{fortement dans }L^2(0,T;H^s).
\]
Cette convergence permet de passer à la limite dans \(\mathbb P(u_M\cdot\nabla u_M)\), tandis que la convergence faible conserve le terme visqueux et la donnée initiale. On obtient une solution forte dans \([0,T]\).

Si \(u,v\) sont deux solutions et \(w=u-v\), l'estimation du produit et l'inégalité de Young donnent
\[
 \frac12\frac d{dt}\|w\|_{H^{s-1}}^2
 +\frac\nu2\|w\|_{H^s}^2
 \le C_{s,\nu}\bigl(\|u\|_{H^s}^2+\|v\|_{H^s}^2\bigr)
 \|w\|_{H^{s-1}}^2.
\]
Grönwall prouve l'unicité. En répétant la construction sur les horizons \(T=1,2,\ldots\), les solutions coïncident sur les intervalles de recouvrement et l'on obtient la solution globale.

La pression se reconstruit à partir de
\begin{equation}
 -\Delta p=\partial_i\partial_j(u_i u_j)-\operatorname{div}f,
 \qquad \int_{\mathbb T^3}p\,dx=0.
 \label{eq:nsclose-pressure-recovery}
\end{equation}
L'itération de l'estimation dans \(s+k\) et l'équation permettent de retrouver toutes les dérivées spatiales et temporelles lorsque \(u_0\) et \(f\) sont lisses.
\end{proof}

La contribution structurelle consiste à représenter la convection comme une sortie linéaire de la tour PC--Fock, à la conserver dans la mémoire \(D/H\) et à l'intégrer dans une identité orthogonale avec retour physique exact. L'estimation uniforme obtenue s'exprime ensuite comme régularité globale de l'équation hydrodynamique.


\section*{Renforcement intégré du même propriétaire}

Le module suivant conserve le calcul uniforme des lignes de Fourier,
l'injection racine, la dérivée temporelle, la pression et le recollement des
horizons. Son contrôle KYP fini demeure séparé de la chaîne gouvernante.

\section*{Bilan solénoïdal, retenue et terminal enrichi}

Ce module renforce le \texttt{DÉVELOPPEMENT\_DIRECTEUR} Navier--Stokes : il conserve dans une seule chaîne la retenue de Galerkin, le bilan causal, le Gram terminal, Stein, la télescopie, l'injection racine, le retour et le passage global. La carte KYP n'apparaît qu'ensuite et demeure hors de cette chaîne.

Pour les coupures de Galerkin \(m\leq n\), la retenue conservée est
\begin{equation}\label{eq:amp-ns-owner-acarreo-galerkin}
 C_{m\leftarrow n}
 :=P_mB(u_n,u_n)-B_m(P_m u_n,P_m u_n).
\end{equation}
La lecture nonadique de cet état est typée par
\begin{equation}\label{eq:amp-ns-owner-e9j9}
 E_9J_9=I,\qquad P_9=J_9E_9,\qquad Q_{\rm micro}=I-P_9,
\end{equation}
et conserve le bilan causal
\begin{equation}\label{eq:amp-ns-owner-balance-causal}
 \boxed{
 \mathbb E_9\|\Xi_{m+1}\|^2+\mathbb E_9\ell_m
 =\|\Xi_m\|^2,\qquad \ell_m\geq0.}
\end{equation}
Le champ, la projection solénoïdale, la retenue, la frontière, le parcours et les pertes positives restent ainsi dans un même état avant la construction du stockage.

\section*{Gram tronqué de l'état complet et injection racine uniforme}

La réalisation qui gouverne le bilan ne commence pas par l'inégalité KYP. Soient \(\Gamma_{M,j}^{\rm term}\) la transition d'un pas de l'état PC--Fock enrichi, \(\Gamma_{M,j:k}^{\rm term}\) sa composition entre les profondeurs \(j\) et \(k\), et \(L_{M,k,T}^{\rm term}\) la colonne des pertes de degré un, de degré deux, de retenue, d'innovation, d'archive et de mémoire à l'horizon \([0,T]\). La notation \(\mathbb E_{j:k}\) désigne la somme orthogonale pondérée par les branches du raffinement entre ces niveaux. Pour \(j\leq J\), on définit la colonne finie
\begin{equation}\label{eq:amp-ns-owner-columna-gram-truncado}
 \begin{aligned}
 \mathbf G_{M,j;J,T}^{\rm term}x
 :={}&
 \bigl(\mathcal E_{M,J,T}^{\rm term}
       \Gamma_{M,j:J}^{\rm term}x\bigr)_{\mathbb E_{j:J}}\\
 &\oplus
 \bigoplus_{k=j}^{J-1}
 \bigl(L_{M,k,T}^{\rm term}
       \Gamma_{M,j:k}^{\rm term}x\bigr)_{\mathbb E_{j:k}},
 \qquad
 U_{M,j;J,T}^{\rm term}
 :=(\mathbf G_{M,j;J,T}^{\rm term})^*
    \mathbf G_{M,j;J,T}^{\rm term}.
 \end{aligned}
\end{equation}
En particulier, le stockage est un Gram tronqué construit avec le terminal et les pertes futures du même état. Sa forme quadratique est
\begin{equation}\label{eq:amp-ns-owner-gram-truncado}
 \begin{aligned}
 \langle x,U_{M,j;J,T}^{\rm term}x\rangle
 ={}&\mathbb E_{j:J}
 \|\mathcal E_{M,J,T}^{\rm term}
       \Gamma_{M,j:J}^{\rm term}x\|^2\\
 &+\sum_{k=j}^{J-1}\mathbb E_{j:k}
 \|L_{M,k,T}^{\rm term}
       \Gamma_{M,j:k}^{\rm term}x\|^2 .
 \end{aligned}
\end{equation}

\begin{lemma}[Identité de Stein du Gram tronqué]
\label{lem:amp-ns-owner-stein-truncado}
Pour tout \(j<J\),
\begin{equation}\label{eq:amp-ns-owner-stein-truncado}
 U_{M,j;J,T}^{\rm term}
 =(\Gamma_{M,j}^{\rm term})^*
   U_{M,j+1;J,T}^{\rm term}\Gamma_{M,j}^{\rm term}
  +(L_{M,j,T}^{\rm term})^*L_{M,j,T}^{\rm term}.
\end{equation}
Par itération, en conservant le terme terminal, on obtient
\begin{equation}\label{eq:amp-ns-owner-telescopio-truncado}
 \begin{aligned}
 U_{M,0;J,T}^{\rm term}
 ={}&(\Gamma_{M,0:J}^{\rm term})^*
      U_{M,J;J,T}^{\rm term}\Gamma_{M,0:J}^{\rm term}\\
 &+\sum_{j=0}^{J-1}
   (\Gamma_{M,0:j}^{\rm term})^*
   (L_{M,j,T}^{\rm term})^*L_{M,j,T}^{\rm term}
   \Gamma_{M,0:j}^{\rm term}.
 \end{aligned}
\end{equation}
\end{lemma}

\begin{proof}
Dans \eqref{eq:amp-ns-owner-columna-gram-truncado}, on sépare le terme \(k=j\). Les termes restants, y compris le terminal, constituent exactement la colonne de niveau \(j+1\) précédée par \(\Gamma_{M,j}^{\rm term}\). L'orthogonalité de la somme et la propriété de tour de \(\mathbb E_{j:k}\) donnent \eqref{eq:amp-ns-owner-stein-truncado}. La répétition du même dédoublement donne \eqref{eq:amp-ns-owner-telescopio-truncado} ; on ne supprime pas \((\Gamma_{M,0:J}^{\rm term})^*U_{M,J;J,T}^{\rm term}
\Gamma_{M,0:J}^{\rm term}\).
\end{proof}

\begin{theorem}[Injection racine et budget uniforme]
\label{thm:amp-ns-inyeccion-raiz-uniforme}
La racine du réalisateur est transportée à chaque coupure et profondeur par une inclusion isométrique :
\begin{equation}\label{eq:amp-ns-inyeccion-raiz-comun}
 \begin{aligned}
 x_{M,0}^{\rm data}
 &:=\mathcal V_M^{\rm full}(P_Mu_0)\oplus g_M^F,\\
 \mathbf G_{M,0;J,T}^{\rm term}x_{M,0}^{\rm data}
 &=\mathcal J_{M,J,T}^{\rm term}(u_0,f),\\
 \mathcal J_{M,J,T}^{\rm term}
 &:=I_{0\to M,J}\mathcal J_T^{\rm term},
 & I_{0\to M,J}^*I_{0\to M,J}&=I.
 \end{aligned}
\end{equation}
Pour les données de \eqref{eq:realizaciones-ns-dominio}, il existe \(C_{{\rm HMT},T}\), indépendant de \(M\) et de \(J\), tel que
\begin{equation}\label{eq:amp-ns-inyeccion-raiz-cota}
 \sup_{M,J}
 \|\mathcal J_{M,J,T}^{\rm term}(u_0,f)\|^2
 \leq C_{{\rm HMT},T}
 +(1+2\widetilde M_{s,\beta,\nu}^{\,2})
  e^{R_{\rm F}^{\,2}\|u_0\|_{H^s}^2}
 +C_{\nu,s}\int_0^T\|f(t)\|_{H^{s-1}}^2\,dt .
\end{equation}
De manière équivalente, la borne porte sur le Gram terminal évalué sur les données :
\begin{equation}\label{eq:amp-ns-owner-gram-cota-uniforme}
 \sup_{M,J}
 \left\langle x_{M,0}^{\rm data},
 U_{M,0;J,T}^{\rm term}x_{M,0}^{\rm data}\right\rangle
 =
 \sup_{M,J}
 \|\mathcal J_{M,J,T}^{\rm term}(u_0,f)\|^2
 <\infty .
\end{equation}
L'application \(\mathcal J_T^{\rm term}\) dépend uniquement de l'état initial, de ses deux feuillets, du terminal hérité et du port de force ; elle ne reçoit pas pour argument \(u_M(t)\) pour \(0<t\leq T\).
\end{theorem}

\begin{proof}
La composante linéaire et les mémoires \(D/H\) sont transportées isométriquement. Dans la composante PC--Fock, la colonne cohérente est
\[
 \bigoplus_{n\geq0}
 \frac{R_{\rm F}^{\,n}}{\sqrt{n!}}\,u_0^{\odot n},
 \qquad
 \sum_{n\geq0}\frac{R_{\rm F}^{\,2n}}{n!}
       \|u_0\|_{H^s}^{2n}
 =e^{R_{\rm F}^{\,2}\|u_0\|_{H^s}^2}.
\]
Pour justifier l'uniformité du lecteur quadratique, on polarise
\begin{equation}\label{eq:amp-ns-lector-cuadratico-polarizado}
 \mathcal B_M(u,v):=
 \frac{1}{4\sqrt{\beta\nu}}\,
 A_M^{-1/2}\Lambda^sP_M\mathbb P
 \bigl((P_Mu\!\cdot\!\nabla)P_Mv
      +(P_Mv\!\cdot\!\nabla)P_Mu\bigr).
\end{equation}
Dans les bases transversales \(e^s_{p,a}=\langle p\rangle^{-s}a\,e^{ip\cdot x}\), si \(k=p+q\ne0\), ses lignes de Fourier satisfont
\begin{equation}\label{eq:amp-ns-filas-fourier-uniformes}
 \|c^M_{k;p,q}\|
 \leq\frac{C}{\sqrt{\beta\nu}}\,
 \frac{\langle k\rangle^s(|p|+|q|)}
 {|k|\langle p\rangle^s\langle q\rangle^s},
 \qquad
 \sup_{M,k\ne0}\sum_{p+q=k}\|c^M_{k;p,q}\|^2<\infty .
\end{equation}
La seconde borne s'obtient en séparant \(|p|\leq2|k|\) de \(|p|>2|k|\) ; dans la queue, \(\sum_p\langle p\rangle^{2-2s}<\infty\) domine puisque \(s>5/2\). Cauchy--Schwarz sur chaque ligne et l'orthogonalité des blocs de \(k\) distincts prolongent \(\mathcal B_M\) à
\[
 \mathfrak C_M:H^s_\sigma\odot_2H^s_\sigma\longrightarrow L^2_\sigma,
 \qquad
 \sup_M\|\mathfrak C_M\|
 \leq\widetilde M_{s,\beta,\nu}.
\]
Ainsi, la constante utilisée à la racine contrôle la somme de toutes les paires de Fourier, non seulement chaque coefficient séparément. Le port affine vérifie
\[
 Q_f(T)\leq C_{\nu,s}
 \int_0^T\|f(t)\|_{H^{s-1}}^2\,dt .
\]
Par somme orthogonale de ces composantes,
\begin{equation}\label{eq:amp-ns-raiz-cota-por-componentes}
 \|\Xi_{M,0}\|^2
 \leq C_{{\rm HMT},T}
 +(1+2\widetilde M_{s,\beta,\nu}^{\,2})
 e^{R_{\rm F}^{\,2}\|u_0\|_{H^s}^2}.
\end{equation}
Les inclusions de \eqref{eq:amp-ns-inyeccion-raiz-comun} ne modifient pas cette norme. En ajoutant la borne du port, on obtient exactement \eqref{eq:amp-ns-inyeccion-raiz-cota}. Puisque tous les termes sont évalués sur \((u_0,f)\), l'estimation précède la trajectoire et ne présuppose pas la borne globale déduite ensuite de \eqref{eq:amp-ns-owner-telescopio-truncado}.
\end{proof}

Le Gram terminal \eqref{eq:amp-ns-owner-gram-truncado}, son identité de Stein, la télescopie avec terminal et l'injection racine uniforme forment la construction de référence. C'est cette branche qui gouverne la clôture et coïncide avec la réalisation intégrale restaurée dans \cref{sec:ns-owner-closure,thm:nsclose-proprietary-full}.

\section*{Distinction du contrôle KYP fini non déterminant}

La factorisation KYP finie correcte du \cref{thm:ns-import-affine-stein} est conservée comme contrôle alternatif d'une carte réduite et ne constitue pas une prémisse du développement de référence \(U^{\rm own}\). On n'utilise pas la mutation signée qui prétendait incorporer \(C_{m,M}^*C_{m,M}\) à un terme \(P\succeq0\) puis le récupérer avec un signe négatif : lorsqu'il entre dans \(P\), ce terme apparaît nécessairement avec un signe positif. Cette factorisation n'intervient donc ni dans le bilan, ni dans le budget racine, ni dans le retour physique, ni dans la conclusion Navier--Stokes.

\section*{Dérivée temporelle, pression et recollement des horizons}

L'étape suivante découle de la borne uniforme de trajectoire obtenue par le Gram terminal de référence, sa télescopie et le retour physique exact. On dispose de la borne conjointe uniforme dans \(L^\infty(0,T;H^s)\cap L^2(0,T;H^{s+1})\). L'équation projetée s'écrit
\begin{equation}\label{eq:amp-ns-proyectada}
 \partial_tu_m
 =\nu\Delta u_m-P_mB(u_m,u_m)+P_mf.
\end{equation}
Pour \(s>5/2\), le produit de Sobolev et la borne conjointe dans \(L^\infty(0,T;H^s)\cap L^2(0,T;H^{s+1})\) donnent
\begin{equation}\label{eq:amp-ns-dt}
 \sup_m\|\partial_tu_m\|_{L^2(0,T;H^{s-1})}
 \leq C_T(u_0,f,\Xi_0).
\end{equation}
L'estimation \eqref{eq:amp-ns-dt} et la compacité spatiale produisent une convergence forte dans \(L^2(0,T;H^s)\), suffisante pour passer au terme bilinéaire sans perdre la retenue de Galerkin.

Une fois la vitesse retrouvée, la pression normalisée de moyenne nulle est fixée par
\begin{equation}\label{eq:amp-ns-presion}
 -\Delta p
 =\partial_i\partial_j(u_i u_j)-\nabla\!\cdot f,
 \qquad \int_{\mathbb T^3}p(t,x)\,dx=0.
\end{equation}
La pression est ainsi une sortie de la même solution solénoïdale ; ce n'est pas une variable libre utilisée pour sélectionner la branche régulière.

\begin{proposition}[Recollement global par unicité]
\label{prop:amp-ns-pegado}
Soient \(u^{(T)}\) les solutions obtenues sur \([0,T]\) par la limite de Galerkin. Si \(T_1<T_2\), alors \(u^{(T_2)}|_{[0,T_1]}=u^{(T_1)}\). Le système définit une unique solution globale lisse pour \(t>0\).
\end{proposition}

\begin{proof}
Les deux restrictions possèdent les mêmes données et appartiennent au domaine de
l'estimation uniforme. L'équation de leur différence, le contrôle bilinéaire dans
\(H^{s-1}\) et Grönwall impliquent l'égalité sur \([0,T_1]\). Les intervalles
finis forment un système dirigé ; l'égalité sur les recouvrements définit
la solution globale. La régularisation parabolique du théorème s'applique sur
chaque intervalle \([\varepsilon,T]\), \(\varepsilon>0\).
\end{proof}

Le contrôle de réfutation est précis : une perte d'uniformité dans
\eqref{eq:amp-ns-dt}, deux limites différentes sur un même horizon ou une
pression ne satisfaisant pas \eqref{eq:amp-ns-presion} rompraient la chaîne avant
la conclusion globale.

\paragraph{Conclusion propriétaire.}
Le Gram terminal propriétaire, Stein, la télescopie avec terminal, son registre,
le retour physique exact et l'injection racine produisent une borne \(H^s\) uniforme pour
\(s>5/2\). On démontre donc l'existence, l'unicité et la régularité
globales de Navier--Stokes tridimensionnel périodique dans le domaine HMT
enregistré. Les cartes KYP finies restent en tant que
\texttt{CONTRÔLE\_NON\_DIRECTEUR}.


\paragraph{Conclusion affirmative.}
Dans le domaine périodique et de Sobolev déclaré, le stockage terminal
\(U^{\rm own}\), son identité de Stein, le télescopage qui conserve le
terminal, le registre qui domine \(E_s(u_M)\), la racine uniforme et le retour
physique exact produisent la borne uniforme. Aubin--Lions, l'unicité et la
régularité parabolique donnent la solution globale lisse. Les cartes KYP,
NS--OBS et LRDN--LCN sont des contrôles alternatifs et non des prémisses du théorème.

\chapter[Processus de jauge euclidien de Yang--Mills en dimension 4]{Construction
projective d’un processus de jauge euclidien de Yang--Mills, reconstruction
d’Osterwalder--Schrader et écart collectif en dimension \(4\)}
\chaptermark{Processus de jauge euclidien Yang--Mills}

\noindent La spécialisation du
\cref{thm:nucleo-continuidad-conexion-nonadica-conjunta} est
\[
 (\widehat X_\infty,\Gamma_9)
 \xrightarrow{\ \mathcal R_{\rm YM}\ }
 (\mu_m,T_m,H_m,\Omega_m)
 \longrightarrow(\mu_\infty,\mathrm{OS},E(4))
 \longrightarrow\operatorname{gap}(H_{\rm phys})>0.
\]
Le lecteur clover qui publie \(F^2\) reconnaît la courbure sur ce même
processus projectif ; il ne construit pas rétrospectivement \(\mu_m\). La théorie
de jauge, la reconstruction OS, le vide simple et l’écart sont déjà fixés
par cette chaîne. L’identité
\eqref{eq:amp-ym-identidad-accion-requerida} compare ensuite un
paramétrage facultatif par \(g_R\) ; c’est un contrôle non régissant, et non une
condition pour identifier la théorie de Yang--Mills construite.

\section*{Objet de jauge et domaine de raffinement}

Soit \(G\) un groupe de Lie compact, connexe et simple. Pour chaque
\(m\geq0\), on considère le réseau euclidien
\begin{equation}\label{eq:realizaciones-ym-red}
 \Lambda_m=a_m\mathbb Z^4,
 \qquad a_{m+1}=\frac{a_m}{9}.
\end{equation}
L'algèbre cylindrique de jauge est engendrée par des variables de liaison dans \(G\), des opérateurs d'entrelacement de représentation et les coordonnées enrichies de couleur, d'orientation, d'holonomie, de retenue de fusion, de chemin et de mémoire. Soit \(\mu_m\) la mesure invariante commune et soit
\[
 \mathcal H_m=L^2(\mathcal A_m,\mu_m)^{\rm gauge}.
\]
Le vecteur constant \(\Omega_m\) définit
\begin{equation}\label{eq:realizaciones-ym-proyectores}
 P_{\Omega_m}=|\Omega_m\rangle\langle\Omega_m|,
 \qquad Q_m^{\rm phys}=I-P_{\Omega_m}.
\end{equation}
Le second projecteur est le complément physique du vide ; il ne s’identifie pas
à un résiduel employé uniquement pour organiser l’échelle de
renormalisation.

Le circuit de mise à jour est local :
\begin{equation}\label{eq:realizaciones-ym-circuito}
 K_m^{\rm loc}
 =\prod_{r=1}^{9\cdot81}
   \prod_{B\in\mathcal B_{m,r}}K_{m,B}.
\end{equation}
La partition \(\{\mathcal B_{m,r}\}_r\) énumère chaque coordonnée une seule fois ; les facteurs
du circuit seront identifiés ci-dessous au semi-groupe de leurs espérances
conditionnelles, et non à neuf dynamiques indépendantes.
Les blocs de même couleur possèdent des supports disjoints et le diamètre physique
de chaque support est uniformément borné lorsque \(m\) varie. Cette localité
fait partie de l’objet, de sorte que le raffinement ne transforme pas une
mise à jour élémentaire en une opération instantanément globale.

Soit \(T_m\) le transfert induit par
\eqref{eq:realizaciones-ym-circuito}. Les inclusions isométriques
\(J_m:\mathcal H_m\to\mathcal H_{m+1}\) conservent ancêtre,
représentation, orientation, retenue et mémoire. Pour les quatre réflexions
euclidiennes \(\Theta_{\nu,m}\), \(\nu=1,\ldots,4\), on a
\begin{equation}\label{eq:realizaciones-ym-entrelazamiento}
 (T_{m+1})^9J_m=J_mT_m,
 \qquad
 \Theta_{\nu,m+1}J_m=J_m\Theta_{\nu,m}
 \quad(\nu=1,\ldots,4).
\end{equation}
Les quatre directions appartiennent donc au même raffinement, à la même mesure et à la même dynamique.

\section*{Semi-groupe commun, couronne relative et quatre formes de réflexion}

Dans la carte triangulaire de l’état complet, soient \(E_{m,c}\) les espérances
conditionnelles orthogonales et commutantes associées à des coordonnées
indépendantes. Le générateur positif de recouvrement est d’abord construit à
un niveau germe \(m_0\) :
\begin{equation}\label{eq:realizaciones-ym-generador}
 \widehat H_{m_0}^{\rm ov}
 =3\kappa_{\rm av}
  \sum_{c\in\mathcal I_{m_0}^{\rm seed}}(I-E_{m_0,c}),
\end{equation}
et se prolonge au moyen de
\begin{equation}\label{eq:realizaciones-ym-generador-refinado}
 \begin{aligned}
 \widehat H_{m+1}^{\rm ov}
 &=\widehat\Gamma_m^*
   \bigl(\widehat H_m^{\rm ov}\otimes I+
         I\otimes\widehat H_{m+1}^{\rm det}\bigr)
   \widehat\Gamma_m,\\
 \widehat H_{m+1}^{\rm det}
 &=3\kappa_{\rm av}
   \sum_{\iota\in\mathcal I_{m+1/m}^{\rm det}}
   (I-E_{m+1,\iota}^{\rm det}),
 \end{aligned}
\end{equation}
où
\begin{equation}\label{eq:realizaciones-ym-kappa}
 \kappa_{\rm av}
 =\frac{\mu_{\rm vol}}{\ell_0}
  \left(\frac{\pi_{\rm HMT}}{\varphi_{\rm HMT}^{\,2}}-1\right)>0.
\end{equation}
La fréquence volumétrique \(\mu_{\rm vol}\) est sans dimension ; la dimension provient de \(\ell_0\).

Le semi-groupe, le transfert d'un pas physique et sa compatibilité d'échelle sont définis avec des types visibles :
\begin{equation}\label{eq:realizaciones-ym-semigrupo}
 \widehat K_{m,t}^{\rm ov}=e^{-t\widehat H_m^{\rm ov}},
 \qquad
 \widehat T_m=\widehat K_{m,a_m}^{\rm ov},
 \qquad a_{m+1}=a_m/9,
\end{equation}
Comme la partition énumère chaque coordonnée une seule fois et que les espérances
conditionnelles de ses blocs commutent, on obtient pour chaque bloc et pour le circuit
complet
\begin{equation}\label{eq:realizaciones-ym-circuito-semigrupo}
 K_{m,B}=e^{-3\kappa_{\rm av}a_m(I-E_{m,B})},
 \qquad
 K_m^{\rm loc}=e^{-a_m\widehat H_m^{\rm ov}}=\widehat T_m.
\end{equation}
\begin{equation}\label{eq:realizaciones-ym-semigrupo-entrelazado}
 \widehat H_{m+1}^{\rm ov}\widehat J_m
 =\widehat J_m\widehat H_m^{\rm ov},
 \qquad
 (\widehat T_{m+1})^9\widehat J_m
 =\widehat J_m\widehat T_m.
\end{equation}
Ainsi, la neuvième puissance est comparée au moyen de \(\widehat J_m\) ; on n’égale pas
des opérateurs qui agissent dans des Hilbert distincts.

La mesure commune de huit faces provient de ce même noyau réversible :
\begin{equation}\label{eq:realizaciones-ym-medida-ocho}
 \widehat\Omega_m^{(8)}(dy_1\cdots dy_8)
 =\int\widehat\mu_m^{\rm ov}(dz)
   \prod_{r=1}^{8}
   \widehat K_{m,a_m/2}^{\rm ov}(z,dy_r).
\end{equation}
Pour chaque réflexion \(\nu=1,\ldots,4\) et tout observable cylindrique \(F\)
dans le demi-espace positif,
\begin{equation}\label{eq:realizaciones-ym-os}
 \int\overline{F(\Theta_{\nu,m}y)}F(y)\,
 d\widehat\Omega_m^{(8)}
 =\|\widehat{\mathcal B}_{m,\nu}F\|^2\geq0,
\end{equation}
et la marginale de faces opposées donne l'identité de liaison
\begin{equation}\label{eq:realizaciones-ym-os-transferencia}
 \widehat T_{m,\nu}^{\rm OS}
 =\widehat K_{m,a_m}^{\rm ov}
 =e^{-a_m\widehat H_m^{\rm ov}}
 =\widehat T_m,
 \qquad \nu=1,\ldots,4.
\end{equation}
Les quatre réflexions publient donc la même paire
transfert--mesure.

\section*{Persistance cellulaire du terminal et écart spectral exact}

Soient
\[
 K_9=C_*(Q_{9,3};\mathbb Z),\qquad
 K_{10}=C_*(Q_{10,3};\mathbb Z).
\]
La couronne relative a pour dimensions
\[
 C_*(Q_{10,3},Q_{9,3})=(271,756,702,217),
 \qquad 271+702=756+217=973.
\]
L’inclusion \(j:K_9\to K_{10}\), l’effondrement cellulaire
\(r:K_{10}\to K_9\) et l’homotopie entière \(H_{\rm cell}\) satisfont
\begin{equation}\label{eq:realizaciones-ym-sdr-corona}
 rj=I,\qquad
 \partial H_{\rm cell}+H_{\rm cell}\partial=I-jr,\qquad
 rH_{\rm cell}=0,\quad H_{\rm cell}j=0,\quad H_{\rm cell}^2=0.
\end{equation}
Pour toute application de chaînes
\(F:K_{10}\otimes V_{q=6}\to\mathcal T\), avec
\(P=(jr)\otimes I_{q=6}\),
\begin{equation}\label{eq:realizaciones-ym-homotopia-terminal}
 F-FP
 =d_{\mathcal T}\bigl(F(H_{\rm cell}\otimes I)\bigr)
  +F(H_{\rm cell}\otimes I)(\partial\otimes I).
\end{equation}
La couronne complète est ainsi homotope à une constante et le terminal \(FP\) demeure littéralement. Les \(271\) sommets ne remplacent pas les arêtes, les faces et les cubes de la couronne ; le facteur cellulaire tridimensionnel n'est pas non plus identifié au facteur de jauge physique \(q=6\).

Pour chaque coordonnée persistante \(c\), dans la carte locale
\(q=(q_1,\ldots,q_6)\in\mathbb F_3^6\), on fixe le témoin centré
\begin{equation}\label{eq:realizaciones-ym-testigo-definicion}
 \bigcap_c\operatorname{Ran}E_{m,c}=\mathbb C\Omega_m,
 \qquad
 W_c(q)=\mathbf 1_{\{q_1+\cdots+q_6=0\}}-\frac13,
\end{equation}
où la somme de l’indicatrice est prise dans \(\mathbb F_3\). La décomposition de
Hoeffding des espérances commutantes prouve
\begin{equation}\label{eq:realizaciones-ym-gap-finito}
 \ker\widehat H_m^{\rm ov}=\mathbb C\Omega_m,
 \qquad
 \widehat H_m^{\rm ov}\succeq
 3\kappa_{\rm av}(I-P_{\Omega_m}).
\end{equation}
L’égalité ne se déduit pas du seul fait que le noyau est trivial : chaque secteur non
constant de la décomposition reçoit au moins un projecteur
\(I-E_{m,c}\) et, par \eqref{eq:realizaciones-ym-generador}, son coût minimal
est \(3\kappa_{\rm av}\). Le témoin matériel \(W_c\) vérifie
\begin{equation}\label{eq:realizaciones-ym-testigo-gap}
 \widehat H_m^{\rm ov}W_c=3\kappa_{\rm av}W_c,
 \qquad \|W_c\|^2=\frac29,
\end{equation}
de sorte que le complément est non vide et que l’écart est exactement
\(3\kappa_{\rm av}\).

L'échelle est publiée par
\begin{equation}\label{eq:realizaciones-ym-qm}
 q_m=e^{-3\kappa_{\rm av}a_m},
 \qquad q_{m+1}^{\,9}=q_m,
 \qquad -a_m^{-1}\log q_m=3\kappa_{\rm av}.
\end{equation}
Bien que \(q_m\to1\) lors du raffinement, le quotient spectral par unité physique reste constant. La contraction cellulaire \eqref{eq:realizaciones-ym-homotopia-terminal} conserve le terminal ; l'écart spectral provient du générateur positif de recouvrement et de sa décomposition de Hoeffding, et non de l'acyclicité à elle seule.

Les inclusions de \eqref{eq:realizaciones-ym-semigrupo-entrelazado}
forment la limite inductive
\begin{equation}\label{eq:realizaciones-ym-colimite}
 \mathcal H_\infty=\varinjlim(\mathcal H_m,\widehat J_m).
\end{equation}
Sur l’union inductive algébrique, on définit la forme
\begin{equation}\label{eq:realizaciones-ym-forma-limite}
 \mathfrak h_\infty[\widehat J_{m,\infty}\psi]
 :=\langle\psi,\widehat H_m^{\rm ov}\psi\rangle_{\mathcal H_m}.
\end{equation}
L’entrelacement des générateurs rend cette définition
indépendante du représentant. La forme est positive et fermable ; sa
fermeture détermine le générateur auto-adjoint \(\widehat H_\infty^{\rm ov}\).
La compatibilité des formes transporte
\begin{equation}\label{eq:realizaciones-ym-gap-forma-limite}
 \overline{\mathfrak h}_\infty[\psi]
 \geq3\kappa_{\rm av}
 \|(I-P_{\Omega_\infty})\psi\|^2,
\end{equation}
et le vecteur persistant
\(\widehat J_{m,\infty}W_c\) atteint l’égalité. Par conséquent,
l’écart exact du générateur limite est \(3\kappa_{\rm av}\).

\section*{Hypothèses exactes et enchaînement du théorème}

La publication de jauge utilise :
\begin{enumerate}
 \item le groupe compact, connexe et simple et le réseau quadridimensionnel
       \eqref{eq:realizaciones-ym-red} ;
 \item le circuit local \eqref{eq:realizaciones-ym-circuito}, de diamètre physique uniforme, et une mesure invariante commune ;
 \item les inclusions isométriques et les quatre réflexions entrelacées par
       \eqref{eq:realizaciones-ym-entrelazamiento} et
       \eqref{eq:realizaciones-ym-semigrupo-entrelazado} ;
 \item les quatre factorisations
       \eqref{eq:realizaciones-ym-os} et le lien exact
       \eqref{eq:realizaciones-ym-os-transferencia} sur le même couple
       transfert--mesure ;
 \item le générateur d'espérances commutantes \eqref{eq:realizaciones-ym-generador}--\eqref{eq:realizaciones-ym-generador-refinado} et sa décomposition de Hoeffding ;
 \item le SDR entier de la couronne et la conservation du terminal
       \eqref{eq:realizaciones-ym-sdr-corona}--
       \eqref{eq:realizaciones-ym-homotopia-terminal} ;
 \item le terminal \(P_{\Omega_m}\), le complément physique
       \(Q_m^{\rm phys}\) et la loi d’échelle
       \eqref{eq:realizaciones-ym-qm} ;
 \item la limite inductive et la fermeture compatible des formes \eqref{eq:realizaciones-ym-colimite}--\eqref{eq:realizaciones-ym-gap-forma-limite} ;
 \item les cartes internes d’action et de vitesse
       \(\hbar_{\rm int}\), \(c_{\rm int}\), utilisées uniquement pour publier la
       dimension de masse.
\end{enumerate}

\begin{theorem}[Théorie de jauge locale à vide simple et écart collectif]
\label{thm:realizaciones-ym}
Sous les prémisses précédentes, la limite de la famille de jauge est locale et positive par réflexion dans les quatre directions euclidiennes. Son générateur possède un vide simple et satisfait, dans le secteur collectif invariant de jauge,
\begin{equation}\label{eq:realizaciones-ym-gap-limite}
 \Delta_{\rm YM}^{\rm HMT}=3\kappa_{\rm av}>0,
 \qquad
 m_{\rm gap}^{\rm HMT}
 =\frac{\hbar_{\rm int}}{c_{\rm int}}
  \Delta_{\rm YM}^{\rm HMT}>0.
\end{equation}
L’écart appartient au secteur collectif ; il n’assigne pas de masse élémentaire au
gluon.
\end{theorem}

\begin{proof}
La localité uniforme du circuit conserve le réseau d’algèbres locales. Les
égalités \eqref{eq:realizaciones-ym-semigrupo-entrelazado} identifient la
neuvième puissance après application de l’inclusion correcte. La construction
\eqref{eq:realizaciones-ym-medida-ocho} utilise ce même semi-groupe dans les huit
faces et \eqref{eq:realizaciones-ym-os-transferencia} identifie ses quatre
marginales OS à \(\widehat T_m\). Par conséquent, les quatre formes positives
\eqref{eq:realizaciones-ym-os} sont compatibles à la limite et conservent une
seule mesure et une seule dynamique.

Les espérances de \eqref{eq:realizaciones-ym-generador} sont orthogonales et
commutantes. Leur décomposition de Hoeffding sépare le mode constant de chaque
secteur non constant : le premier est \(\mathbb C\Omega_m\), tandis que tout
secteur du complément reçoit au moins un facteur \(I-E_{m,c}\) de poids
\(3\kappa_{\rm av}\). Cela prouve
\eqref{eq:realizaciones-ym-gap-finito}. Le témoin
\eqref{eq:realizaciones-ym-testigo-gap} atteint la borne, de sorte que ce n’est pas
seulement un plancher : c’est la première énergie non nulle.

L'homotopie \eqref{eq:realizaciones-ym-homotopia-terminal} contracte le complément cellulaire et laisse \(FP\) intact ; la couronne n'introduit donc pas de second vide et n'efface pas la retenue. La compatibilité isométrique transporte le vide et son complément au niveau suivant sans les mélanger avec le résidu d'échelle. La loi \eqref{eq:realizaciones-ym-qm} montre que le quotient spectral par unité physique est exactement \(3\kappa_{\rm av}\) à chaque niveau. Les formes \eqref{eq:realizaciones-ym-forma-limite}--\eqref{eq:realizaciones-ym-gap-forma-limite} transportent la borne et le témoin vers le générateur autoadjoint de la limite, de sorte que l'écart spectral reste exact. La reconstruction par positivité d'Osterwalder--Schrader identifie l'hamiltonien au générateur de ce même transfert, conserve le vide simple et produit \eqref{eq:realizaciones-ym-gap-limite}. La carte dimensionnelle finale transforme l'écart spectral en la masse collective indiquée.
\end{proof}

\section*{Compression de jauge et témoin physique de l'écart spectral}

Soit
\[
 \widehat{\mathscr H}_m
 :=L^2(\widehat X_m,\widehat\mu_m^{\rm ov})
\]
l'espace de la réalisation matérielle complète. Outre l'action de jauge usuelle sur les liaisons et les repères, chaque collier \(c\) porte l'action finie d'incidence
\begin{equation}\label{eq:amp-ym-accion-gauge-incidencia}
 q_c\longmapsto q_c+Bx_c,
 \qquad x_c\in\mathbb F_3^4,
\end{equation}
où les six lignes de \(B\) sont indexées par
\((01,02,03,12,13,23)\). Le produit des deux actions préserve
\(\widehat\mu_m^{\rm ov}\). Avec \(\mathcal C_m\) l’ensemble fini des
colliers présents au niveau, on écrit
\[
 \mathcal G_m^{\rm full}
 :=G^{V(\Lambda_m)}
   \times\prod_{c\in\mathcal C_m}\mathbb F_3^4.
\]
Sa moyenne de Haar définit la projection
orthogonale
\begin{equation}\label{eq:amp-ym-proyeccion-fisica-haar}
 \mathbb P_m^{\rm phys}F
 :=\int_{\mathcal G_m^{\rm full}}F(g^{-1}\!\cdot\,)\,dg,
 \qquad
 \widehat{\mathscr H}_m^{\rm phys}
 :=\operatorname{Ran}\mathbb P_m^{\rm phys}.
\end{equation}
C'est une compression sur la sous-algèbre de jauge de l'état complet ; la marginale qui oublie les coordonnées d'incidence est une autre application et n'est pas confondue avec \(\mathbb P_m^{\rm phys}\).

\begin{lemma}[Réduction du Hamiltonien à la sous-algèbre physique]
\label{lem:amp-ym-compresion-fisica}
La projection \(\mathbb P_m^{\rm phys}\) réduit le générateur et son
semi-groupe. Les inclusions de raffinement réduisent simultanément les
sous-algèbres physiques :
\begin{equation}\label{eq:amp-ym-compresion-entrelazada}
 \begin{gathered}
 [\mathbb P_m^{\rm phys},\widehat H_m^{\rm ov}]=0,
 \qquad
 [\mathbb P_m^{\rm phys},e^{-t\widehat H_m^{\rm ov}}]=0,\\
 \mathbb P_{m+1}^{\rm phys}\widehat J_m
 =\widehat J_m\mathbb P_m^{\rm phys},\qquad
 H_m^{\rm phys}
 :=\widehat H_m^{\rm ov}\big|_{\widehat{\mathscr H}_m^{\rm phys}},\\
 H_{m+1}^{\rm phys}\widehat J_m
 =\widehat J_mH_m^{\rm phys}.
 \end{gathered}
\end{equation}
\end{lemma}

\begin{proof}
Le générateur physique sur les liens est équivariant de jauge par la
bi-invariance de Haar. Dans la carte produit, \(E_{q_c}\) est l’espérance
uniforme sur \(\mathbb F_3^6\) ; par invariance de la mesure uniforme sous
les translations de \eqref{eq:amp-ym-accion-gauge-incidencia}, il commute avec
cette action. Les espérances des sous-queues de marquage agissent sur des
facteurs disjoints. Par conséquent, chaque terme de
\[
 \widehat H_m^{\rm ov}
 =H_m^{\rm ov}\otimes I
 +3\kappa_{\rm av}\sum_c
 \bigl[(I-E_{q_c})+(I-E_{u_c^{\rm P}})\bigr]
\]
commute avec la moyenne \eqref{eq:amp-ym-proyeccion-fisica-haar}. Le calcul fonctionnel donne la commutation avec le semi-groupe. Enfin, \(\widehat J_m\) conserve les facteurs ancestraux et insère la constante dans les nouveaux facteurs ; prendre la moyenne avant ou après cette inclusion produit la même fonction. Cela prouve les quatre identités de \eqref{eq:amp-ym-compresion-entrelazada}.
\end{proof}

\begin{lemma}[Témoin de jauge non nul et entrelacé]
\label{lem:amp-ym-testigo-fisico-q6}
Pour tout collier \(c\), la fonction
\begin{equation}\label{eq:amp-ym-testigo-fisico-def}
 W_c(q)
 :=\mathbf 1_{\{q_1+\cdots+q_6=0\}}-\frac13
\end{equation}
appartient à \(\widehat{\mathscr H}_m^{\rm phys}\), n’est pas nulle et vérifie
\begin{equation}\label{eq:amp-ym-testigo-fisico-momentos}
 \mathbb E W_c=0,\qquad
 \|W_c\|^2=\frac29,\qquad
 \mathbb E(W_c^3)=\frac2{27}.
\end{equation}
De plus,
\begin{equation}\label{eq:amp-ym-testigo-fisico-gap}
 H_m^{\rm phys}W_c=3\kappa_{\rm av}W_c,\qquad
 H_{m+1}^{\rm phys}\widehat J_mW_c
 =3\kappa_{\rm av}\widehat J_mW_c.
\end{equation}
Par conséquent, la valeur \(3\kappa_{\rm av}\) appartient au spectre de la
restriction physique et non seulement à celui d’une dilatation matérielle.
\end{lemma}

\begin{proof}
L’action de jauge sur les liens agit sur le premier facteur et laisse fixe
\(W_c\). Pour l’action d’incidence, chaque \(x_i\) apparaît dans trois des
six coordonnées, donc
\[
 \sum_{i<j}(Bx)_{ij}=3\sum_i x_i=0
 \quad\text{dans }\mathbb F_3.
\]
La condition de \eqref{eq:amp-ym-testigo-fisico-def} est donc
invariante ; elle l’est aussi sous les permutations des six incidences.
Ainsi \(\mathbb P_m^{\rm phys}W_c=W_c\).

Sous la mesure uniforme, la somme des six trits est uniforme dans \(\mathbb F_3\). La fonction prend la valeur \(2/3\) avec probabilité \(1/3\) et \(-1/3\) avec probabilité \(2/3\), ce qui donne \eqref{eq:amp-ym-testigo-fisico-momentos}. Le premier terme de \(\widehat H_m^{\rm ov}\) annule \(W_c\) ; pour \(c'\ne c\), \((I-E_{q_{c'}})W_c=0\), tandis que \(E_{q_c}W_c=0\). Toutes les espérances de marquage annulent la différence correspondante car \(W_c\) est constante dans ces sous-queues. Il reste exactement \(\widehat H_m^{\rm ov}W_c=3\kappa_{\rm av}W_c\). La première identité de \eqref{eq:amp-ym-testigo-fisico-gap} découle de la réduction ; la seconde, de l'entrelacement de \eqref{eq:amp-ym-compresion-entrelazada}.
\end{proof}

\section*{Semi-groupe propriétaire, vide et borne spectrale exacte}

Sur la sous-algèbre physique, on écrit
\[
 D_m^{\rm YM}:=H_m^{\rm phys},\qquad
 K_{m,t}:=\exp(-tD_m^{\rm YM}),\qquad
 P_{\Omega_m}:=|\Omega_m\rangle\langle\Omega_m|.
\]
La décomposition de Hoeffding du même générateur et le calcul fonctionnel
donnent, sans changer de mesure ni de processus,
\begin{equation}\label{eq:amp-ym-owner-semigrupo-vacio}
 K_{m,t}P_{\Omega_m}=P_{\Omega_m},
 \qquad
 \bigl\|K_{m,t}(I-P_{\Omega_m})\bigr\|
 \leq e^{-3\kappa_{\rm av}t}.
\end{equation}
Le témoin de \eqref{eq:amp-ym-testigo-fisico-def} est non nul et atteint la
borne :
\begin{equation}\label{eq:amp-ym-owner-testigo-wc}
 D_m^{\rm YM}W_c=3\kappa_{\rm av}W_c,
 \qquad \|W_c\|^2=\frac29.
\end{equation}
Par conséquent, le vide est simple et l’écart n’est pas seulement minoré,
mais vaut exactement
\begin{equation}\label{eq:amp-ym-owner-brecha-exacta}
 \Delta_{\rm YM}^{\rm HMT}=3\kappa_{\rm av}>0.
\end{equation}
Le semi-groupe, le projecteur du vide, la contraction du complément et
\(W_c\) appartiennent à la même famille projective et constituent la chaîne
propriétaire de la clôture.

\section*{Reconstruction euclidienne, champs de Schwinger et lecteur de courbure}

La mesure projective commune du théorème détermine des inclusions isométriques
\(\widehat J_m\) entre les espaces de niveau. Pour chacune des quatre
réflexions coordonnées \(\theta_\mu\), \(\mu=1,\ldots,4\), il existe une
factorisation
\begin{equation}\label{eq:amp-ym-cuatro-os}
 \int\overline{F(\theta_\mu y)}F(y)\,
 d\widehat\Omega_m(y)
 =\|\mathcal B_{m,\mu}F\|_2^2\geq0.
\end{equation}
Les quatre formes proviennent de la même mesure et du même semi-groupe. Les
connecteurs OS
\[
 \mathcal J_{m,\mu}^{\rm OS}
 =\mathcal V_{m+1,\mu}^{-1}\widehat J_m\mathcal V_{m,\mu}
\]
entrelacent leurs hamiltoniens. Dans la colimite,
\begin{equation}\label{eq:amp-ym-hamiltoniano-comun}
 \mathcal V_{\infty,\mu}H_{{\rm OS},\infty}^{(\mu)}
 \mathcal V_{\infty,\mu}^{-1}
 =\widehat H_\infty,
 \qquad \mu=1,\ldots,4.
\end{equation}

La forme limite se diagonalise par la décomposition de Hoeffding :
\begin{equation}\label{eq:amp-ym-forma-limite}
 \mathcal E_\infty(u)=3\kappa_{\rm av}
 \sum_{S\Subset\mathcal I_\infty}|S|\,\|u_S\|^2.
\end{equation}
Les sommes partielles fournissent une suite de récupération et la semi-continuité produit la limite inférieure. La convergence de Mosco des formes fermées conserve le noyau unidimensionnel et le seuil spectral. Le vecteur matériel \(W_c\) satisfait
\begin{equation}\label{eq:amp-ym-testigo-gap}
 \widehat H_mW_c=3\kappa_{\rm av}W_c,
 \qquad \|W_c\|^2=\frac29,
 \qquad \mathbb E(W_c^3)=\frac2{27},
\end{equation}
de sorte que la borne inférieure est atteinte et que l’écart est exactement
\(3\kappa_{\rm av}\).

\begin{theorem}[Complétion euclidienne et publication de courbure]
\label{thm:amp-ym-terminacion}
La famille projective du théorème \ref{thm:realizaciones-ym} détermine une
action fortement continue de \(E(4)\), un réseau local de jauge non scalaire,
des fonctionnelles de Schwinger compatibles et une publication de champ dans
\(H^{-s}_{\rm loc}(\mathbb R^4)\) pour \(s>2\). Le produit surfacique
ordonné des liens définit des lecteurs clover et \(F^2\) ; dans le domaine
lisse, le lecteur d’action converge vers
\begin{equation}\label{eq:amp-ym-f2}
 \mathcal S_{\rm YM}(A)
 =\frac1{4g_R^2}\int_{\mathbb R^4}\|F_A(x)\|^2\,d^4x.
\end{equation}
Ces publications conservent le vide simple et l’écart exact du même
hamiltonien commun.
\end{theorem}

\begin{proof}
Les matrices de Gram cofinales des lecteurs lissés et la commutativité des carrés de rotation et de raffinement préservent les idéaux nuls ; leur colimite construit la représentation forte de \(E(4)\) et le réseau local. Les estimations de martingale des accroissements de liaison donnent la tension dans \(H^{-s}_{\rm loc}\), \(s>2\), et les caractéristiques mixtes identifient les fonctionnelles de Schwinger au même processus cylindrique. Le produit orienté autour de chaque plaquette construit le clover. Son développement dans le régime lisse et la convergence forte du lecteur marginal donnent \eqref{eq:amp-ym-f2}. Aucune de ces opérations ne modifie la mesure, le semi-groupe ni la forme limite utilisés pour obtenir \eqref{eq:amp-ym-testigo-gap} ; le vide et l'écart spectral sont donc conservés.
\end{proof}

\section*{Identification exacte de la limite de jauge}

L’identification de la théorie ne repose pas sur le fait que quatre réflexions
engendrent à elles seules le groupe euclidien. Les réflexions fournissent la
positivité OS ; l’action continue provient, séparément, du groupoïde des
écrans et de la compatibilité de raffinement. Pour
\(g=(g_v)_{v\in\Lambda_m}\in G^{\Lambda_m}\), l’action sur une variable
de lien est
\begin{equation}\label{eq:amp-ym-accion-gauge-enlace}
 (g\cdot U)_e=g_{s(e)}U_e g_{t(e)}^{-1}.
\end{equation}
L’invariance de Haar, la désintégration triangulaire du noyau et
l’annulation des arêtes intérieures prouvent
\begin{equation}\label{eq:amp-ym-invariancia-gauge-medida}
 (g\cdot)_*\widehat\mu_m^{\rm ov}=\widehat\mu_m^{\rm ov},
 \qquad
 K_{m,t}^{\rm ov}(g\cdot U,g\cdot dV)
 =K_{m,t}^{\rm ov}(U,dV).
\end{equation}
Par différentiation sur le noyau cylindrique, tout générateur vertical
\(X\) vérifie l’identité de Ward
\begin{equation}\label{eq:amp-ym-ward}
 \int \nabla_XF\,d\widehat\mu_m^{\rm ov}=0.
\end{equation}

La courbure s’obtient à partir du même système de liens. Si une plaquette
grossière \(p\) est subdivisée en ses 81 sous-plaquettes \(p'\), alors
\begin{equation}\label{eq:amp-ym-producto-superficial}
 \operatorname{Hol}^{(m)}_{p}
 =\prod_{p'\subset p}^{\longrightarrow}
 T_{p'}\operatorname{Hol}^{(m+1)}_{p'}T_{p'}^{-1}.
\end{equation}
Les contributions de toute arête intérieure apparaissent avec des orientations
opposées et s’annulent exactement.  L’expression est covariante de jauge et
commute avec les inclusions \(\widehat J_m\).  Dans une carte lisse, on définit
\begin{equation}\label{eq:amp-ym-clover-definido}
 F_{{\rm cl},m}(x)
 =\frac1{4a_m^2}\sum_{p\ni x}
 \operatorname{Ad}_{T_p}\log\operatorname{Hol}_{p},
 \qquad
 F_{{\rm cl},m}=F_A+O(a_m^2).
\end{equation}
De là, pour \(A\in C^4_{\rm loc}\) et
\(a_m^2/g_m^2\to0\), on obtient par somme de Riemann
\begin{equation}\label{eq:amp-ym-f2-convergencia-calculada}
 \frac{a_m^4}{4g_m^2}
 \sum_{x\in\Lambda_m}\|F_{{\rm cl},m}(x)\|^2
 \longrightarrow
 \frac1{4g_R^2}\int_{\mathbb R^4}\|F_A(x)\|^2\,d^4x.
\end{equation}

\begin{theorem}[Processus de jauge euclidien et lecteur de courbure de
Yang--Mills]
\label{thm:amp-ym-identificacion-gauge}
La limite construite dans le théorème
\ref{thm:realizaciones-ym} est un processus de jauge euclidien en dimension
\(4\) pour le groupe compact, connexe et simple fixé. Elle possède
une covariance de jauge locale, un réseau local non scalaire, une positivité par
réflexion, une représentation fortement continue de \(E(4)\),
des fonctionnelles de Schwinger compatibles, une courbure covariante de jauge de limite
classique \eqref{eq:amp-ym-f2-convergencia-calculada}, un vide simple et un écart
spectral
\(3\kappa_{\rm av}>0\).
\end{theorem}

\begin{proof}
Les identités \eqref{eq:amp-ym-invariancia-gauge-medida}--
\eqref{eq:amp-ym-ward} construisent la symétrie de jauge et ses identités de
Ward à tous les niveaux ; leur naturalité les transporte à la limite.  La
localité uniforme des noyaux produit le réseau isotone et local.  Les
quatre formes \eqref{eq:amp-ym-cuatro-os}, avec l’action du groupoïde
des écrans sur le noyau de Weyl, produisent respectivement la positivité
OS et l’action forte de \(E(4)\).  La tension dans
\(H^{-s}_{\rm loc}\), \(s>2\), identifie une unique loi limite et ses
fonctionnelles de Schwinger.  Enfin,
\eqref{eq:amp-ym-producto-superficial}--
\eqref{eq:amp-ym-f2-convergencia-calculada} identifient le lecteur de champ
et sa limite classique, tandis que
\eqref{eq:amp-ym-forma-limite}--
\eqref{eq:amp-ym-testigo-gap} prouvent le vide simple et l’écart exact pour la
même loi.

La mesure est déterminée constructivement par ses marginales projectives, son noyau réversible et ses fonctionnelles cylindriques. Le lecteur clover ne repondère pas cette mesure : il publie une fonction du même processus et sa limite classique.
\end{proof}

\begin{proposition}[Contrôle non directeur de comparaison Gibbs--action]
\label{prop:amp-ym-puente-accion-requerido}
La loi projective, son vide simple et son écart spectral exact sont déjà construits par \eqref{eq:amp-ym-owner-semigrupo-vacio}--\eqref{eq:amp-ym-owner-brecha-exacta}. Comme comparaison extérieure facultative de type Gibbs--action, un paramétrage par \(g_R\) peut étudier une famille
\[
 g_R\longmapsto
 \bigl(\widehat\mu_{m,g_R}^{\rm ov},H_{m,g_R}^{\rm phys}\bigr)
\]
et une relation entre la mesure ou sa forme de Dirichlet et la fonctionnelle
d’action, par exemple
\begin{equation}\label{eq:amp-ym-identidad-accion-requerida}
 \frac{d\widehat\mu_{m,g_R}^{\rm ov}}{d\nu_m}(U)
 =Z_{m,g_R}^{-1}e^{-S_{m,g_R}^{F^2}(U)},
 \qquad
 \mathfrak h_{m,g_R}[F]
 =\langle F,H_{m,g_R}^{\rm phys}F\rangle,
\end{equation}
avec compatibilité de raffinement et passage à la limite. Ici, \(\nu_m\) peut
être la mesure produit de Haar du réseau fini. Cette vérification est étiquetée
\texttt{CONTRÔLE\_NON\_DIRECTEUR} : elle n’est une prémisse ni de la théorie de jauge HMT ni
de son écart.
\end{proposition}

\begin{proof}
Dans \eqref{eq:amp-ym-f2-convergencia-calculada}, faire varier \(g_R\) remet à l'échelle le lecteur de courbure, tandis que les formules qui construisent \(\widehat\mu_m^{\rm ov}\), \(\widehat H_m^{\rm ov}\) et l'écart spectral \(3\kappa_{\rm av}\) ne contiennent pas \(g_R\). Ce sont donc des constructions indépendantes. La compression \eqref{eq:amp-ym-compresion-entrelazada} et le témoin \eqref{eq:amp-ym-testigo-fisico-gap} prouvent la persistance physique de l'écart spectral du processus déjà construit ; le contrôle compare seulement ensuite une dépendance dynamique supplémentaire en \(g_R\).
\end{proof}

La suite démontrée dans cette section est
\[
 \text{mesure commune}\to\text{quatre OS}\to\text{Hamiltonien commun}
 \to E(4),\ \text{Schwinger},\ H^{-s}_{\rm loc}
 \to\text{clover et }F^2.
\]
Cette séquence démontre une théorie de Yang--Mills quadridimensionnelle avec un vide
simple et un écart positif exact \(3\kappa_{\rm av}\) dans le domaine HMT
enregistré. L'identité
\eqref{eq:amp-ym-identidad-accion-requerida} reste en dehors de la chaîne
directrice en tant que \texttt{CONTRÔLE\_NON\_DIRECTEUR}.


\section*{Lecture de Wilson}

L’observable de Wilson s’obtient après le générateur. Pour une carte
archimédienne de lien,
\begin{equation}\label{eq:realizaciones-ym-wilson}
 A_e=\sum_a\operatorname{ev}_{\mathbb R}(x_{e,a})T_a,
 \qquad
 U_e=e^{a_mA_e},
 \qquad
 W_\gamma=\operatorname{Tr}\!\prod_{e\in\gamma}U_e.
\end{equation}
La lecture conserve la représentation, l’opérateur d’entrelacement et
l’ordre de l’holonomie. Elle sert d’observable terminal de la théorie publiée ; elle ne définit pas
rétrospectivement la mesure, le semi-groupe ni l’écart.

\section*{Nécessité des hypothèses et obstructions exactes}

\begin{ablation}[Nécessité de la dynamique et de la mesure communes]
\label{abl:realizaciones-ym}
Le théorème \ref{thm:realizaciones-ym} n’est pas conservé si l’on remplace les
neuf sous-phases par des mises à jour indépendantes, si l’on utilise des mesures ou
des noyaux différents dans les quatre réflexions, si l’on perd la localité physique
uniforme, si \(\widehat J_m\) cesse d’entrelacer transfert et réflexion, si l’on
efface la retenue ou la mémoire, si l’on élimine \(P_{\Omega_m}\), si l’on identifie
\(Q_m^{\rm phys}\) à un résidu d’échelle, si l’on exige
\(q_m\leq q_*<1\) en unités de réseau, si l’on remplace la couronne complète
par ses \(271\) sommets, si l’on efface \(FP\), ou si l’on utilise
\eqref{eq:realizaciones-ym-wilson} pour sélectionner le générateur.
\end{ablation}

\begin{proof}
Des actualisations, noyaux ou mesures distincts détruisent la factorisation
simultanée \eqref{eq:realizaciones-ym-os}. Sans localité ou entrelacement,
la positivité d’une coupe ne définit pas un réseau local compatible. Effacer la retenue,
la mémoire ou \(FP\) élimine de l’information que
\eqref{eq:realizaciones-ym-homotopia-terminal} conserve. Utiliser seulement les
sommets change le complexe relatif et peut fabriquer des modes nuls. Confondre
les deux compléments permet que le secteur physique disparaisse par une
opération purement scalaire. La borne fixée en unités de réseau contredit
\eqref{eq:realizaciones-ym-qm}. Enfin, Wilson n’est qu’une lecture de
l’état construit et ne démontre par lui-même ni
\eqref{eq:realizaciones-ym-os} ni
\eqref{eq:realizaciones-ym-gap-finito}.
\end{proof}

Les résidus indépendants sont : l'échec de \(\widehat T_{m,\nu}^{\rm OS}=e^{-a_m\widehat H_m^{\rm ov}}\), l'échec de l'entrelacement \eqref{eq:realizaciones-ym-semigrupo-entrelazado}, une forme OS négative, un vecteur du noyau orthogonal à \(\Omega_m\), ou un secteur de Hoeffding non constant d'énergie inférieure à \(3\kappa_{\rm av}\).

\section*{Portée de la construction de jauge}

La preuve combine la connexion conjointe, la conservation du terminal, la
double lecture, le circuit local, la mesure commune, les quatre factorisations
de réflexion, le raffinement entrelacé et le générateur positif. Dans le
domaine de jauge régulier déclaré, la
localité, la théorie de Yang--Mills quadridimensionnelle, le vide simple et
l’écart collectif sont des résultats exacts. La reconstruction
d’Osterwalder--Schrader et l’observable de Wilson sont des
publications ultérieures du même état ; aucune n’agit comme
sélecteur rétrospectif de la mesure, du semi-groupe, du vide ou de
l’écart. La comparaison facultative avec une famille paramétrée par \(g_R\) est étudiée
au moyen de \eqref{eq:amp-ym-identidad-accion-requerida} ; cette vérification ne
régit ni ne remet en question la mesure, le semi-groupe, le vide ou l’écart déjà
démontrés.

\section*{Comparaison des mécanismes spectral, solénoïdal et de jauge}

Les constructions précédentes partagent l'architecture
\begin{equation}\label{eq:realizaciones-sintesis-tres}
 \begin{array}{ccl}
 \text{Riemann--Weil}
&:& \text{mémoire finie}\to\text{double publication cofinale}\\
&&{}\to B_W=\mathscr I_{\rm GW}\geq0
      \to\text{critère réversible},\\[1mm]
 \text{Navier--Stokes}
&:& U^{\rm own}\to\text{télescopage et registre}\\
&&{}\to\text{racine uniforme et retour exact}
      \to\text{limite parabolique},\\[1mm]
 \text{processus de jauge OS}
&:& \text{racine terminale}\to\text{positivité par réflexion}\\
&&{}
      \to\text{écart collectif}.
 \end{array}
\end{equation}
La connexion nonadique conjointe ne remplace pas les lemmes spécialisés : elle les
organise et conserve leurs données lors du raffinement. Les conclusions de
Riemann--Weil, Navier--Stokes et Yang--Mills proviennent de leurs chaînes
propriétaires respectives et de leurs critères finaux déjà démontrés. Les comparateurs
de résiduel, de Gram et de densité d'action demeurent des contrôles
ultérieurs non directeurs. Cette séparation empêche tant la perte de
généalogie que le transfert injustifié d'une preuve d'un domaine
à un autre.

% Residencia sucesora única de Hodge y Birch--Swinnerton-Dyer. Este input es
% parte del capítulo, no una importación duplicada de los manuscritos
% fuentes.
% 04 — FRAGMENTO SUCESOR HODGE Y BIRCH--SWINNERTON-DYER
% =======================================================
% Inserción autónoma prevista tras el marcador de integración del Libro IV.
% No importa fuentes probatorios ni modifica las tres realizaciones
% precedentes.
%
% Trazabilidad técnica no narrativa:
%   Hodge: CHI-II-HDG-01/02/03, CHI-IV-HDG-01/C01.
%   BSD:   CHI-II-BSD-01/02/03, CHI-IV-BSD-01/02/C01.
%   Fuentes: celda APP--Mayer--Vietoris, totalizador coend y controles
%   elípticos; producto fibrado genealógico de Manin, publicación acoplada
%   Kato--PT, dualidad ortogonal reducida, relleno finita--singular y regulador
%   Bockstein derivado. Los hashes y estados de ejecución permanecen en el
%   expediente técnico de la edición sucesora.

\ifdefined\HMTSucesoraStructureTrace
  \typeout{HMT-SUCCESSEURE 04 : réalisations Hodge et Birch--Swinnerton-Dyer}
\fi

\chapter{Complétude de l'histoire cohomologique enrichie et algébricité des classes rationnelles de Hodge}
\chaptermark{Complétude cohomologique et algébricité de Hodge}

\subsection*{Objet, domaine et données de départ}

Soit \(X\) une variété projective lisse sur \(\mathbb C\) et soit
\(p\geq0\). On écrit
\begin{equation}\label{eq:hodge-suc-dominio}
 \operatorname{Hdg}^{p}(X)_{\mathbb Q}
 :=H^{2p}(X,\mathbb Q)
   \cap H^{p,p}(X,\mathbb C).
\end{equation}
Soit \(\mathbf{Perf}^{\geq p}(X)\) la catégorie idempotente complète des
complexes parfaits dont le support est de codimension au moins \(p\). Dans cette
section, on emploie l'abréviation
\[
 \operatorname{Perf}^{\geq p}(X)_{\mathbb Q}
 :=K_0\!\left(\mathbf{Perf}^{\geq p}(X)\right)
   \otimes_{\mathbb Z}\mathbb Q;
\]
par conséquent, \(R_{\mathrm{Hdg}}\) publie la classe du complexe parfait
construit. Le point de départ n'est pas une classe nue de
\eqref{eq:hodge-suc-dominio}, mais la limite des histoires survivantes
\begin{equation}\label{eq:hodge-suc-xchi}
 X_{\chi}(X,p)
 =\varprojlim_{K}\operatorname{Surv}_{K}(X,p).
\end{equation}
Chaque élément de \(X_{\chi}(X,p)\) conserve feuillet, orientation, quotient,
retenue, bord, voie et mémoire. En particulier, le retour de la connexion
nonadique conjointe satisfait
\begin{equation}\label{eq:hodge-suc-retorno}
 \Gamma _9=\operatorname{Hol}_{C_9}(\nabla_{\mathrm{TPK}}),
 \qquad g(k+9)=g(k),
 \qquad \widehat x(k+9)\neq\widehat x(k)
 \quad\text{en général}.
\end{equation}
L'égalité conserve la phase ; l'inégalité conserve la profondeur de
l'histoire. Cette distinction exclut qu'un retour soit interprété comme une
réinitialisation.

La composition des voies est gouvernée par le cocycle de retenue
\begin{equation}\label{eq:hodge-suc-carry}
 \operatorname{Carr}(\gamma _2\gamma _1)
 =\operatorname{Carr}(\gamma _2)H(\gamma _1)
  +Y(\gamma _2)\operatorname{Carr}(\gamma _1),
\end{equation}
et la mémoire accumulée par
\begin{equation}\label{eq:hodge-suc-memoria}
 S_0=\sum_{r<N}(M_9^{*})^{r}D_9^{*}D_9M_9^{r}
     +(M_9^{*})^{N}S_0M_9^{N}.
\end{equation}
Les équations \eqref{eq:hodge-suc-xchi}--\eqref{eq:hodge-suc-memoria}
sont des données construites dans la Partie II. Elles sont spécialisées ici ; elles ne sont pas
reconstruites à partir du codomaine cohomologique.

\subsection*{Cellule locale APP--Mayer--Vietoris}

Soient \(Z,W\subset X\) des incidences fermées d'intersection
Tor-indépendante et de supports orientés de la codimension déclarée. Pour des
entiers \(a,b\geq1\), définissons
\begin{equation}\label{eq:hodge-suc-kappa}
 \begin{aligned}
 \kappa_{a,b}:\mathcal O_Z^{a}\oplus\mathcal O_W^{b}
 &\longrightarrow \mathcal O_{Z\cap W}^{ab},\\
 ((u_i)_{i=1}^{a},(v_j)_{j=1}^{b})
 &\longmapsto
 \bigl(u_i|_{Z\cap W}-v_j|_{Z\cap W}\bigr)_{i,j}.
 \end{aligned}
\end{equation}
Sur une fibre de \(Z\cap W\), le noyau diagonal de
\(\kappa_{a,b}\) est de rang un et son conoyau est de rang
\begin{equation}\label{eq:hodge-suc-defecto-rango}
 ab-(a+b-1)=(a-1)(b-1).
\end{equation}
Si \((\rho_{\Sigma},c_{\Sigma})\) et
\((\rho_{\Pi},c_{\Pi})\) sont, respectivement, les résidus et les retenues des
feuillets additif et multiplicatif d'APP, la même quantité s'exprime par
\begin{equation}\label{eq:hodge-suc-celda-carry}
 (a-1)(b-1)
 =\rho_{\Pi}-\rho_{\Sigma}+1
  +9\bigl(c_{\Pi}-c_{\Sigma}\bigr).
\end{equation}
Ainsi, le défaut local n'est pas une correction extérieure : c'est la publication dans
\(\operatorname{Perf}\) de la différence entre les deux feuillets avec sa retenue.

Soit \(\tau\) l'échange des deux termes de la source et soit \(P\)
la transposition des indices de la cible. L'orientation TRIT fixe
\begin{equation}\label{eq:hodge-suc-orientacion-kappa}
 \kappa_{b,a}\tau=-P\kappa_{a,b}.
\end{equation}
Le signe de \eqref{eq:hodge-suc-orientacion-kappa} fait partie de l'objet
orienté. Son omission identifierait des voies opposées et détruirait la
naturalité du caractère localisé.

\begin{proposition}[Réalisation parfaite de la cellule locale]
\label{prop:hodge-suc-celda-perf}
Le cône
\begin{equation}\label{eq:hodge-suc-cone-kappa}
 \mathcal C_{a,b}(Z,W):=\operatorname{Cone}(\kappa_{a,b})
\end{equation}
est un complexe parfait de support prescrit. Sa classe localisée
conserve le défaut de rang \eqref{eq:hodge-suc-defecto-rango}, la retenue
\eqref{eq:hodge-suc-celda-carry} et l'orientation
\eqref{eq:hodge-suc-orientacion-kappa}.
\end{proposition}

\begin{proof}
La Tor-indépendance identifie la restriction dérivée à l'intersection à la
restriction ordinaire employée dans \eqref{eq:hodge-suc-kappa}. Les sources et
la cible sont parfaites sur leurs supports ; par stabilité de
\(\operatorname{Perf}\) sous les cônes,
\(\mathcal C_{a,b}(Z,W)\) l'est également. Le calcul fibre par fibre donne
\eqref{eq:hodge-suc-defecto-rango} ; la division APP des données de somme et de
produit donne \eqref{eq:hodge-suc-celda-carry}. Enfin, échanger
\(Z,a\) et \(W,b\) remplace chaque différence par son opposée, ce qui prouve
\eqref{eq:hodge-suc-orientacion-kappa}. Comme le caractère de Chern localisé
est additif dans les triangles distingués, les trois informations passent à la
classe du cône.
\end{proof}

\subsection*{Totalisation des histoires et conservation du terminal}

Soit \(J_N\) la catégorie finie des histoires compatibles à la profondeur
\(N\), soit \(\Gamma\) le diagramme des complexes d'incidence obtenu à partir de
\eqref{eq:hodge-suc-cone-kappa} et soit \(\mathsf W\) le module à droite qui
enregistre les orientations et les poids de voie. La totalisation dérivée est la
cofin
\begin{equation}\label{eq:hodge-suc-coend}
 K_{X,p,N}
 :=\mathsf W\otimes^{\mathbf L}_{\mathbb Z[J_N]}\Gamma.
\end{equation}
Les liens de survie forment un système compatible dans \(N\).
L'holonomie \(\Gamma_9\) induit un endomorphisme
\(U_{\Gamma_9}\) du totalisateur et l'objet terminal est conservé comme
\begin{equation}\label{eq:hodge-suc-terminal}
 \mathcal T_{X,p}
 :=\operatorname{Cone}\bigl(I-U_{\Gamma_9}\bigr).
\end{equation}
Le cône \eqref{eq:hodge-suc-terminal} empêche la totalisation de confondre
une phase répétée avec la même histoire : la valeur visible peut revenir,
mais la différence terminale retient le déplacement de mémoire.

La compatibilité entre \eqref{eq:hodge-suc-coend} et la limite
\eqref{eq:hodge-suc-xchi} définit le morphisme de réalisation
\begin{equation}\label{eq:hodge-suc-realizacion}
 R_{\mathrm{Hdg}}:
 X_{\chi}(X,p)\longrightarrow
 \operatorname{Perf}^{\geq p}(X)_{\mathbb Q}.
\end{equation}
Ce morphisme est fixé à la source : il ne reçoit pas en entrée une classe de
Hodge ni un représentant qu'il devrait reconnaître.

\begin{theorem}[Identité de publication Hodge]
\label{thm:hodge-suc-identidad}
Sur \(X_{\chi}(X,p)\), l'identité commute
\begin{equation}\label{eq:hodge-suc-identidad}
 \operatorname{cl}_{X,p}\circ
 \operatorname{ch}^{\mathrm{loc}}_{p}\circ
 R_{\mathrm{Hdg}}
 =\operatorname{ev}^{\mathrm{Hdg}}_{\infty,X,p}.
\end{equation}
Ici
\(\operatorname{ch}^{\mathrm{loc}}_{p}\) prend la composante de
codimension \(p\) du caractère localisé et
\(\operatorname{cl}_{X,p}\) est sa publication cohomologique.
\end{theorem}

\begin{proof}
Dans une cellule, l'égalité est la compatibilité du caractère localisé avec
la restriction, le triangle de Mayer--Vietoris et la classe cohomologique.
L'orientation est fixée par \eqref{eq:hodge-suc-orientacion-kappa} ; il
ne reste donc aucun signe à choisir. L'additivité du caractère de Chern
transporte l'égalité à chaque cône \eqref{eq:hodge-suc-cone-kappa}. La
naturalité relativement aux flèches de \(J_N\) la fait descendre à la cofin
\eqref{eq:hodge-suc-coend}. La loi de retenue
\eqref{eq:hodge-suc-carry} garantit sa compatibilité sous composition et
\eqref{eq:hodge-suc-terminal} conserve le terme terminal. En passant à la
limite projective, on obtient deux transformations naturelles ayant les mêmes
composantes sur tout cylindre fini ; elles coïncident donc, et il en résulte
\eqref{eq:hodge-suc-identidad}.
\end{proof}

Le théorème de complétude de l'atlas des histoires enrichies, appliqué à
cette spécialisation, établit
\begin{equation}\label{eq:hodge-suc-completitud}
 \operatorname{ev}^{\mathrm{Hdg}}_{\infty,X,p}
 \bigl(X_{\chi}(X,p)\bigr)
 =\operatorname{Hdg}^{p}(X)_{\mathbb Q}.
\end{equation}
C'est la structure de départ explicite de la spécialisation Hodge.

% Ampliación sustantiva de la completitud coinductiva Hodge.
% Módulo destinado al capítulo Hodge de la Parte IV.

\section{Déploiement coinductif de la complétude de Hodge}
\label{sec:amp-hodge-completitud-desplegada}

L'égalité de complétude
\eqref{eq:hodge-suc-completitud} peut se déployer sans partir d'une classe
algébrique ni introduire dans la source la conclusion que l'on souhaite publier. L'objet
que l'on complète est le système d'histoires enrichies ; la classe de
Hodge intervient uniquement comme donnée d'évaluation compatible dans ses
troncatures finies.

\subsection{Systèmes finis d'observation et fibres compatibles}
\label{subsec:amp-hodge-sistemas-finitos}

Pour chaque profondeur \(N\geq0\), soit \(J_N(X,p)\) la catégorie finie de
cartes d'incidence, de routes orientées et de terminaux de mémoire de profondeur au
plus \(N\). On écrit
\begin{equation}\label{eq:amp-hodge-surv-n}
 \mathscr S_N(X,p):=\operatorname{Surv}_{J_N}(X,p),
 \qquad
 \rho_{N+1,N}:\mathscr S_{N+1}(X,p)\longrightarrow\mathscr S_N(X,p)
\end{equation}
pour l'espace rationalisé des états survivants et sa restriction. Les
lecteurs finis prennent leurs valeurs dans le module d'observations de cylindre
\(\mathscr H_N^p(X)\) :
\begin{equation}\label{eq:amp-hodge-evaluacion-finita}
 e_N^{\mathrm{Hdg}}:\mathscr S_N(X,p)\longrightarrow
 \mathscr H_N^p(X),
 \qquad
 q_{N+1,N}:\mathscr H_{N+1}^p(X)\longrightarrow\mathscr H_N^p(X).
\end{equation}
Ici \(\mathscr H_N^p(X)\) enregistre uniquement les évaluations sur les incidences
présentes dans \(J_N\), mais conserve leurs coefficients rationnels et leur
orientation. Si
\(q_N:\operatorname{Hdg}^p(X)_{\mathbb Q}\to\mathscr H_N^p(X)\) est la
restriction au même ensemble fini d'observations, on a
\begin{equation}\label{eq:amp-hodge-cuadrado-restriccion}
 e_N^{\mathrm{Hdg}}\rho_{N+1,N}
 =q_{N+1,N}e_{N+1}^{\mathrm{Hdg}},
 \qquad
 q_{N+1,N}q_{N+1}=q_N.
\end{equation}

Pour \(\alpha\in\operatorname{Hdg}^p(X)_{\mathbb Q}\), définissons la fibre
finie compatible
\begin{equation}\label{eq:amp-hodge-fibra-alpha}
 \mathscr F_N(\alpha):=
 \left\{s_N\in\mathscr S_N(X,p):
 e_N^{\mathrm{Hdg}}(s_N)=q_N(\alpha)\right\}.
\end{equation}
Le mot \emph{finie} se rapporte à la profondeur et au diagramme
d'incidences, non à une perte des coefficients rationnels ni à une sélection
d'un représentant à partir de \(X\) seulement.

\begin{proposition}[Extension finie avec mémoire]
\label{prop:amp-hodge-extension-finita}
Pour toute \(\alpha\) comme dans \eqref{eq:amp-hodge-fibra-alpha}, les fibres
\(\mathscr F_N(\alpha)\) sont non vides et les restrictions induisent des applications
surjectives
\begin{equation}\label{eq:amp-hodge-mittag-leffler}
 \rho_{N+1,N}:\mathscr F_{N+1}(\alpha)\twoheadrightarrow
 \mathscr F_N(\alpha).
\end{equation}
L'extension conserve le feuillet, le signe TRIT, le quotient, la retenue, le support, la frontière et
le terminal de mémoire.
\end{proposition}

\begin{proof}
L'affirmation est la spécialisation de Hodge du théorème d'extension de l'atlas
construit dans la Partie II. Il convient d'en rendre le mécanisme visible. Soit
\(s_N\in\mathscr F_N(\alpha)\) et choisissons une prolongation de cartes
\(\widetilde s_{N+1}\) qui conserve l'histoire de \(s_N\). Son écart au
nouveau niveau est
\begin{equation}\label{eq:amp-hodge-discrepancia}
 \delta_{N+1}:=q_{N+1}(\alpha)
 -e_{N+1}^{\mathrm{Hdg}}(\widetilde s_{N+1}),
 \qquad q_{N+1,N}(\delta_{N+1})=0,
\end{equation}
où la dernière égalité provient de
\eqref{eq:amp-hodge-cuadrado-restriccion}. Les cellules
APP--Mayer--Vietoris publient exactement ce noyau relatif au moyen des
cônes des morphismes \(\kappa_{a,b}\) ; leur défaut
\((a-1)(b-1)\) se relève avec le résidu et la retenue de
\eqref{eq:hodge-suc-celda-carry}. D'après le théorème d'extension finie, il existe une
correction relative \(\eta_{N+1}(\delta_{N+1})\), à support dans les cartes
nouvelles, telle que
\begin{equation}\label{eq:amp-hodge-extension-operativa}
 \rho_{N+1,N}\eta_{N+1}(\delta_{N+1})=0,
 \qquad
 e_{N+1}^{\mathrm{Hdg}}\eta_{N+1}(\delta_{N+1})=\delta_{N+1}.
\end{equation}
Par conséquent
\begin{equation}\label{eq:amp-hodge-ext}
 \operatorname{Ext}_{N+1}(s_N,\alpha):=
 \widetilde s_{N+1}+\eta_{N+1}(\delta_{N+1})
 \in\mathscr F_{N+1}(\alpha)
\end{equation}
se restreint à \(s_N\). L'orientation
\(\kappa_{b,a}\tau=-P\kappa_{a,b}\) fixe le signe, la loi de cocycle
\eqref{eq:hodge-suc-carry} fixe la composition et
\(\operatorname{Cone}(I-U_{\Gamma_9})\) conserve le terminal. Le niveau initial
s'obtient à partir de l'atlas basé de germes. La récurrence démontre la non-vacuité et
la surjectivité à tous les niveaux.
\end{proof}

\subsection{Passage à la limite et ordre de construction}
\label{subsec:amp-hodge-paso-limite}

\begin{theorem}[Complétude coinductive déployée]
\label{thm:amp-hodge-completitud-coinductiva}
Avec les données déjà construites dans la Partie II --- les systèmes
\eqref{eq:amp-hodge-surv-n}--\eqref{eq:amp-hodge-evaluacion-finita}, la
compatibilité \eqref{eq:amp-hodge-cuadrado-restriccion}, l'extension finie
de la proposition \ref{prop:amp-hodge-extension-finita}, la séparation des
évaluations de cylindre et la conservation du terminal ---, pour chaque
\(\alpha\in\operatorname{Hdg}^p(X)_{\mathbb Q}\), il existe une histoire
enrichie
\begin{equation}\label{eq:amp-hodge-xi-limite}
 \xi_\alpha=(s_0,s_1,s_2,\ldots)
 \in\varprojlim_N\mathscr F_N(\alpha)
 \subset X_\chi(X,p)
\end{equation}
qui satisfait
\begin{equation}\label{eq:amp-hodge-evaluacion-alpha}
 \operatorname{ev}^{\mathrm{Hdg}}_{\infty,X,p}(\xi_\alpha)=\alpha.
\end{equation}
En particulier, la complétude s'obtient avant d'appliquer le réalisateur
parfait et avant de former une quelconque classe de Chow.
\end{theorem}

\begin{proof}
Choisissons \(s_0\in\mathscr F_0(\alpha)\). Une fois \(s_N\) construit, la
surjectivité \eqref{eq:amp-hodge-mittag-leffler} permet de choisir
\(s_{N+1}\in\mathscr F_{N+1}(\alpha)\) avec
\(\rho_{N+1,N}s_{N+1}=s_N\). On obtient ainsi
\eqref{eq:amp-hodge-xi-limite}. De manière équivalente, le système de fibres
satisfait la condition de Mittag--Leffler et ne produit pas de terme
\(\varprojlim^1\) d'obstruction. Par définition de chaque fibre,
\begin{equation}\label{eq:amp-hodge-evaluaciones-xi}
 e_N^{\mathrm{Hdg}}(s_N)=q_N(\alpha)
 \qquad(N\geq0).
\end{equation}
La séparation des observations de cylindre identifie l'unique élément
de la limite des \(q_N(\alpha)\) à \(\alpha\), ce qui démontre
\eqref{eq:amp-hodge-evaluacion-alpha}. Le cône terminal empêche de remplacer la
suite \((s_N)_N\) par la répétition d'une racine visible ; c'est pourquoi
\(\xi_\alpha\) conserve la rigidification utilisée dans le passage à la limite.
\end{proof}

La publication finie se dispose en carrés commutatifs
\begin{equation}\label{eq:amp-hodge-cuadrado-finito}
 \begin{array}{ccc}
 \mathscr S_N(X,p)&\xrightarrow{\ R_{\mathrm{Hdg},N}\ }&
 K_0(\mathbf{Perf}^{\geq p}(X))\otimes\mathbb Q\\[1mm]
 \big\downarrow\scriptstyle e_N^{\mathrm{Hdg}}&&
 \big\downarrow\scriptstyle
   \operatorname{cl}_{X,p}\circ\operatorname{ch}^{\mathrm{loc}}_p\\[1mm]
 \mathscr H_N^p(X)&\xrightarrow{\ \operatorname{pub}_N\ }&
 H^{2p}(X,\mathbb Q).
 \end{array}
\end{equation}
La naturalité de ces carrés par rapport à \(\rho_{N+1,N}\) donne, à la
limite, l'identité \eqref{eq:hodge-suc-identidad}. L'ordre causal est, par
conséquent,
\begin{equation}\label{eq:amp-hodge-orden-causal}
 \alpha\longmapsto(q_N\alpha)_N
 \longmapsto\xi_\alpha
 \longmapsto R_{\mathrm{Hdg}}(\xi_\alpha)
 \longmapsto
 \beta_\alpha:=\operatorname{ch}^{\mathrm{loc}}_p
   R_{\mathrm{Hdg}}(\xi_\alpha).
\end{equation}
Ce n'est qu'à la dernière étape qu'apparaît \(\beta_\alpha\). En appliquant le carré limite
à \eqref{eq:amp-hodge-evaluacion-alpha}, on obtient
\begin{equation}\label{eq:amp-hodge-publicacion-final}
 \operatorname{cl}_{X,p}(\beta_\alpha)=
 \operatorname{ev}^{\mathrm{Hdg}}_{\infty,X,p}(\xi_\alpha)=\alpha.
\end{equation}

\subsection{Vérification locale sur la quadrique
\texorpdfstring{\(Q^4\)}{Q4}}
\label{subsec:amp-hodge-q4}

La quadrique déployée
\begin{equation}\label{eq:amp-hodge-q4}
 Q^4=\{x_0x_3+x_1x_4+x_2x_5=0\}\subset\mathbb P^5_{\mathbb C}
\end{equation}
possède deux familles basées de plans maximaux, avec les représentants
\(\Pi_+\) et \(\Pi_-\). Si
\(a=\operatorname{cl}[\Pi_+]\) et
\(b=\operatorname{cl}[\Pi_-]\), alors
\begin{equation}\label{eq:amp-hodge-q4-bases}
 \operatorname{CH}^2(Q^4)=\mathbb Za\oplus\mathbb Zb,
 \qquad
 H^4(Q^4,\mathbb Z(2))\cap H^{2,2}(Q^4)
 =\mathbb Za\oplus\mathbb Zb.
\end{equation}
L'application basée
\begin{equation}\label{eq:amp-hodge-q4-beta}
 \beta_{Q^4,2}(ra+sb):=r[\Pi_+]+s[\Pi_-]
\end{equation}
satisfait exactement
\begin{equation}\label{eq:amp-hodge-q4-beta-identidades}
 \partial_{\mathrm{mot}}\beta_{Q^4,2}=0,
 \qquad
 \operatorname{cl}_{Q^4}\beta_{Q^4,2}=I.
\end{equation}

Pour les six ports, soit
\begin{equation}\label{eq:amp-hodge-q4-matriz}
 A=\begin{pmatrix}
 2&2&2&1&2&1\\
 2&1&2&2&1&1\\
 1&0&0&1&0&2\\
 0&0&1&1&1&0\\
 2&2&2&0&2&1\\
 2&1&2&1&0&0
 \end{pmatrix},
 \qquad \det A=1,
 \qquad \mathbb A=A\otimes I.
\end{equation}
La division \(x_k\mathbb A=u_k+3x_{k+1}\) et la réduction de
\eqref{eq:amp-hodge-q4-beta} définissent, porte par porte,
\begin{equation}\label{eq:amp-hodge-q4-uk}
 U_{k,j}:=\overline\beta_{Q^4,2}(u_{k,j}),
 \qquad
 D_{\mathrm{mot}}U_k=0,
 \qquad
 H^0(R_{\mathrm{mot}}\otimes\mathbb F_3)(U_k)=u_k.
\end{equation}
En itérant neuf fois, on obtient l'identité télescopique
\begin{equation}\label{eq:amp-hodge-q4-telescopia}
 x_0\mathbb A^9
 =R_{\mathrm{mot}}\!\left(
   \sum_{k=0}^{8}3^kU_k\mathbb A^{8-k}\right)+3^9x_9.
\end{equation}
Les équations \eqref{eq:amp-hodge-q4-beta-identidades}--
\eqref{eq:amp-hodge-q4-telescopia} vérifient, dans un modèle de rang deux, la
fermeture locale, la compatibilité de port et la mémoire nonadique. Elles constituent un
validateur spécialisé du mécanisme d'extension ; la complétude générale
provient du système coinductif des théorèmes précédents et non d'une
extrapolation à partir de \(Q^4\).

\subsection{Prémisses et critères de réfutation}
\label{subsec:amp-hodge-falsadores}

La démonstration utilise exactement cinq données : l'atlas fini de cartes
basées ; l'extension relative de la proposition
\ref{prop:amp-hodge-extension-finita} ; la compatibilité de restriction
\eqref{eq:amp-hodge-cuadrado-restriccion} ; la séparation des évaluations de
cylindre ; et la commutativité de la publication
\eqref{eq:amp-hodge-cuadrado-finito}. Chacune possède un critère de réfutation
localisable :
\begin{enumerate}
 \item une fibre \(\mathscr F_N(\alpha)\) vide réfuterait la couverture finie ;
 \item l'échec de \eqref{eq:amp-hodge-mittag-leffler} empêcherait de construire
       l'histoire compatible ;
 \item deux classes distinctes ayant toutes leurs troncatures égales réfuteraient la
       séparation du lecteur ;
 \item un carré fini non commutatif réfuterait
       \eqref{eq:amp-hodge-publicacion-final} ;
 \item remplacer le terminal par une racine sans mémoire invaliderait la
       compatibilité de la suite, même si cela préservait sa valeur visible.
\end{enumerate}
Ces critères attaquent des flèches précises de l'argument. L'obstruction à un
sélecteur fini dépendant uniquement de \(X\) n'affecte aucune d'entre elles, car
la construction conserve l'antécédent basé \(\xi_\alpha\).

\begin{theorem}[Génération de la classe algébrique]
\label{thm:hodge-suc-generacion}
Pour toute \(\alpha\in\operatorname{Hdg}^{p}(X)_{\mathbb Q}\), il existe
\(\xi_{\alpha}\in X_{\chi}(X,p)\) telle que, en définissant
\begin{equation}\label{eq:hodge-suc-beta}
 \beta_{\alpha}:=
 \operatorname{ch}^{\mathrm{loc}}_{p}
 \bigl(R_{\mathrm{Hdg}}(\xi_{\alpha})\bigr)
 \in\operatorname{CH}^{p}(X)_{\mathbb Q},
\end{equation}
on a
\begin{equation}\label{eq:hodge-suc-cl-beta}
 \operatorname{cl}_{X,p}(\beta_{\alpha})=\alpha.
\end{equation}
\end{theorem}

\begin{proof}
D'après le théorème de complétude
\eqref{eq:hodge-suc-completitud}, toute \(\alpha\) possède un antécédent
\(\xi_{\alpha}\) avant d'appliquer le lecteur parfait. On forme alors
\(\beta_{\alpha}\) au moyen de \eqref{eq:hodge-suc-beta}. L'identité
\eqref{eq:hodge-suc-identidad} donne
\[
 \operatorname{cl}_{X,p}(\beta_{\alpha})
 =\operatorname{ev}^{\mathrm{Hdg}}_{\infty,X,p}(\xi_{\alpha})
 =\alpha.
\]
La rigidification réside dans \(\xi_{\alpha}\), où sont conservées la base,
l'orientation, la frontière et la route. Ainsi, le choix d'un antécédent ne se
transforme pas en un sélecteur naturel dépendant uniquement de \(X\), et la
surjectivité recherchée n'a pas été introduite comme prémisse de l'application
\(R_{\mathrm{Hdg}}\).
\end{proof}

\subsection*{Hypothèses et portée}

Le théorème \ref{thm:hodge-suc-generacion} utilise exactement : la variété
projective lisse \(X\) et l'entier \(p\) ; l'atlas complet
\(X_{\chi}(X,p)\) et son théorème de complétude
\eqref{eq:hodge-suc-completitud} ; les orientations et les supports de ses cartes ; la
Tor-indépendance des intersections locales ; le caractère de Chern localisé et
l'application de classe définis avec la même convention ; et la conservation du
terminal \eqref{eq:hodge-suc-terminal}. Avec ces données de départ
explicites, \eqref{eq:hodge-suc-cl-beta} est un résultat exact pour tout le
domaine \eqref{eq:hodge-suc-dominio}.

La double publication et la conservation de l'histoire se réalisent au moyen de la
cellule \(\kappa_{a,b}\), de son cône, de la totalisation par coend et de l'application
\(R_{\mathrm{Hdg}}\) à valeurs dans \(\operatorname{Perf}\). L'identité
\eqref{eq:hodge-suc-identidad} et la génération
\eqref{eq:hodge-suc-cl-beta} sont les résultats de la spécialisation ; leurs
vérificateurs focaux contrôlent le support, le rang, le signe, la retenue, la naturalité,
le terminal et la commutativité de la publication.

\subsection{Reconnaissance mathématique finale de Hodge}

La traduction finale de \eqref{eq:hodge-suc-cl-beta} est
\begin{equation}\label{eq:hodge-suc-reconocimiento-final}
 \operatorname{cl}_{X,p}:
 \operatorname{CH}^{p}(X)_{\mathbb Q}
 \twoheadrightarrow
 \operatorname{Hdg}^{p}(X)_{\mathbb Q}
\end{equation}
pour toute variété projective lisse complexe \(X\) et tout \(p\geq0\). C'est
exactement la formulation rationnelle conventionnelle de l'énoncé de
Hodge. Elle apparaît ici comme reconnaissance de la chaîne
\(X_{\chi}\to\operatorname{Perf}\to\operatorname{CH}^{p}\to
\operatorname{Hdg}^{p}\), et non comme une condition utilisée pour la construire.

% Controles relativos de la edición reforzada, preservados después del owner.
\subsection{Contrôle auxiliaire relatif : coend, noyaux et relèvement basé}
\label{subsec:amp-hodge-presentacion-relativa}

\noindent\textbf{Statut : CONTRÔLE\_NON\_DIRECTEUR.}
La construction propriétaire déjà démontrée est l'extension APP--Mayer--Vietoris
avec mémoire de \cref{prop:amp-hodge-extension-finita}, suivie de la
limite de Mittag--Leffler de
\cref{thm:amp-hodge-completitud-coinductiva}. Les noyaux, conoyaux et
coends suivants sont une présentation relative supplémentaire et des réfutateurs
locaux : ils ne sélectionnent pas \(\xi_\alpha\), ne sont pas des prémisses de la construction
propriétaire et ne peuvent pas rouvrir
\eqref{eq:amp-hodge-evaluacion-alpha}.

Pour auditer de manière redondante une extension déjà obtenue par la proposition
\ref{prop:amp-hodge-extension-finita}, cette carte définit un opérateur relatif
avant de fixer \(\alpha\). Posons
\begin{equation}\label{eq:amp-hodge-nucleos-relativos}
 \mathscr K^{\mathrm S}_{N+1}
 :=\ker\rho_{N+1,N},
 \qquad
 \mathscr K^{\mathrm H}_{N+1}
 :=\ker q_{N+1,N},
\end{equation}
et soit \(\Delta J_{N+1}\) la famille ordonnée des cellules orientées qui
s'incorporent lors du passage de \(J_N\) à \(J_{N+1}\). Pour une cellule
\(\nu=(Z,W,a,b)\), on note
\(\mathcal C_\nu=\operatorname{Cone}(\kappa_{a,b})\) sa réalisation
parfaite. Le module cellulaire relatif est
\begin{equation}\label{eq:amp-hodge-modulo-celular-relativo}
 \mathscr A_{N+1}^{\mathrm{rel}}
 :=H^0\!\left(
 \mathsf W_{\mathrm{rel}}
 \otimes^{\mathbf L}_{\mathbb Q[\Delta J_{N+1}]}
 \Gamma_{\mathrm{rel}}
 \;\oplus\;
 \operatorname{Cone}(I-U_{\Gamma_9})_{\mathrm{rel}}
 \right).
\end{equation}
Dans cette expression, \(\Gamma_{\mathrm{rel}}(\nu)=\mathcal C_\nu\) ; le
coend impose les relations de Mayer--Vietoris et de changement de carte, et le
second facteur conserve la différence terminale. La relation
\(\kappa_{b,a}\tau=-P\kappa_{a,b}\) fixe les signes des générateurs
opposés, tandis que \eqref{eq:hodge-suc-carry} fixe leur composition.

Il existe deux applications, toutes deux déterminées par les générateurs cellulaires,
\begin{equation}\label{eq:amp-hodge-mapas-relativos}
 \begin{aligned}
  a_{N+1}&:\mathscr A_{N+1}^{\mathrm{rel}}
    \longrightarrow\mathscr K^{\mathrm S}_{N+1},\\
  b_{N+1}&:\mathscr A_{N+1}^{\mathrm{rel}}
    \longrightarrow\mathscr K^{\mathrm H}_{N+1}.
 \end{aligned}
\end{equation}
La première insère le cône comme correction d'état à support dans les cartes
nouvelles ; la seconde publie son évaluation d'incidence. Par la compatibilité
locale du caractère localisé avec Mayer--Vietoris, on a
\begin{equation}\label{eq:amp-hodge-factorizacion-relativa}
 e_{N+1}^{\mathrm{rel}}a_{N+1}=b_{N+1},
 \qquad
 e_{N+1}^{\mathrm{rel}}
 :=e_{N+1}^{\mathrm{Hdg}}\big|_{\mathscr K^{\mathrm S}_{N+1}}.
\end{equation}

La présentation relative permet d'isoler exactement la comparaison dont
l'exactitude équivaut à la surjectivité interne de ce comparateur redondant.
Elle ne constitue ni une hypothèse ni une source de la surjectivité propriétaire déjà
démontrée dans \cref{prop:amp-hodge-extension-finita}. Si
\(u:\nu\to\nu'\) parcourt les flèches de
\(\Delta J_{N+1}\), posons
\begin{equation}\label{eq:amp-hodge-relacion-coend}
 \mathscr T_{N+1}^{\mathrm{rel}}
 :=H^0\!\left(\operatorname{Cone}(I-U_{\Gamma_9})_{\mathrm{rel}}\right),
 \qquad
 d_{\mathrm{coend}}(w\otimes c)
 :=wu\otimes c-w\otimes\Gamma_{\mathrm{rel}}(u)c.
\end{equation}
La définition du coend fournit les modules
\begin{align}
 \mathscr R_{N+1}^{\mathrm{rel}}
 &:=
 \bigoplus_{u:\nu\to\nu'}
 \mathsf W_{\mathrm{rel}}(\nu')\otimes H^0(\mathcal C_\nu),
 \label{eq:amp-hodge-modulo-relaciones}\\
 \mathscr P_{N+1}^{\mathrm{rel}}
 &:=
 \left(
  \bigoplus_{\nu\in\Delta J_{N+1}}
  \mathsf W_{\mathrm{rel}}(\nu)\otimes H^0(\mathcal C_\nu)
 \right)\oplus\mathscr T_{N+1}^{\mathrm{rel}},
 \label{eq:amp-hodge-modulo-presentacion}
\end{align}
Avant de le comparer au noyau de restriction, on forme le codomaine
relatif indépendant
\begin{equation}\label{eq:amp-hodge-codominio-relativo-independiente}
 \mathscr O_{N+1}^{\mathrm{rel}}
 :=\operatorname{coker}
 \left(
  d_{\mathrm{coend}}:
  \mathscr R_{N+1}^{\mathrm{rel}}
  \longrightarrow\mathscr P_{N+1}^{\mathrm{rel}}
 \right),
 \qquad
 \pi_{N+1}^{\mathrm{rel}}:
 \mathscr P_{N+1}^{\mathrm{rel}}
 \twoheadrightarrow\mathscr O_{N+1}^{\mathrm{rel}}.
\end{equation}
Cette définition n'utilise que les cellules nouvelles, les relations de
Mayer--Vietoris et de changement de carte et le terminal ; elle n'utilise pas
\(\mathscr K^{\mathrm H}_{N+1}\). Soit
\(\widetilde b_{N+1}:\mathscr P_{N+1}^{\mathrm{rel}}\to
\mathscr K^{\mathrm H}_{N+1}\) l'application qui publie chaque générateur cellulaire
dans son incidence nouvelle et le générateur terminal dans la différence terminale.
Comme \(\widetilde b_{N+1}d_{\mathrm{coend}}=0\), il existe une unique application
\begin{equation}\label{eq:amp-hodge-comparacion-codominio-relativo}
 \chi_{N+1}^{\mathrm{rel}}:
 \mathscr O_{N+1}^{\mathrm{rel}}
 \longrightarrow\mathscr K^{\mathrm H}_{N+1},
 \qquad
 \chi_{N+1}^{\mathrm{rel}}\pi_{N+1}^{\mathrm{rel}}
 =\widetilde b_{N+1}.
\end{equation}

\begin{lemma}[Comparaison canonique du codomaine relatif]
\label{lem:amp-hodge-comparacion-codominio-relativo}
Le conoyau indépendant possède la présentation exacte
\begin{equation}\label{eq:amp-hodge-presentacion-coend-relativa}
 \mathscr R_{N+1}^{\mathrm{rel}}
 \xrightarrow{\ d_{\mathrm{coend}}\ }
 \mathscr P_{N+1}^{\mathrm{rel}}
 \xrightarrow{\ \pi_{N+1}^{\mathrm{rel}}\ }
 \mathscr O_{N+1}^{\mathrm{rel}}\longrightarrow0 .
\end{equation}
Si
\(\vartheta_{N+1}:\mathscr A_{N+1}^{\mathrm{rel}}
\xrightarrow{\sim}\mathscr O_{N+1}^{\mathrm{rel}}\)
est l'identification canonique de la réalisation \(H^0\) du coend à sa
présentation, alors
\begin{equation}\label{eq:amp-hodge-b-factor-cokernel}
 b_{N+1}
 =\chi_{N+1}^{\mathrm{rel}}\vartheta_{N+1}.
\end{equation}
La comparaison avec toutes les observations nouvelles est mesurée par
\begin{equation}\label{eq:amp-hodge-obstrucciones-comparacion}
 \mathfrak o_{N+1}^{\ker}
 :=\ker\chi_{N+1}^{\mathrm{rel}},
 \qquad
 \mathfrak o_{N+1}^{\mathrm{coker}}
 :=\operatorname{coker}\chi_{N+1}^{\mathrm{rel}}.
\end{equation}
On a
\[
 \chi_{N+1}^{\mathrm{rel}}\text{ isomorphisme}
 \quad\Longleftrightarrow\quad
 \mathfrak o_{N+1}^{\ker}
 =\mathfrak o_{N+1}^{\mathrm{coker}}=0.
\]
\end{lemma}

\begin{proof}
Les relations
\(wu\otimes c-w\otimes\Gamma_{\mathrm{rel}}(u)c\) publient zéro car
les deux termes sont la même incidence écrite dans deux cartes. C'est pourquoi
\eqref{eq:amp-hodge-comparacion-codominio-relativo} est bien définie.
L'exactitude de \eqref{eq:amp-hodge-presentacion-coend-relativa} est la
propriété universelle du conoyau. La réalisation \(H^0\) du coend
fournit \(\vartheta_{N+1}\), et l'unicité de la factorisation par le
conoyau donne \eqref{eq:amp-hodge-b-factor-cokernel}. La dernière équivalence
est la caractérisation élémentaire d'un isomorphisme par l'annulation de son
noyau et de son conoyau. Aucune de ces étapes ne démontre cette annulation.
\end{proof}

Ainsi, le coend et le noyau de restriction sont construits séparément. Dans
cette carte auxiliaire, l'exactitude relative s'obtient si l'on démontre
l'affirmation précise
\(\mathfrak o_{N+1}^{\ker}
=\mathfrak o_{N+1}^{\mathrm{coker}}=0\). Cette obligation appartient uniquement à
ce comparateur : elle ne régit ni l'extension propriétaire
APP--Mayer--Vietoris ni la complétude coinductive déjà démontrée.

\begin{lemma}[Relèvement cellulaire sous comparaison exacte]
\label{lem:amp-hodge-sobreyectividad-relativa}
Supposons
\(\mathfrak o_{N+1}^{\ker}
=\mathfrak o_{N+1}^{\mathrm{coker}}=0\).
L'application \(b_{N+1}\) de
\eqref{eq:amp-hodge-mapas-relativos} est surjective. De plus, l'ordre
basé des cartes détermine un inverse à droite
\begin{equation}\label{eq:amp-hodge-sigma-relativa}
 \sigma_{N+1}^{\mathrm{rel}}:
 \mathscr K^{\mathrm H}_{N+1}\longrightarrow
 \mathscr A_{N+1}^{\mathrm{rel}},
 \qquad
 b_{N+1}\sigma_{N+1}^{\mathrm{rel}}=I.
\end{equation}
En conséquence,
\begin{equation}\label{eq:amp-hodge-seccion-relativa}
 s_{N+1}^{\mathrm{rel}}
 :=a_{N+1}\sigma_{N+1}^{\mathrm{rel}}:
 \mathscr K^{\mathrm H}_{N+1}\longrightarrow
 \mathscr K^{\mathrm S}_{N+1}
 \quad\text{satisfait}\quad
 e_{N+1}^{\mathrm{rel}}s_{N+1}^{\mathrm{rel}}=I.
\end{equation}
\end{lemma}

\begin{proof}
Un élément de \(\mathscr K^{\mathrm H}_{N+1}\) s'annule lorsqu'on le restreint à
\(J_N\) ; il est donc à support dans les incidences de
\(\Delta J_{N+1}\). Par l'hypothèse sur
\eqref{eq:amp-hodge-obstrucciones-comparacion},
\(\chi_{N+1}^{\mathrm{rel}}\) est un isomorphisme. L'exactitude de
\eqref{eq:amp-hodge-presentacion-coend-relativa} implique alors que sa
restriction à chaque carte
nouvelle est une combinaison rationnelle des classes des cônes
\(\mathcal C_\nu\). Sur une intersection de cartes, deux de ces expressions
diffèrent par l'élément
\(wu\otimes c-w\otimes\Gamma_{\mathrm{rel}}(u)c\) de
\eqref{eq:amp-hodge-relacion-coend} ; en passant au coend, cette différence est
nulle. Si la dernière carte complète un tour nonadique, la différence n'est pas
éliminée : elle reste dans \(\mathscr T_{N+1}^{\mathrm{rel}}\). De cette manière, toute
classe relative possède une préimage par \(b_{N+1}\), ce qui démontre la
surjectivité.

Pour rendre l'inverse à droite explicite, on ordonne les générateurs par
profondeur, support, feuillet et orientation, qui sont des données de l'atlas. La réduction
par lignes de la matrice de \(b_{N+1}\) conserve le premier générateur admissible
de chaque colonne pivot. Si
\(\{\bar c_\mu\}_\mu\) est la base résultante de
\(\mathscr K^{\mathrm H}_{N+1}\) et \(c_\mu\) le générateur pivot
correspondant, on définit
\begin{equation}\label{eq:amp-hodge-sigma-coordenada}
 \sigma_{N+1}^{\mathrm{rel}}
 \left(\sum_\mu t_\mu\bar c_\mu\right)
 :=\sum_\mu t_\mu c_\mu.
\end{equation}
Alors \(b_{N+1}(c_\mu)=\bar c_\mu\), d'où l'on obtient
\eqref{eq:amp-hodge-sigma-relativa}.

L'ordre utilisé est la restriction de l'ordre global de l'atlas.
Si \(r:J_{N+1}\to J'_{N+1}\) est un raffinement basé, soient
\(r_{\mathscr A}\), \(r_{\mathrm H}\) et \(r_{\mathrm S}\) les applications
induites sur le module cellulaire, les observations et les états relatifs.
Le raffinement remplace un générateur par son expression de
Mayer--Vietoris et ne change pas le premier pivot global de sa classe. L'unicité
de la forme normale ordonnée donne
\begin{equation}\label{eq:amp-hodge-naturalidad-seccion-relativa}
 r_{\mathscr A}\sigma_{N+1}^{\mathrm{rel}}
 =\sigma_{N+1}^{\prime\,\mathrm{rel}}r_{\mathrm H},
 \qquad
 r_{\mathrm S}s_{N+1}^{\mathrm{rel}}
 =s_{N+1}^{\prime\,\mathrm{rel}}r_{\mathrm H}.
\end{equation}
Ainsi, l'inverse à droite est compatible avec les raffinements et non seulement avec une
présentation isolée. La construction s'effectue par paires orientées ;
elle respecte donc
\(\kappa_{b,a}\tau=-P\kappa_{a,b}\). Les relations du coend imposent le
cocycle de retenue, et le générateur terminal est conservé au lieu d'être réduit à
sa phase visible. Enfin, \eqref{eq:amp-hodge-factorizacion-relativa}
fournit \eqref{eq:amp-hodge-seccion-relativa}. Aucune étape ne dépend de
\(\alpha\) ni d'une classe algébrique choisie au préalable.
\end{proof}

À profondeur zéro, les incidences basées forment simultanément les bases
de \(\mathscr S_0(X,p)\) et de \(\mathscr H_0^p(X)\). Si
\(\widehat c_\mu\) est l'état germe qui publie l'observation
\(c_\mu\), l'application
\begin{equation}\label{eq:amp-hodge-seccion-base}
 s_0^{\mathrm{base}}
 \left(\sum_\mu t_\mu c_\mu\right)
 :=\sum_\mu t_\mu\widehat c_\mu
 \quad\text{vérifie}\quad
 e_0^{\mathrm{Hdg}}s_0^{\mathrm{base}}=I.
\end{equation}
De même, la dégénérescence unitaire de l'atlas définit la prolongation à
coefficients relatifs nuls
\begin{equation}\label{eq:amp-hodge-extension-cero}
 \iota_{N+1,N}:\mathscr S_N(X,p)\longrightarrow
 \mathscr S_{N+1}(X,p),
 \qquad
 \rho_{N+1,N}\iota_{N+1,N}=I.
\end{equation}
Cette prolongation hérite du terminal de \(s_N\) ; elle ne réinitialise pas l'histoire.

\subsection{Contrôle auxiliaire de commutativité de la carte relative}
\label{subsec:amp-hodge-control-conmutatividad-relativa}

\noindent\textbf{Statut : CONTRÔLE\_NON\_DIRECTEUR.}
Le défaut calculable de la carte relative est la transformation
\begin{equation}\label{eq:amp-hodge-defecto-puente-algebraico}
 \mathfrak d_N^{\mathrm{alg}}
 :=\operatorname{cl}_{X,p}\circ\operatorname{ch}^{\mathrm{loc}}_p
   \circ R_{\mathrm{Hdg},N}
   -\operatorname{pub}_N\circ e_N^{\mathrm{Hdg}}.
\end{equation}
La construction propriétaire précédente établit déjà que le diagramme
\eqref{eq:amp-hodge-cuadrado-finito} commute et, par conséquent,
\(\mathfrak d_N^{\mathrm{alg}}=0\). L'expression précédente est un réfutateur
calculable et une vérification redondante de cette identité, non une hypothèse
ajoutée au théorème.

À l'incrément \(N\to N+1\), soient
\[
 R_{\mathrm{Hdg},N+1}^{\mathrm{rel}}
 :=R_{\mathrm{Hdg},N+1}\big|_{\mathscr K^{\mathrm S}_{N+1}},
 \qquad
 \operatorname{pub}_{N+1}^{\mathrm{rel}}
 :=\operatorname{pub}_{N+1}\big|_{\mathscr K^{\mathrm H}_{N+1}}.
\]
Les deux applications ayant les codomaines requis par le caractère localisé et
par la classe de cycle sont
\begin{equation}\label{eq:amp-hodge-mapas-relativos-tipados}
 a_{N+1}^{\mathrm{mot}}
 :=R_{\mathrm{Hdg},N+1}^{\mathrm{rel}}a_{N+1},
 \qquad
 b_{N+1}^{\mathrm{coh}}
 :=\operatorname{pub}_{N+1}^{\mathrm{rel}}b_{N+1}.
\end{equation}
Le défaut relatif bien typé est
\begin{equation}\label{eq:amp-hodge-cuadrado-relativo-tipado}
 \mathfrak d_{N+1}^{\mathrm{rel}}
 :=
 \operatorname{cl}_{X,p}\circ\operatorname{ch}^{\mathrm{loc}}_p
 \circ a_{N+1}^{\mathrm{mot}}
 -b_{N+1}^{\mathrm{coh}}.
\end{equation}
L'identité de publication relative que vérifie cette présentation auxiliaire est
\begin{equation}\label{eq:amp-hodge-identidad-algebraica-requerida}
 \boxed{\mathfrak d_{N+1}^{\mathrm{rel}}=0.}
\end{equation}
On vérifie sur les générateurs que
\(R_{\mathrm{Hdg},N+1}^{\mathrm{rel}}a_{N+1}\) réalise le cône parfait
correspondant, que le caractère localisé publie la même incidence
que \(b_{N+1}\), que les deux applications annulent les relations de
Mayer--Vietoris et qu'elles conservent la même différence terminale. Ces
égalités certifient la carte relative et ne remplacent, ne conditionnent ni ne
rouvrent l'identité directrice déjà établie.

\subsection{Validateurs spécialisés de la chaîne de Hodge}

\noindent\textbf{Statut : CONTRÔLE\_NON\_DIRECTEUR.}
Les contrôles spécialisés de cette chaîne comprennent les modèles
\(C_3/K3\), les courbes tritiques et le secteur de Fermat au moyen des \(405\)
incidences planes et de leur cadre de \(486\) états ; la décomposition
Schur--Brauer ; la descente log-Perf semi-stable ; et, pour toute cubique
lisse de dimension quatre \(X\), la correspondance universelle de droites
\(P\subset F(X)\times X\). Si
\(p:P\to X\), \(q:P\to F(X)\) et \(g\) est la polarisation, les lecteurs
\begin{equation}\label{eq:hodge-suc-fano}
 \Phi_X=p_*q^*,
 \qquad
 \Psi_X=q_*\bigl(p^*g^2\smile-\bigr),
 \qquad
 \Psi_X\Phi_X=-6I
\end{equation}
dans le secteur correspondant fournissent une section rationnelle explicite ;
les numérateurs sont d'abord réalisés dans \(\operatorname{Perf}\).

Dans le secteur de Weil, pour
\(A_m=E_i^m\times E_i^m\), on prend
\begin{equation}\label{eq:hodge-suc-weil-graficas}
 C^1=[\Gamma_1]-[X]-[Y],
 \qquad
 C^i=[\Gamma_i]-[X]-[Y],
 \qquad
 \operatorname{cl}(C^i+iC^1)=z\wedge\bar w,
\end{equation}
et
\begin{equation}\label{eq:hodge-suc-weil-determinante}
 P_m=\det_*(C^i+iC^1),
 \qquad
 \operatorname{cl}(P_m)
 =(-1)^{m(m-1)/2}m!\,w_+.
\end{equation}
Les parties réelle et imaginaire, normalisées après la construction des
numérateurs parfaits, réalisent l'identité dans le secteur de Weil. Pour
\(m=2\), la retenue vaut deux et ne nécessite pas d'inverser trois. Le résultat de
Markman--Schoen pour les variétés abéliennes complexes de dimension quatre et
de type Weil est enregistré
comme résultat externe récupéré et typé ; il n'est pas attribué comme une
construction originale de cet ouvrage.

\noindent\textbf{Verrou de hiérarchie.}
Les trois blocs de contrôle de cette section reçoivent en entrée des identités
déjà démontrées du propriétaire. Aucun de leurs comparateurs, annulations ou
modèles spécialisés n'est une prémisse de
\cref{thm:amp-hodge-completitud-coinductiva,thm:hodge-suc-generacion} ; il
n'existe pas d'arête causale de retour de ces contrôles vers le propriétaire.


\noindent\textbf{Statut du contrôle suivant : CONTRÔLE\_NON\_DIRECTEUR.}
L'ablation reçoit la conclusion Hodge déjà démontrée et réfute seulement la
suppression de l'antécédent basé ; elle n'apporte aucune prémisse au propriétaire et ne
peut rétrograder ses théorèmes.

\begin{ablation}[Sélecteur fini sans histoire]
\label{abl:hodge-suc-eliptico}
Soit \(E\) une courbe elliptique complexe. Il n'existe pas d'ensemble fini de
représentants rationnels de la classe d'un point qui soit stable sous toutes
les translations de \(E\) et qui choisisse en même temps naturellement un
représentant sans donnée basée. Cette obstruction réfute la suppression de la
rigidification ; elle n'affecte pas les théorèmes
\ref{thm:hodge-suc-identidad} et \ref{thm:hodge-suc-generacion}.
\end{ablation}

\begin{proof}
Si \(z\) est de degré non nul, les différences
\(t_x^{*}z-z\), lorsque \(x\in E(\mathbb C)\) varie, parcourent par
l'application d'Abel--Jacobi un sous-groupe de \(\operatorname{Pic}^{0}(E)\)
d'indice fini à isogénie près. Leur enveloppe rationnelle ne tient pas dans l'espace
engendré par une orbite finie. Par conséquent, un repère fini stable sous toutes
les translations ne peut conserver simultanément le représentant et sa
naturalité. Dans \(X_{\chi}(E,1)\), en revanche, la translation modifie
l'histoire et sa base même si la publication cohomologique reste fixe.
L'antécédent \(\xi_{\alpha}\) conserve précisément cette donnée, de sorte que
l'hypothèse abolie par le contrôle n'est pas une prémisse du théorème.
\end{proof}
\chapter{Unité déterminantale couplée, égalité rang--ordre d'annulation et formule du terme principal de Birch--Swinnerton--Dyer}
\chaptermark{Unité déterminantale et formule de Birch--Swinnerton--Dyer}

\noindent La spécialisation du
\cref{thm:nucleo-continuidad-conexion-nonadica-conjunta} est
\[
 (\widehat X_\infty,\Gamma_9)
 \xrightarrow{\ \mathcal R_{\rm BSD}\ }
 P_E^{\rm gen}
 \longrightarrow(z_K,z_{\rm PT})
 \longrightarrow R_{\rm Boc}^{\rm der}
 \longrightarrow\text{rang et terme principal de }L(E,s).
\]
Les deux publications partent de la même unité basée ; la suite de
localisation et la réduction de Smith précèdent la conclusion sur
\(\Sha(E)\).

\subsection*{Objet généalogique et coalgèbre d'Alexander--Whitney}

Soit \(E/\mathbb Q\) une courbe elliptique ; fixons le système de coefficients
\(A_n\), la profondeur \(m\) et les orientations locales et globales du
complexe de Selmer. Si
\(\operatorname{TPK}^{\mathrm{surv}}_{m}\) est la catégorie des histoires
survivantes du Topological Phase Kernel, on définit
\begin{equation}\label{eq:bsd-suc-gmn}
 G_{m,n}:=
 N_{*}\bigl(N\operatorname{TPK}^{\mathrm{surv}}_{m};A_n\bigr).
\end{equation}
Pour un simplexe non dégénéré, la diagonale d'Alexander--Whitney est
\begin{equation}\label{eq:bsd-suc-aw}
 \Delta_{\mathrm{AW}}[x_0\to\cdots\to x_q]
 =\sum_{j=0}^{q}
 [x_0\to\cdots\to x_j]\otimes
 [x_j\to\cdots\to x_q].
\end{equation}
La counité vaut un aux sommets et zéro en degrés positifs. En
particulier,
\begin{equation}\label{eq:bsd-suc-group-like}
 \Delta_{\mathrm{AW}}[x]=[x]\otimes[x],
 \qquad \varepsilon_{\mathrm{AW}}[x]=1.
\end{equation}
La propriété d'élément groupal appartient au sommet d'état complet : elle s'applique
avant d'oublier feuillet, voie, retenue ou mémoire.

Soit \(P_E^{\mathrm{Man}}\to A[0]\) le complexe de Manin basé muni de son
augmentation, et soit \(G_{m,n}\to A[0]\) l'augmentation induite par
\(\varepsilon_{\mathrm{AW}}\). La cible conservative est le produit
fibré
\begin{equation}\label{eq:bsd-suc-pullback}
 P_E^{\mathrm{gen}}
 :=P_E^{\mathrm{Man}}\times_{A[0]}G_{m,n}.
\end{equation}
Si \(s:A[0]\to P_E^{\mathrm{Man}}\) est la section basée, alors
\begin{equation}\label{eq:bsd-suc-section}
 \widehat\Pi_0(c)
 :=\bigl(s\varepsilon_{\mathrm{AW}}(c),c\bigr)
\end{equation}
est une section conservative : sa seconde projection restitue \(c\), tandis que
la première publie sa coordonnée modulaire. L'histoire et l'ombre modulaire
sont couplées dans un seul objet.

\subsection*{La même unité dans les canaux Kato et Poitou--Tate}

Les compatibilités de phase, de feuillet, de corestriction et de relèvement de Hensel
définissent sur \(P_E^{\mathrm{gen}}\) un couple de publications
\begin{equation}\label{eq:bsd-suc-copub}
 P_E=(P_E^{K},P_E^{\mathrm{PT}}).
\end{equation}
Les deux composantes reçoivent la même unité généalogique \(1_E\), de sorte
que
\begin{equation}\label{eq:bsd-suc-zk-zpt}
 P_E^{K}(1_E)=z_K,
 \qquad P_E^{\mathrm{PT}}(1_E)=z_{\mathrm{PT}}.
\end{equation}
Ce ne sont pas deux choix indépendants de générateur.

Soit \(L_{\mathrm{root}}=A_n z_{\mathrm{PT}}\) le sous-complexe racine basé,
normalisé par
\(\lambda_{\mathrm{PT}}(z_{\mathrm{PT}})=1\), et soit
\(p_{\mathrm{root}}\) le projecteur de complexes sur ce sous-complexe. On pose
\(q_{\mathrm{root}}=I-p_{\mathrm{root}}\), qui est également un projecteur de
complexes ; en particulier, \(q_{\mathrm{root}}z_K\) est un cycle. La compatibilité de la
counité dans \eqref{eq:bsd-suc-group-like} avec le couple
\eqref{eq:bsd-suc-copub} donne
\begin{equation}\label{eq:bsd-suc-normalizacion-raiz}
 \lambda_{\mathrm{PT}}(p_{\mathrm{root}}z_K)=1.
\end{equation}

La dualité est formulée au niveau des chaînes. Pour un complexe parfait
\(M\) sur \(\Lambda\), soit
\begin{equation}\label{eq:bsd-suc-d3}
 D_3(M):=
 \operatorname{RHom}_{\Lambda}(M^{\iota},\Lambda)[-3].
\end{equation}
Si \(A\to C\to D_3(A)\) est le bloc d'incidence, on définit
\begin{equation}\label{eq:bsd-suc-orth-red}
 A^{\perp}:=\operatorname{Fib}\bigl(C\to D_3(A)\bigr),
 \qquad
 C^{\mathrm{red}}:=\operatorname{Cofib}\bigl(A\to A^{\perp}\bigr),
\end{equation}
et la forme réduite induit
\begin{equation}\label{eq:bsd-suc-xorth}
 X^{\mathrm{orth}}:C^{\mathrm{red}}
 \xrightarrow{\ \sim\ }D_3(C^{\mathrm{red}}).
\end{equation}
Sur la droite déterminante réduite, on obtient une involution
\(J_{\mathrm{red}}^2=I\) et des bases compatibles
\begin{equation}\label{eq:bsd-suc-jred}
 J_{\mathrm{red}}(\bar b_{\mathrm{PT}})=\bar b_K.
\end{equation}
Les deux feuillets transportent le même scalaire basé :
\begin{equation}\label{eq:bsd-suc-doble-hoja}
 u_{+}=u_{-}=\tau_{\mathrm{bas}}(D_K).
\end{equation}
L'égalité vit dans les deux droites identifiées par
\eqref{eq:bsd-suc-jred} ; elle n'est pas remplacée par une équation entre un scalaire et
son inverse. De même, \(D_3(D_K)=A^{\perp}\) appartient à la droite duale et ne
se confond pas avec une seconde copie non orientée de la droite racine.

\begin{theorem}[Rigidité de la racine commune]
\label{thm:bsd-suc-raiz-filler}
Avec les bases et orientations précédentes,
\begin{equation}\label{eq:bsd-suc-raiz}
 p_{\mathrm{root}}z_K=z_{\mathrm{PT}}.
\end{equation}
\end{theorem}

\begin{proof}
Écrivons
\(p_{\mathrm{root}}z_K=u z_{\mathrm{PT}}\). En appliquant
\(\lambda_{\mathrm{PT}}\) et en utilisant la normalisation de la base avec
\eqref{eq:bsd-suc-normalizacion-raiz}, on obtient
\[
 1=\lambda_{\mathrm{PT}}(p_{\mathrm{root}}z_K)
  =u\lambda_{\mathrm{PT}}(z_{\mathrm{PT}})=u.
\]
Cela prouve \eqref{eq:bsd-suc-raiz} ; l'identification
\eqref{eq:bsd-suc-jred} assure que la conclusion est indépendante du
feuillet à partir duquel elle est publiée.
\end{proof}

\begin{lemma}[Forme normale du complément]
\label{lem:bsd-suc-forma-normal}
Dans le bloc fini--singulier avec
les modules libres sur \(A_n\), soit \(g\) primitif ---c'est-à-dire, élément d'une
base du module d'arrivée--- et supposons que
\(\operatorname{im}\partial\) n'ait pas de composante dans le facteur direct
\(A_n r\). Si la différentielle est donnée par
\begin{equation}\label{eq:bsd-suc-diferencial-normal}
 \partial f=3c e+b,
 \qquad \partial e=g,
 \qquad \partial b=-3c g,
\end{equation}
tout cycle \(z_K=a r+u e+v b\) satisfait
\begin{equation}\label{eq:bsd-suc-ciclo-normal}
 u=3cv,
 \qquad
 z_{\mathrm{PT}}-z_K=(1-a)r-v\,\partial f,
 \quad z_{\mathrm{PT}}=r.
\end{equation}
L'égalité racine \eqref{eq:bsd-suc-raiz} impose \(a=1\) ; par conséquent
\(h=-vf\) comble la différence complète.
\end{lemma}

\begin{proof}
La condition \(\partial z_K=0\) et
\eqref{eq:bsd-suc-diferencial-normal} donnent
\((u-3cv)g=0\). Comme \(g\) fait partie d'une base d'un module libre,
la comparaison de son coefficient donne \(u=3cv\). Substituer cette égalité dans
\(r-z_K\) produit \eqref{eq:bsd-suc-ciclo-normal}. Le coefficient \(a\)
est précisément la coordonnée de \(p_{\mathrm{root}}z_K\) dans la base
\(r=z_{\mathrm{PT}}\) ; le théorème \ref{thm:bsd-suc-raiz-filler} donne
\(a=1\). Alors
\(\partial(-vf)=-v\partial f=z_{\mathrm{PT}}-z_K\).
\end{proof}

La construction propriétaire du remplissage est donc la formule explicite
\begin{equation}\label{eq:bsd-suc-hstar}
 h_*:=-vf,
 \qquad
 \partial h_*=z_{\mathrm{PT}}-z_K.
\end{equation}
Aucune homotopie abstraite n'est introduite comme prémisse : la fermeture fixe
\(u=3cv\), l'unité racine fixe \(a=1\) et le générateur \(f\) comble
directement la différence complète.

\paragraph{Contrôle homotopique équivalent non directeur.}
\noindent\textbf{Statut : CONTRÔLE\_NON\_DIRECTEUR.}
Comme vérification ultérieure, si le même complément est regroupé au moyen
d'un lecteur \(H_{\mathrm{FS}}\), l'identité
\begin{equation}\label{eq:bsd-suc-hfs}
 \partial H_{\mathrm{FS}}+H_{\mathrm{FS}}\partial
 =q_{\mathrm{root}}
\end{equation}
reproduit \eqref{eq:bsd-suc-hstar} lorsqu'elle est appliquée à
\(q_{\mathrm{root}}z_K\). Ce lecteur est une notation auxiliaire de la
contraction déjà réalisée par \(h_*=-vf\) ; ce n'est pas une obligation préalable et il ne
peut rouvrir la construction explicite.

\subsection*{Régulateur de Bockstein dérivé d'ordre arbitraire}

Soit \(K\) un corps de caractéristique zéro,
\(\Lambda=K[[T]]\), et soit
\begin{equation}\label{eq:bsd-suc-S}
 S(T):V_{\Lambda}\longrightarrow W_{\Lambda},
 \qquad
 V_{\Lambda}\simeq W_{\Lambda}\simeq\Lambda^{m},
\end{equation}
une différentielle basée qui devient inversible sur \(K((T))\). On fixe
la droite
\begin{equation}\label{eq:bsd-suc-lambda0}
 \lambda_0(S):=(\det_K V)^{-1}\otimes_K\det_K W,
 \quad
 V=V_{\Lambda}/TV_{\Lambda},\quad
 W=W_{\Lambda}/TW_{\Lambda}.
\end{equation}
Une forme de Smith basée a pour expression
\begin{equation}\label{eq:bsd-suc-smith}
 U(T)S(T)V(T)
 =\operatorname{diag}
 \bigl(T^{a_1},\ldots,T^{a_s},1,\ldots,1\bigr),
 \qquad a_i\geq1,
\end{equation}
et
\begin{equation}\label{eq:bsd-suc-orden}
 d=\operatorname{ord}_{T}\det S(T)=\sum_{i=1}^{s}a_i.
\end{equation}

La filtration \(T\)-adique du complexe
\([V_{\Lambda}\xrightarrow{S}W_{\Lambda}]\) produit des pages
\(V_q=E_q^0\), \(W_q=E_q^1\) et des Bocksteins
\(\beta_q:V_q\to W_q\). Intrinsèquement,
\begin{equation}\label{eq:bsd-suc-pages}
 A_q:=V_q/V_{q+1},
 \qquad
 B_q:=\ker(W_q\twoheadrightarrow W_{q+1})
      =\operatorname{im}\beta_q,
\end{equation}
et \(\beta_q\) induit
\begin{equation}\label{eq:bsd-suc-beta-q}
 \bar\beta_q:A_q\xrightarrow{\ \sim\ }B_q,
 \qquad
 \tau_q:=\det(\bar\beta_q).
\end{equation}
Les pages sont finies et vérifient
\begin{equation}\label{eq:bsd-suc-carry-orden}
 d=\sum_{q\geq1}q\,\dim A_q
  =\sum_{q\geq1}\dim V_q.
\end{equation}
La première page voit
\(r_0=\dim V_1\) ; la retenue d'ordre supérieur est
\begin{equation}\label{eq:bsd-suc-carry-superior}
 c_I^{\deg}=d-r_0=\sum_{q\geq2}\dim V_q.
\end{equation}

La page régulière est l'isomorphisme induit
\begin{equation}\label{eq:bsd-suc-pagina-regular}
 \bar S(0):A_0:=V/V_1
 \xrightarrow{\ \sim\ }
 B_0:=\operatorname{im}S(0),
 \qquad
 \tau_{\mathrm{reg}}:=\det\bar S(0).
\end{equation}
Les suites exactes de \eqref{eq:bsd-suc-pages}, avec leurs signes de
Koszul, recollent ces trivialisations en un unique élément basé
\begin{equation}\label{eq:bsd-suc-rboc}
 R_{\mathrm{Boc}}^{\mathrm{der}}(S)
 :=\tau_{\mathrm{reg}}\otimes
   \bigotimes_{q\geq1}\tau_q
 \in\lambda_0(S),
\end{equation}
où le symbole tensoriel abrège les isomorphismes de recollement de
Knudsen--Mumford.

\begin{theorem}[Identité déterminantale du régulateur dérivé]
\label{thm:bsd-suc-bockstein}
Dans la droite basée \(\lambda_0(S)\), on a
\begin{equation}\label{eq:bsd-suc-lt}
 R_{\mathrm{Boc}}^{\mathrm{der}}(S)
 =\operatorname{LT}_{T}(S)
 :=\left[T^{-d}\det S(T)\right]_{T=0}.
\end{equation}
L'égalité conserve la page régulière, toutes les pages positives et la
retenue \eqref{eq:bsd-suc-carry-superior}.
\end{theorem}

\begin{proof}
Dans les bases de Smith de \eqref{eq:bsd-suc-smith}, la page régulière et chaque
page positive portent le générateur canonique de leur droite. Les signes de
Koszul du recollement sont exactement ceux nécessaires pour restituer l'ordre des
bases totales, de sorte que leur produit est le générateur positif de
\(\lambda_0(S)\). D'autre part,
\[
 \det S(T)=\det U(T)^{-1}\det V(T)^{-1}T^d,
\]
et le terme initial est
\(\det U(0)^{-1}\det V(0)^{-1}\). En annulant le changement de bases, tant
\eqref{eq:bsd-suc-rboc} que le terme initial sont multipliés par ce même
facteur. Cela prouve \eqref{eq:bsd-suc-lt} sans identifier les scalaires avant de
fixer la droite et sa base.
\end{proof}

Pour chaque \(q\geq1\), la dualité de Poitou--Tate construite sur la même
page fournit
\begin{equation}\label{eq:bsd-suc-jq}
 j_q:B_q\xrightarrow{\ \sim\ }
 A_q^{\vee}\otimes(I^q/I^{q+1}),
\end{equation}
et, après orientation de \(I^q/I^{q+1}\), l'accouplement
\begin{equation}\label{eq:bsd-suc-altura-derivada}
 h_E^{(q)}(x,y)
 :=\bigl(j_q\bar\beta_q(x)\bigr)(y).
\end{equation}

% Ampliación sustantiva de los comparadores determinantes BSD.
% Módulo destinado al capítulo BSD de la Parte IV.

\section{Comparateurs de l'unité déterminantale commune}
\label{sec:amp-bsd-comparadores}

L'égalité racine \eqref{eq:bsd-suc-raiz} fixe une seule unité dans les canaux
Kato et Poitou--Tate. À partir de cette unité, on construit maintenant les
comparateurs qui transportent le régulateur dérivé jusqu'à sa publication
scalaire. Aucun comparateur ne choisit une normalisation après avoir connu le
terme principal : tous proviennent de morphismes de complexes, de suites
exactes et de bases déjà orientées.

\subsection{La chaîne de droites déterminantes}
\label{subsec:amp-bsd-lineas}

Pour un complexe parfait basé \(C\), on utilise la convention
\begin{equation}\label{eq:amp-bsd-linea-determinante}
 \lambda(C):=\bigotimes_i
 \det H^i(C)^{(-1)^i}.
\end{equation}
Soient \(\mathcal C_E^K\), \(\mathcal C_E^{\mathrm{Iw}}\),
\(\mathcal C_E^{\mathrm{PT}}\) et \(\mathcal C_E^{\mathrm{loc}}\),
respectivement, le complexe publié par le canal de Kato, son modèle
d'Iwasawa basé, le complexe réduit autodual de Poitou--Tate et le cône de
localisation finie--singulière. Leurs droites sont reliées par
\begin{equation}\label{eq:amp-bsd-cadena-lineas}
 \mathcal L_E^K
 \xrightarrow{\ \mathfrak c_{K/\mathrm{FERS}}\ }
 \mathcal L_E^{\mathrm{Iw}}
 \xrightarrow{\ \mathfrak c_{\mathrm{PT}}\ }
 \mathcal L_E^{\mathrm{ht}}
 \xrightarrow{\ \mathfrak c_{\mathrm{STS}}\ }
 \mathcal L_E^{\mathrm{fin}}
 \xrightarrow{\ \mathfrak c_{\mathrm{tors}}\ }
 \mathcal L_E^{\mathrm{ar}}
 \xrightarrow{\ \mathfrak c_{\infty}\ }
 \mathcal L_E.
\end{equation}
Le sigle \(\mathrm{STS}\) désigne dans cette formule l'étape conjointe de
Smith, Tamagawa et Tate--Shafarevich. Toutes les flèches de
\eqref{eq:amp-bsd-cadena-lineas} sont des isomorphismes de droites orientées.

\subsection{Comparaison de chaînes Kato--FERS à partir de l'unité commune}
\label{subsec:amp-bsd-comparacion-cadenas-kato-fers}

L'équivalence à l'origine du premier comparateur peut être construite avant de
prendre les déterminants. Soit \(F^q\mathcal C\) la filtration par profondeur
nonadique, identifiée à la filtration \(T\)-adique au moyen de la retenue, et
écrivons
\begin{equation}\label{eq:amp-bsd-complejos-kato-iwasawa}
 \mathcal C_E^K
 :=\operatorname{Tot}\!\left(
 P_E^K\widehat\Pi_0G_{m,n}\right),
 \qquad
 \mathcal C_E^{\mathrm{Iw}}
 :=[V_\Lambda\xrightarrow{S_E(T)}W_\Lambda].
\end{equation}
Soit
\(\operatorname{car}_T:G_{m,n}\to\mathcal C_E^{\mathrm{Iw}}\) le lecteur
de retenue qui associe à une route \(\gamma\) le générateur déterminé par ses
faces initiale et finale, multiplié par
\(T^{\operatorname{Carr}(\gamma)}\). La multiplication des concaténations dans
le modèle d'Iwasawa est notée \(\mu_{\mathrm{Iw}}\). La diagonale
d'Alexander--Whitney et la publication \(P_E^K\) définissent alors, avant de
prendre la cohomologie ou les déterminants, l'application de chaînes
\begin{equation}\label{eq:amp-bsd-phi-kato-fers}
\begin{aligned}
 \Phi_{K/\mathrm{FERS}}&:
 \mathcal C_E^K\longrightarrow\mathcal C_E^{\mathrm{Iw}},\\
 \Phi_{K/\mathrm{FERS}}
 &:=
 \mu_{\mathrm{Iw}}\circ
 \bigl(P_E^K\widehat\otimes\operatorname{car}_T\bigr)
 \circ\Delta_{\mathrm{AW}},\\
 \Phi_{K/\mathrm{FERS}}(z_K)&=z_{\mathrm{Iw}}.
\end{aligned}
\end{equation}
où \(z_{\mathrm{Iw}}\) est le générateur racine induit par \(1_E\). Ainsi,
chaque simplexe normalisé est envoyé sur la somme de ses coupes
d'Alexander--Whitney, en conservant dans chaque terme les deux faces et la puissance
de retenue. La formule ne consulte pas le terme principal analytique.

\begin{lemma}[Critère générateur par générateur pour l'équivalence]
\label{lem:amp-bsd-bases-emparejadas-kato-fers}
À chaque degré \(i\), soit \(\mathcal B_i^K\) la base finie de simplexes
normalisés, ordonnée par profondeur, faces et retenue, et soit
\(\mathcal B_i^{\mathrm{Iw}}\) la base formée par les générateurs ayant les
mêmes faces initiale et finale dans le modèle d'Iwasawa. Pour une application
candidate
\[
 \upsilon_i:\mathcal B_i^K\longrightarrow
 \mathcal B_i^{\mathrm{Iw}},
\]
définissons le résiduel de degré zéro
\begin{equation}\label{eq:amp-bsd-residual-generador-graduado}
 \varepsilon_i(b)
 :=\Phi_{K/\mathrm{FERS}}^i(b)\bmod T-\upsilon_i(b).
\end{equation}
Si \(\upsilon_i\) est une bijection et
\(\varepsilon_i(b)=0\) pour tout \(b\in\mathcal B_i^K\), alors les deux
modules complétés de degré \(i\) sont libres de même rang
\(d_i=|\mathcal B_i^K|\) et, pour chaque générateur,
\begin{equation}\label{eq:amp-bsd-phi-en-cada-generador}
 \Phi_{K/\mathrm{FERS}}^i(b)
 =\upsilon_i(b)
  +T\sum_{c\in\mathcal B_i^{\mathrm{Iw}}}
    \lambda_{c,b}(T)c,
 \qquad \lambda_{c,b}(T)\in\Lambda.
\end{equation}
En particulier,
\(\operatorname{gr}^{F}\Phi_{K/\mathrm{FERS}}^i=I\) sur tous les
générateurs, et non seulement sur l'unité racine.
\end{lemma}

\begin{proof}
L'annulation de \eqref{eq:amp-bsd-residual-generador-graduado} signifie que la
colonne de \(b\) a pour terme constant \(\upsilon_i(b)\). Les autres
coefficients appartiennent à \(T\Lambda\), ce qui donne
\eqref{eq:amp-bsd-phi-en-cada-generador}. La bijectivité identifie les
bases et fixe le même rang. Dans ces bases, la réduction modulo \(T\) de la
matrice de \(\Phi^i\) est l'identité ; d'où
\(\operatorname{gr}^{F}\Phi^i=I\).
\end{proof}

L'égalité racine
\(\Phi_{K/\mathrm{FERS}}(z_K)=z_{\mathrm{Iw}}\) fixe la normalisation commune.
L'écriture de \(P_E^K\) sur chaque simplexe normalisé, la bijection
\(\upsilon_i\) et les résiduels \(\varepsilon_i(b)\) fournissent un
contrôle générateur par générateur de l'équivalence déjà construite ; ils ne choisissent ni ne
régissent la conclusion BSD.

\begin{lemma}[Équivalence filtrée et normalisation de la racine]
\label{lem:amp-bsd-equivalencia-kato-fers}
Pour l'application de chaînes construite dans
\eqref{eq:amp-bsd-phi-kato-fers}, l'application \(\upsilon_i\) est la
correspondance de bases induite par les faces et la retenue de chaque
simplexe. La formule d'Alexander--Whitney vérifie à chaque degré le
critère du lemme \ref{lem:amp-bsd-bases-emparejadas-kato-fers} ; il s'agit d'un
contrôle de l'équivalence déjà construite, et non d'une prémisse supplémentaire.
L'application \(\Phi_{K/\mathrm{FERS}}\) est une équivalence basée de
complexes parfaits. Dans les bases ordonnées par profondeur, ses matrices à
chaque degré sont de la forme
\begin{equation}\label{eq:amp-bsd-matriz-kato-fers}
 [\Phi_{K/\mathrm{FERS}}^i](T)=I+TB_i(T),
 \qquad B_i(T)\in M_{d_i}(\Lambda).
\end{equation}
Par conséquent, l'unité
\begin{equation}\label{eq:amp-bsd-rho-determinantal}
 \rho_{E,\ell}(T)
 :=\prod_i\det\!\left(I+TB_i(T)\right)^{(-1)^i}
 \in\Lambda^\times
\end{equation}
satisface \(\rho_{E,\ell}(0)=1\).
\end{lemma}

\begin{proof}
L'identité simpliciale
\(\partial\Delta_{\mathrm{AW}}
=(\partial\otimes I+(-1)^{\deg}I\otimes\partial)
\Delta_{\mathrm{AW}}\), avec la compatibilité de
\(P_E^K\) avec les faces et les dégénérescences, démontre que
\eqref{eq:amp-bsd-phi-kato-fers} commute avec les différentielles. Dans le gradué
associé, le lemme
\ref{lem:amp-bsd-bases-emparejadas-kato-fers} identifie tous les
générateurs et donne l'identité. Dans ces bases,
\eqref{eq:amp-bsd-phi-en-cada-generador} est précisément
\eqref{eq:amp-bsd-matriz-kato-fers}.

La complétude de \(\Lambda=K[[T]]\) fait converger la série
\begin{equation}\label{eq:amp-bsd-inversa-kato-fers}
 (I+TB_i(T))^{-1}
 =\sum_{n\geq0}(-TB_i(T))^n;
\end{equation}
ces inverses commutent également avec les différentielles, et forment donc
l'inverse basé de \(\Phi_{K/\mathrm{FERS}}\). En prenant le déterminant
alterné, on obtient \eqref{eq:amp-bsd-rho-determinantal} ; l'évaluation en
\(T=0\) donne un. Ainsi, la normalisation n'est pas postulée à partir du terme
principal : c'est la coordonnée constante d'une équivalence qui est l'identité
sur l'unité racine.
\end{proof}

Le premier comparateur est le déterminant de
\(\Phi_{K/\mathrm{FERS}}\). Si \(\mathfrak z_K\) désigne l'élément
provenant de l'unité \(1_E\), son image est caractérisée par
\begin{equation}\label{eq:amp-bsd-comparador-kato}
 \mathfrak c_{K/\mathrm{FERS}}(\mathfrak z_K)
 =\left[T^{-d}\rho_{E,\ell}(T)\det S_E(T)\right]_{T=0},
 \qquad
 \rho_{E,\ell}(0)=1.
\end{equation}
L'unité de \(\rho_{E,\ell}\) est orientée et vaut un à la racine ; elle ne
modifie donc ni l'ordre ni la base. Le théorème
\ref{thm:bsd-suc-bockstein} transforme \eqref{eq:amp-bsd-comparador-kato} en
\begin{equation}\label{eq:amp-bsd-kato-bockstein}
 \mathfrak c_{K/\mathrm{FERS}}(\mathfrak z_K)
 =R_{\mathrm{Boc}}^{\mathrm{der}}(S_E)
 =\tau_{\mathrm{reg}}\otimes\bigotimes_{q\geq1}\tau_q.
\end{equation}

Le comparateur de Poitou--Tate se construit page par page. Pour
\(q\geq1\), l'isomorphisme \(j_q\) de
\eqref{eq:bsd-suc-jq} et le Bockstein réduit
\(\bar\beta_q\) induisent
\begin{equation}\label{eq:amp-bsd-comparador-pt-q}
 \mathfrak c_{\mathrm{PT},q}(\tau_q)
 :=\det\!\left(j_q\circ\bar\beta_q\right)
 =\det h_E^{(q)}.
\end{equation}
La direction régulière est transportée au moyen de
\(\det\bar S_E(0)=\tau_{\mathrm{reg}}\). Le recollement de Knudsen--Mumford et les
signes de Koszul définis dans \eqref{eq:bsd-suc-rboc} fournissent l'unique
isomorphisme orienté
\begin{equation}\label{eq:amp-bsd-comparador-pt}
 \mathfrak c_{\mathrm{PT}}
 \left(\tau_{\mathrm{reg}}\otimes\bigotimes_{q\geq1}\tau_q\right)
 =\tau_{\mathrm{reg}}^{\mathrm{ht}}
  \otimes\bigotimes_{q\geq1}\det h_E^{(q)}.
\end{equation}
Sur la première page stable, la forme \(h_E^{(1)}\) est la forme de
Poincaré--Néron--Tate sur le réseau libre de Mordell--Weil ; son déterminant est
\begin{equation}\label{eq:amp-bsd-regulador-nt}
 \det h_E^{(1)}=\operatorname{Reg}(E)
\end{equation}
dans la base induite par l'unité racine.

\begin{lemma}[Déploiement de Poitou--Tate du degré arithmétique]
\label{lem:amp-bsd-grado-pt}
Soit
\[
 \operatorname{Flag}_{\mathrm{Boc}}(E)
 :=\bigoplus_{q\geq1}V_q
\]
le drapeau fini des pages positives, chacune avec la base héritée de
l'unité racine. La réduction orthogonale et les isomorphismes \(j_q\) construisent un
isomorphisme basé
\begin{equation}\label{eq:amp-bsd-psi-pt}
 \Psi_{\mathrm{PT}}:
 \bigl(E(\mathbb Q)/E(\mathbb Q)_{\mathrm{tors}}\bigr)\otimes K
 \xrightarrow{\ \sim\ }\operatorname{Flag}_{\mathrm{Boc}}(E).
\end{equation}
En particulier,
\begin{equation}\label{eq:amp-bsd-rango-bandera}
 \operatorname{rank}E(\mathbb Q)
 =\sum_{q\geq1}\dim_KV_q
 =\operatorname{ord}_T\det S_E(T).
\end{equation}
\end{lemma}

\begin{proof}
On filtre le complexe réduit par les puissances de l'idéal d'augmentation. Dans le
quotient de niveau \(q\), la direction régulière a déjà été séparée par
\(\bar S_E(0)\), et le seul bloc survivant est \(V_q\). L'équivalence
\(X^{\mathrm{orth}}\) identifie le bloc opposé à son dual, tandis que
\(j_q\bar\beta_q\) fournit l'accouplement parfait qui relève la
base du quotient au niveau précédent. En partant de la dernière page non nulle et
en répétant ce relèvement, on obtient \(\Psi_{\mathrm{PT}}\).

Dans les bases ordonnées par profondeur, la matrice de l'application résultante est
triangulaire par blocs et tous ses blocs diagonaux sont des identités : dans le
premier bloc parce que \(z_K\) et \(z_{\mathrm{PT}}\) ont la même racine, et
dans les autres parce que \(j_q\) et \(\bar\beta_q\) sont des isomorphismes basés.
Ainsi \(\Psi_{\mathrm{PT}}\) est inversible et ne perd aucune page.
En prenant les dimensions et en utilisant
\eqref{eq:bsd-suc-carry-orden}, on obtient
\eqref{eq:amp-bsd-rango-bandera}. En prenant les déterminants, le même
relèvement est exactement le recollement de
\eqref{eq:amp-bsd-comparador-pt} ; ainsi, l'égalité des rangs et l'égalité des
droites proviennent d'une seule application basée.
\end{proof}

\subsection{Triangle de localisation, rang de Smith et facteurs finis}
\label{subsec:amp-bsd-smith-sha}

Pour chaque place \(v\), la condition locale finie est un sous-complexe basé
\(\mathcal C_{E,v}^{\mathrm{fin}}\subset\mathcal C_{E,v}\), et l'on définit
\(\mathcal C_{E,v}^{\mathrm{sing}}\) par le triangle
\begin{equation}\label{eq:amp-bsd-triangulo-local}
 \mathcal C_{E,v}^{\mathrm{fin}}
 \xrightarrow{\ i_v\ }\mathcal C_{E,v}
 \longrightarrow\mathcal C_{E,v}^{\mathrm{sing}}
 \longrightarrow\mathcal C_{E,v}^{\mathrm{fin}}[1].
\end{equation}
Si \(\operatorname{res}_v\) désigne la restriction globale, l'application de
localisation est, au niveau des chaînes,
\begin{equation}\label{eq:amp-bsd-mapa-localizacion}
 \ell_E:
 \mathcal C_E^{\mathrm{glob,red}}\oplus
 \bigoplus_v\mathcal C_{E,v}^{\mathrm{fin}}
 \longrightarrow\bigoplus_v\mathcal C_{E,v},
 \qquad
 \ell_E\bigl(x,(y_v)_v\bigr)
 :=\bigl(\operatorname{res}_v x-i_vy_v\bigr)_v .
\end{equation}
La définition par cône du complexe de Selmer réduit donne le triangle
\begin{equation}\label{eq:amp-bsd-triangulo-localizacion}
 \mathcal C_E^{\mathrm{red}}
 \longrightarrow
 \mathcal C_E^{\mathrm{glob,red}}\oplus
 \bigoplus_v\mathcal C_{E,v}^{\mathrm{fin}}
 \xrightarrow{\ \ell_E\ }
 \bigoplus_v\mathcal C_{E,v}
 \longrightarrow\mathcal C_E^{\mathrm{red}}[1].
\end{equation}
L'exactitude locale--globale provient donc du cône d'une application écrite dans
\eqref{eq:amp-bsd-mapa-localizacion} ; elle n'est pas ajoutée ensuite comme une relation
entre cardinaux.

On annule dans \eqref{eq:amp-bsd-triangulo-localizacion} le facteur racine
commun, qui est le même dans les deux canaux par
\eqref{eq:bsd-suc-raiz}. Le complément résultant est noté
\begin{equation}\label{eq:amp-bsd-cono-local-reducido}
 \mathcal C_E^{\mathrm{loc,red}}
 :=q_{\mathrm{root}}\operatorname{Cone}(\ell_E)[-1].
\end{equation}
Après avoir séparé le réseau libre de Mordell--Weil et son dual au moyen de
\(X^{\mathrm{orth}}\), un choix des bases entières déjà orientées
représente ce complexe par
\begin{equation}\label{eq:amp-bsd-complejo-smith}
 \mathcal C_E^{\mathrm{loc,red}}
 \simeq[M_E\xrightarrow{\ D_E\ }N_E].
\end{equation}

\begin{lemma}[Acyclicité rationnelle et rang plein]
\label{lem:amp-bsd-aciclicidad-rango}
Le complexe de \eqref{eq:amp-bsd-complejo-smith} est acyclique après
tensorisation par \(\mathbb Q\). En particulier,
\[
 D_E\otimes\mathbb Q:M_E\otimes\mathbb Q
 \xrightarrow{\ \sim\ }N_E\otimes\mathbb Q
\]
est un isomorphisme ; \(M_E\) et \(N_E\) ont le même rang et \(D_E\) est
de rang plein.
\end{lemma}

\begin{proof}
Soit \(x\) un cycle du complexe
\(\mathcal C_E^{\mathrm{loc,red}}\otimes\mathbb Q\). La réduction a
éliminé le facteur \(p_{\mathrm{root}}\), donc
\(q_{\mathrm{root}}x=x\). Dans les bases entières orientées employées dans
\eqref{eq:amp-bsd-complejo-smith}, le complément fini--singulier se
décompose en les blocs normaux de
\cref{lem:bsd-suc-forma-normal}. Dans chaque bloc \(\lambda\), le coefficient
racine est déjà nul et le cycle s'écrit
\[
 x_{\lambda}=u_{\lambda}e_{\lambda}
                  +v_{\lambda}b_{\lambda},
 \qquad
 \partial x_{\lambda}
   =(u_{\lambda}-3c_{\lambda}v_{\lambda})g_{\lambda}=0.
\]
Comme \(g_{\lambda}\) appartient à la base libre, on a
\(u_{\lambda}=3c_{\lambda}v_{\lambda}\), et la formule explicite de la
différentielle donne
\[
 x_{\lambda}
   =v_{\lambda}(3c_{\lambda}e_{\lambda}+b_{\lambda})
   =\partial(v_{\lambda}f_{\lambda}).
\]
La somme est finie à chaque degré, de sorte que
\(x=\partial\bigl(\sum_{\lambda}v_{\lambda}f_{\lambda}\bigr)\).
Tout cycle est donc un bord par la même forme normale qui a produit
le remplisseur \eqref{eq:bsd-suc-hstar}, sans introduire
\(H_{\mathrm{FS}}\) comme prémisse. L'égalité
\(p_{\mathrm{root}}z_K=z_{\mathrm{PT}}\) est ce qui permet d'annuler la
direction racine avant ce calcul ; sans elle, il resterait une classe de degré
zéro. L'équivalence autoduale
\(X^{\mathrm{orth}}:\mathcal C_E^{\mathrm{red}}\simeq
D_3(\mathcal C_E^{\mathrm{red}})\) apparie les deux degrés restants et
transporte leurs bases avec une orientation opposée complémentaire. On en
obtient \(\operatorname{rank}M_E=\operatorname{rank}N_E\).
Enfin, l'acyclicité d'un complexe à deux termes de même rang
équivaut à l'inversibilité de \(D_E\otimes\mathbb Q\).
\end{proof}

Le lemme permet d'appliquer la forme normale de Smith sans présupposer aucun de
ses diviseurs. Il existe des changements de base orientés \(U,V\) tels que
\begin{equation}\label{eq:amp-bsd-morfismo-smith}
 U D_E V=\operatorname{diag}(s_1,\ldots,s_t),
 \qquad s_i\in\mathbb Z\setminus\{0\}.
\end{equation}
Par conséquent,
\begin{equation}\label{eq:amp-bsd-grupo-smith}
 F_E:=\operatorname{coker}D_E
 \simeq\bigoplus_{i=1}^{t}\mathbb Z/|s_i|\mathbb Z,
 \qquad
 |F_E|=\prod_{i=1}^{t}|s_i|<\infty.
\end{equation}

\begin{lemma}[Suite finie de localisation]
\label{lem:amp-bsd-exactitud-sha-tamagawa}
Le fragment de torsion de la suite longue de cohomologie de
\eqref{eq:amp-bsd-triangulo-localizacion} est
\begin{equation}\label{eq:amp-bsd-exacta-sha-tamagawa}
 0\longrightarrow\operatorname{Sha}(E)
 \xrightarrow{\ \iota_{\mathrm{Sha}}\ }F_E
 \xrightarrow{\ \partial_{\mathrm{loc}}\ }
 \bigoplus_v\Phi_v(\mathbb F_v)
 \longrightarrow0,
 \qquad c_v=|\Phi_v(\mathbb F_v)|.
\end{equation}
\end{lemma}

\begin{proof}
Un cocycle global représente une classe de
\(\operatorname{Sha}(E)\) précisément lorsque sa restriction à chaque place
admet une primitive dans
\(\mathcal C_{E,v}^{\mathrm{fin}}\). La classe de la paire formée par ce
cocycle et ses primitives locales dans le cône de
\eqref{eq:amp-bsd-mapa-localizacion} définit
\(\iota_{\mathrm{Sha}}\). Si cette classe est un bord, la composante globale est un
bord et la classe de Tate--Shafarevich est nulle ; ainsi,
\(\iota_{\mathrm{Sha}}\) est injective.

L'application \(\partial_{\mathrm{loc}}\) prend une classe du cône, la restreint à
chaque quotient
\(\mathcal C_{E,v}^{\mathrm{sing}}\) de
\eqref{eq:amp-bsd-triangulo-local} et conserve sa classe de composante.
L'identification
\[
 H^0(\mathcal C_{E,v}^{\mathrm{sing}})_{\mathrm{tors}}
 \simeq\Phi_v(\mathbb F_v)
\]
est induite par le quotient fini--singulier, non par l'ordre du groupe.
Une classe du cône a un résidu nul exactement lorsque ses représentants
locaux peuvent être choisis dans les sous-complexes finis ; d'après la définition
précédente, ce noyau est l'image de \(\iota_{\mathrm{Sha}}\).
Enfin, la dernière flèche de
\eqref{eq:amp-bsd-triangulo-local} relève chaque classe de composante au
complexe local. L'exactitude du cône
\eqref{eq:amp-bsd-triangulo-localizacion}, après retrait des deux
facteurs libres appariés par \(X^{\mathrm{orth}}\), transforme ce
relèvement en une préimage par \(\partial_{\mathrm{loc}}\). Cela démontre
la surjectivité et, par conséquent, toute
\eqref{eq:amp-bsd-exacta-sha-tamagawa}.
\end{proof}

Maintenant, et seulement après le lemme
\ref{lem:amp-bsd-aciclicidad-rango}, le groupe \(F_E\) est fini.
L'injection de \eqref{eq:amp-bsd-exacta-sha-tamagawa} démontre alors la
finitude de \(\operatorname{Sha}(E)\). En prenant les ordres dans cette suite
exacte, on obtient
\begin{equation}\label{eq:amp-bsd-cardinal-smith}
 |F_E|=|\operatorname{Sha}(E)|\prod_v c_v.
\end{equation}
C'est le comparateur
\(\mathfrak c_{\mathrm{STS}}\) : il transporte les facteurs élémentaires de Smith
vers les facteurs de Tate--Shafarevich et de Tamagawa en conservant leur ordre dans la droite,
au lieu d'insérer leurs cardinaux une fois le terme principal calculé.

La partie de torsion se construit en appliquant le déterminant au réseau de
Mordell--Weil et à son dual de Pontryagin. Les deux facteurs finis sont duaux et
apportent
\begin{equation}\label{eq:amp-bsd-comparador-torsion}
 \mathfrak c_{\mathrm{tors}}:quad
 u\longmapsto
 \frac{u}{|E(\mathbb Q)_{\mathrm{tors}}|^2}.
\end{equation}
Le comparateur archimédien s'obtient en intégrant la différentielle de Néron
orientée sur le cycle réel basé :
\begin{equation}\label{eq:amp-bsd-comparador-periodo}
 \Omega_E:=\int_{E(\mathbb R)}|\omega_E|,
 \qquad
 \mathfrak c_\infty:quad u\longmapsto\frac{u}{\Omega_E}.
\end{equation}
La différentielle et le cycle sont fixés avant d'évaluer \(L(E,s)\) ; ainsi
\eqref{eq:amp-bsd-comparador-periodo} ne normalise pas rétrospectivement la
section analytique.

\subsection{Séparation causale des primitives, des identités et des conséquences}
\label{subsec:amp-bsd-separacion-causal}

Le critère générateur par générateur du
lemme \ref{lem:amp-bsd-bases-emparejadas-kato-fers} audite à chaque degré
l'équivalence déjà construite. Il ne conditionne pas son existence. L'ordre logique
du paquet est
\begin{equation}\label{eq:amp-bsd-diagrama-causal}
\begin{array}{c}
 \begin{gathered}
  1_E,\ \Delta_{\mathrm{AW}},\ P_E^K,\ P_E^{\mathrm{PT}},
  \ h_*=-vf,\ X^{\mathrm{orth}},\\
  \{\,\mathcal C_{E,v}^{\mathrm{fin}}\to\mathcal C_{E,v}\,\}_v,\\
  \text{bases de Mordell--Weil et de Pontryagin ;}\\
  \text{différentielle et cycle de Néron}
 \end{gathered}
 \\[2mm]\Downarrow\\[-1mm]
 \begin{gathered}
  \Phi_{K/\mathrm{FERS}}\text{ équivalence},\quad
  \rho_{E,\ell}(0)=1,\quad
  p_{\mathrm{root}}z_K=z_{\mathrm{PT}},\\
  R_{\mathrm{Boc}}^{\mathrm{der}}=\operatorname{LT}_T(S_E),\quad
  \det(j_q\bar\beta_q)=\det h_E^{(q)},\\
  H^*(\mathcal C_E^{\mathrm{loc,red}}\otimes\mathbb Q)=0,\quad
  D_E\otimes\mathbb Q\text{ isomorphisme},\\
  0\to\operatorname{Sha}(E)\to F_E
  \to\bigoplus_v\Phi_v(\mathbb F_v)\to0
 \end{gathered}
 \\[2mm]\Downarrow\\[-1mm]
 \begin{gathered}
  d_{\mathrm{an}}=\operatorname{rank}E(\mathbb Q),\qquad
  \operatorname{Reg}(E),\qquad
  |\operatorname{Sha}(E)|\prod_vc_v,\\
  |E(\mathbb Q)_{\mathrm{tors}}|^{-2},\qquad
  \Omega_E^{-1},\qquad
  \text{égalité du terme principal}.
 \end{gathered}
\end{array}
\end{equation}
La première ligne contient uniquement des objets et des applications déjà basés, avec
la vérification générateur par générateur indiquée ; la deuxième, des identités
démontrées par ces applications ; la troisième, leurs conséquences sur la droite
déterminante. En particulier, aucun facteur de la dernière ligne n'est
employé pour normaliser rétroactivement une flèche de la première.

\begin{proposition}[Paquet exact de comparateurs]
\label{prop:amp-bsd-paquete-comparadores-exacto}
Les comparateurs de
\eqref{eq:amp-bsd-cadena-lineas} proviennent d'un unique paquet d'applications de
complexes et satisfont, dans l'ordre de construction,
\begin{equation}\label{eq:amp-bsd-paquete-identidades}
 \begin{gathered}
 \Phi_{K/\mathrm{FERS}}(z_K)=z_{\mathrm{Iw}},\qquad
 \operatorname{gr}^{F}\Phi_{K/\mathrm{FERS}}=I,\\
 \det\nolimits_{\mathrm{alt}}\Phi_{K/\mathrm{FERS}}
   =\rho_{E,\ell}(T),\qquad \rho_{E,\ell}(0)=1,\\
 \det(j_q\bar\beta_q)=\det h_E^{(q)}\quad(q\geq1),\\
 \Psi_{\mathrm{PT}}:
  \bigl(E(\mathbb Q)/E(\mathbb Q)_{\mathrm{tors}}\bigr)\otimes K
  \xrightarrow{\sim}\operatorname{Flag}_{\mathrm{Boc}}(E),\\
 p_{\mathrm{root}}z_K=z_{\mathrm{PT}},\qquad
 h_*=-vf,\quad \partial h_*=z_{\mathrm{PT}}-z_K,\\
 H^*(\mathcal C_E^{\mathrm{loc,red}}\otimes\mathbb Q)=0,\qquad
 U D_E V=\operatorname{diag}(s_1,\ldots,s_t),\\
 0\longrightarrow\operatorname{Sha}(E)
  \longrightarrow F_E
  \longrightarrow\displaystyle\bigoplus_v\Phi_v(\mathbb F_v)
  \longrightarrow0 .
 \end{gathered}
\end{equation}
La composition induite sur les droites déterminantes est exactement
\begin{equation}\label{eq:amp-bsd-paquete-composicion}
 \det(\Phi_{K/\mathrm{FERS}})
 \xrightarrow{\ \det(j_q\bar\beta_q)\ }
 \det(D_E)
 \xrightarrow{\ \mathrm{Smith}\ }
 \frac{|\operatorname{Sha}(E)|\prod_vc_v}
      {|E(\mathbb Q)_{\mathrm{tors}}|^2\,\Omega_E},
\end{equation}
avec les signes de Knudsen--Mumford et les orientations de
\eqref{eq:amp-bsd-cadena-lineas}. En particulier, le membre droit de
\eqref{eq:amp-bsd-paquete-composicion} est la coordonnée finale de l'élément
transporté et non une donnée servant à construire l'une des flèches.
\end{proposition}

\begin{proof}
La première ligne de \eqref{eq:amp-bsd-paquete-identidades} est
\eqref{eq:amp-bsd-phi-kato-fers} et
\eqref{eq:amp-bsd-matriz-kato-fers}--
\eqref{eq:amp-bsd-rho-determinantal}. La deuxième est
\eqref{eq:amp-bsd-comparador-pt-q} et le lemme
\ref{lem:amp-bsd-grado-pt}. La troisième combine l'égalité de la racine avec
le remplisseur explicite \eqref{eq:bsd-suc-hstar} ; son extension linéaire dans la
forme normale induit le lecteur auxiliaire utilisé dans le lemme
\ref{lem:amp-bsd-aciclicidad-rango}. En se restreignant au complément de la
racine, cette réduction contracte tout cycle rationnel. La dernière ligne est
\eqref{eq:amp-bsd-morfismo-smith} et
\eqref{eq:amp-bsd-exacta-sha-tamagawa}. Appliquer le foncteur déterminant à
ces identités, dans ce même ordre, donne
\eqref{eq:amp-bsd-paquete-composicion}. Aucune égalité ne s'obtient
en comparant d'abord les deux scalaires finaux.
\end{proof}

\begin{lemma}[Indépendance causale du terme principal]
\label{lem:amp-bsd-independencia-termino-principal}
Soit \(\mathscr P_E\) l'ensemble de données
\begin{equation}\label{eq:amp-bsd-primitivas-paquete}
 \mathscr P_E:=
 \left(
  1_E,\Delta_{\mathrm{AW}},P_E^K,P_E^{\mathrm{PT}},
  h_*=-vf,X^{\mathrm{orth}},
  \{\mathcal C_{E,v}^{\mathrm{fin}}\xrightarrow{i_v}
     \mathcal C_{E,v}\}_v,
  \text{bases et orientations},
  \omega_E,\gamma_E^{\mathbb R}
 \right),
\end{equation}
où \(\omega_E\) est la différentielle de Néron et
\(\gamma_E^{\mathbb R}\) le cycle réel orienté. Toutes les applications du
paquet de la proposition
\ref{prop:amp-bsd-paquete-comparadores-exacto} sont déterminées par
\(\mathscr P_E\). Ne font pas partie de \(\mathscr P_E\) les numéraux
\[
 L^{(r)}(E,1),\quad r,\quad\operatorname{Reg}(E),\quad
 |\operatorname{Sha}(E)|,\quad(c_v)_v,\quad
 |E(\mathbb Q)_{\mathrm{tors}}|,\quad\Omega_E.
\]
Ces valeurs s'obtiennent, respectivement, en évaluant la section analytique,
en prenant le degré du drapeau, en calculant des déterminants et des formes de Smith,
en prenant les ordres de groupes finis et en intégrant
\(\omega_E\) sur \(\gamma_E^{\mathbb R}\).
\end{lemma}

\begin{proof}
La formule \eqref{eq:amp-bsd-phi-kato-fers} utilise uniquement
\(1_E,\Delta_{\mathrm{AW}},P_E^K\) et la retenue. Les applications de Bockstein et de
Poitou--Tate utilisent la filtration, \(P_E^{\mathrm{PT}}\), les bases et
l'autodualité. Le cône de \eqref{eq:amp-bsd-mapa-localizacion} utilise les
applications \(i_v\) et les restrictions locales ; sa contraction et sa forme de
Smith utilisent le remplisseur explicite \eqref{eq:bsd-suc-hstar}, son extension
linéaire auxiliaire, \(X^{\mathrm{orth}}\) et les bases entières. Les deux
derniers comparateurs utilisent les réseaux de torsion et la
paire \((\omega_E,\gamma_E^{\mathbb R})\). Il s'agit d'une énumération exhaustive
des arguments qui figurent dans les formules de construction. Les
numéraux énumérés dans l'énoncé n'apparaissent qu'après application de la
cohomologie, du déterminant, du cardinal ou de l'intégration, de sorte qu'ils ne peuvent pas
sélectionner rétrospectivement les applications.
\end{proof}

\subsection{Composition, degrés et terme principal}
\label{subsec:amp-bsd-composicion}

Définissons le comparateur total par la composition dans l'ordre écrit :
\begin{equation}\label{eq:amp-bsd-comparador-total}
 \mathfrak C_E:=
 \mathfrak c_\infty\circ
 \mathfrak c_{\mathrm{tors}}\circ
 \mathfrak c_{\mathrm{STS}}\circ
 \mathfrak c_{\mathrm{PT}}\circ
 \mathfrak c_{K/\mathrm{FERS}}.
\end{equation}
Comme \(p_{\mathrm{root}}z_K=z_{\mathrm{PT}}\), toutes les flèches reçoivent la
même unité basée. Le remplisseur \(h_*=-vf\) de
\eqref{eq:bsd-suc-hstar} remplit explicitement la différence ; son extension
linéaire dans la base normale produit le lecteur homotopique auxiliaire et démontre
qu'aucune classe supplémentaire ne modifie l'élément de la droite pendant le
transport. Le lecteur n'est pas une prémisse distincte.

\begin{theorem}[Construction des comparateurs basés]
\label{thm:amp-bsd-comparadores-construidos}
Pour toute courbe elliptique \(E/\mathbb Q\) munie des complexes, des bases,
des orientations et des données locales définis précédemment, les applications
\eqref{eq:amp-bsd-phi-kato-fers},
\eqref{eq:amp-bsd-comparador-pt},
\eqref{eq:amp-bsd-triangulo-localizacion},
\eqref{eq:amp-bsd-comparador-torsion} et
\eqref{eq:amp-bsd-comparador-periodo}, avec les lemmes
\ref{lem:amp-bsd-aciclicidad-rango} et
\ref{lem:amp-bsd-exactitud-sha-tamagawa}, construisent les cinq comparateurs de
\eqref{eq:amp-bsd-cadena-lineas}. Leur composition préserve l'unité, le degré
et l'orientation, et satisfait
\begin{equation}\label{eq:amp-bsd-igualdad-grados}
 d_{\mathrm{an}}=\operatorname{ord}_T\det S_E(T)
 =d_{\mathrm{ar}}=\operatorname{rank}E(\mathbb Q).
\end{equation}
De plus, \(\operatorname{Sha}(E)\) est fini et l'égalité de l'élément
distingué dans \(\mathcal L_E\) se publie sous la forme
\begin{equation}\label{eq:amp-bsd-termino-principal}
 \frac{L^{(r)}(E,1)}{r!\,\Omega_E}
 =\frac{|\operatorname{Sha}(E)|\,\operatorname{Reg}(E)
        \prod_v c_v}
       {|E(\mathbb Q)_{\mathrm{tors}}|^2},
 \qquad r=\operatorname{rank}E(\mathbb Q).
\end{equation}
\end{theorem}

\begin{proof}
Le comparateur Kato--FERS et la condition
\(\rho_{E,\ell}(0)=1\) identifient l'ordre analytique à
\(d=\operatorname{ord}_T\det S_E(T)\) et le terme initial à
\(R_{\mathrm{Boc}}^{\mathrm{der}}(S_E)\). L'identité
\eqref{eq:bsd-suc-lt} conserve la page régulière et chaque page positive. Les
isomorphismes \(j_q\) transportent ces pages vers les hauteurs dérivées. Le
lemme \ref{lem:amp-bsd-grado-pt} construit l'isomorphisme du drapeau avec le
réseau libre de Mordell--Weil et, avec
\eqref{eq:bsd-suc-carry-orden}, démontre
\eqref{eq:amp-bsd-igualdad-grados}.

Le lemme \ref{lem:amp-bsd-aciclicidad-rango} déduit l'acyclicité rationnelle
directement du remplisseur explicite \eqref{eq:bsd-suc-hstar}, de sa réduction
linéaire du complément, de l'annulation de la racine commune et de
\(X^{\mathrm{orth}}\) ; on en obtient le rang plein de \(D_E\), non
comme hypothèse supplémentaire. La forme de Smith
\eqref{eq:amp-bsd-morfismo-smith} produit alors le groupe fini \(F_E\).
Le lemme \ref{lem:amp-bsd-exactitud-sha-tamagawa} extrait des triangles
\eqref{eq:amp-bsd-triangulo-local} et
\eqref{eq:amp-bsd-triangulo-localizacion} la suite exacte finie ; ce n'est
qu'ensuite que l'on déduit la finitude de \(\operatorname{Sha}(E)\) et
\eqref{eq:amp-bsd-cardinal-smith}. Enfin,
\eqref{eq:amp-bsd-regulador-nt},
\eqref{eq:amp-bsd-comparador-torsion} et
\eqref{eq:amp-bsd-comparador-periodo} publient, respectivement, le régulateur,
le quotient de torsion et la période. L'égalité racine élimine toute unité
relative entre les deux publications. En évaluant le même élément basé aux
deux extrémités de \eqref{eq:amp-bsd-comparador-total}, on obtient
\eqref{eq:amp-bsd-termino-principal}.
\end{proof}

\subsection{Prémisses et critères de réfutation}
\label{subsec:amp-bsd-falsadores}

La construction part de l'unité généalogique \(1_E\) et de la diagonale
d'Alexander--Whitney. Elle utilise également les
publications couplées, tous les
Bockstein et la page régulière, le remplisseur explicite \(h_*=-vf\),
l'autodualité réduite \(X^{\mathrm{orth}}\), les applications de chaînes du
triangle de localisation et les bases orientées de torsion et de période. L'équivalence
Kato--FERS et sa normalisation à un, le rang plein de \(D_E\),
l'exactitude de \eqref{eq:amp-bsd-exacta-sha-tamagawa} et la finitude de
\(\operatorname{Sha}(E)\) sont des conclusions des lemmes précédents, non des
prémisses indépendantes. Le critère générateur par générateur du lemme
\ref{lem:amp-bsd-bases-emparejadas-kato-fers} est un contrôle non directeur
qui audite localement le premier comparateur. Les critères qui réfuteraient une flèche précise
sont :
\begin{enumerate}
 \item un coefficient racine différent de un, qui introduirait une unité
       relative entre \(z_K\) et \(z_{\mathrm{PT}}\) ;
 \item un \(j_q\) non inversible ou l'omission d'une page de Bockstein, qui
       modifierait le degré ou le déterminant de hauteur ;
 \item un diviseur de Smith nul, qui empêcherait la finitude de \(F_E\) ;
 \item l'échec de l'exactitude de
       \eqref{eq:amp-bsd-exacta-sha-tamagawa}, qui romprait l'identification
       des facteurs finis ;
 \item une inversion d'orientation dans la torsion ou la période, qui changerait
       l'élément basé même si elle préservait sa droite non orientée.
\end{enumerate}
La projection de Manin sans histoire et le changement d'échelle visible modulo trois
violent le premier critère ; ils constituent donc des preuves de nécessité pour
ces modèles réduits et non des prémisses du théorème
\ref{thm:amp-bsd-comparadores-construidos}.


\begin{theorem}[Comparaison déterminantale locale--globale]
\label{thm:bsd-suc-local-global}
Soient les comparateurs basés déjà construits dans le
\cref{thm:amp-bsd-comparadores-construidos} :
Kato--FERS/Iwasawa,
Poitou--Tate/Poincaré--Néron--Tate,
Smith--Tamagawa--Shafarevich, torsion et période.
Si \(d_{\mathrm{an}}\) est le degré de la section analytique dans la
spécialisation centrale et \(d_{\mathrm{ar}}\) le degré de la droite de Selmer
basée, alors
\begin{equation}\label{eq:bsd-suc-grados}
 d_{\mathrm{an}}=d=d_{\mathrm{ar}},
\end{equation}
et l'on obtient l'égalité d'éléments distingués
\begin{equation}\label{eq:bsd-suc-linea-final}
 \mathfrak z_{\mathrm{an}}(E)=\mathfrak z_{\mathrm{ar}}(E)
 \quad\text{dans}\quad\mathcal L_E.
\end{equation}
Le membre analytique est la publication du terme initial ; le membre
arithmétique est la publication du régulateur complet, des facteurs
élémentaires de Smith,
des facteurs locaux, de la torsion et de la période.
\end{theorem}

\begin{proof}
La construction basée donne
\([\Phi_{K/\mathrm{FERS}}^i](T)=I+TB_i(T)\) en chaque degré. Le comparateur
basé de Kato--FERS exprime alors la section analytique corrigée
comme
\begin{equation}\label{eq:bsd-suc-fers}
 L_{\ell}^{\mathrm{corr}}(E,T)
 =\rho_{E,\ell}(T)\det S_E^{\mathrm{corr}}(T),
 \qquad \rho_{E,\ell}(0)=1.
\end{equation}
Par \eqref{eq:bsd-suc-orden}, son degré central est \(d\) ; par
\eqref{eq:bsd-suc-lt}, son terme initial est exactement
\(R_{\mathrm{Boc}}^{\mathrm{der}}\), et non ce régulateur multiplié par une
unité indéterminée. Les isomorphismes \(j_q\) de
\eqref{eq:bsd-suc-jq} transportent chaque \(\tau_q\) vers le déterminant de la
hauteur dérivée de sa page. La page régulière transporte le terme
acyclique que les pages positives ne voient pas.

La comparaison de Poincaré--Néron--Tate identifie le degré de cette droite de
hauteurs à \(d_{\mathrm{ar}}\). Les facteurs élémentaires de Smith décomposent sa partie
finie en facteurs de Shafarevich et de Tamagawa ; le quotient des réseaux
entiers apporte le carré de la torsion, et le lecteur archimédien apporte
la période. Toutes ces applications agissent sur la même droite racine. Par le
théorème \ref{thm:bsd-suc-raiz-filler}, il ne reste pas d'unité relative entre
les publications Kato et Poitou--Tate. La composition préserve le degré et
l'élément basé, ce qui prouve simultanément
\eqref{eq:bsd-suc-grados} et \eqref{eq:bsd-suc-linea-final}. La formule
scalaire correspondante est réservée à la reconnaissance finale.
\end{proof}

\subsection*{Hypothèses et portée}

La chaîne BSD utilise : l'état enrichi complet et sa coalgèbre
d'Alexander--Whitney ; le produit fibré \eqref{eq:bsd-suc-pullback} ; les deux
publications construites sur la même unité ; la droite racine, sa base et sa
counité ; le quotient réduit parfait et saturé, avec
\(X^{\mathrm{orth}}\) ; le complexe fini--singulier complet ; les drapeaux
de Bockstein finis, leurs signes de Koszul et la page régulière ; et les
comparaisons de hauteur, de Smith, de Tamagawa, de Shafarevich, de torsion et de périodes
avec une orientation commune. La vérification degré par degré des bases
appariées et des résiduels \(\varepsilon_i(b)\) est enregistrée comme
contrôle local de la réalisation, non comme prémisse gouvernante de la fermeture.
Sous cette structure de départ explicite,
\eqref{eq:bsd-suc-raiz}, \eqref{eq:bsd-suc-hstar},
\eqref{eq:bsd-suc-lt} et \eqref{eq:bsd-suc-linea-final} sont des égalités
exactes.

La conservation de l'unité, le parcours à deux feuillets et la publication
locale--globale sont formalisés au moyen de la coalgèbre de chaînes, du produit
fibré généalogique, de la publication couplée, de la réduction orthogonale et de la
tour de Bockstein. Les identités auxiliaires vérifient séparément la propriété d'élément groupal,
la conservation de l'histoire sous \(\widehat\Pi_0\), l'identité de
l'unité dans les deux canaux, l'involution réduite, la forme normale
finie--singulière et l'identité déterminantale d'ordre arbitraire.

\paragraph{Contrôles d'ablation BSD.}
\noindent\textbf{Statut : CONTRÔLE\_NON\_DIRECTEUR.}
Les trois ablations suivantes reçoivent les identités déjà démontrées et
réfutent uniquement des modèles réduits. Elles n'apportent pas de prémisses au propriétaire,
ne conditionnent pas sa fermeture et ne possèdent pas d'arête causale de retour vers lui.

\begin{ablation}[Projection de Manin sans histoire]
\label{abl:bsd-suc-manin}
Remplacer \(P_E^{\mathrm{gen}}\) par \(P_E^{\mathrm{Man}}\) efface des données dont
la publication couplée a besoin. En effet, la voie
\(\gamma=(\mathsf N\mathsf E)^{54}\) peut avoir la même cellule et la même
phase visible que la voie triviale, mais conserve un enroulement et une mémoire
distincts dans \(G_{m,n}\). L'augmentation de Manin les identifie ; la seconde
projection de \eqref{eq:bsd-suc-pullback} ne le fait pas.
\end{ablation}

\begin{proof}
Tout relèvement basé au complexe de Manin se factorise, à homotopie près, par la
section de l'augmentation \(s\varepsilon_{\mathrm{AW}}\). Il ne dépend donc que
de la coordonnée visible. En revanche, la seconde composante de
\(\widehat\Pi_0(c)\) est \(c\) lui-même, de sorte qu'elle conserve le simplexe
d'histoire et distingue l'enroulement de \(\gamma\). L'ablation réfute la
projection nue ; elle ne réfute pas le produit fibré employé dans le théorème.
\end{proof}

\begin{ablation}[Changement d'échelle de la racine visible seulement modulo trois]
\label{abl:bsd-suc-twist}
Dans la spécialisation \(c=1\) de
\eqref{eq:bsd-suc-diferencial-normal}, le cycle altéré
\begin{equation}\label{eq:bsd-suc-twist4}
 z_K^{(4)}=4r+3e+b
\end{equation}
satisfait
\begin{equation}\label{eq:bsd-suc-twist-defecto}
 z_{\mathrm{PT}}-z_K^{(4)}
 =\partial(-f)-3r.
\end{equation}
Le défaut \(-3r\) disparaît modulo trois et réapparaît modulo neuf. Par
conséquent, \eqref{eq:bsd-suc-twist4} conserve une ombre partielle, mais non
l'unité basée ni \eqref{eq:bsd-suc-raiz}.
\end{ablation}

\begin{proof}
L'égalité \(\partial f=3e+b\) donne immédiatement
\eqref{eq:bsd-suc-twist-defecto}. Modulo trois, le second terme est nul ;
modulo neuf, c'est une classe racine non divisible par un bord du complément.
En outre, sa coordonnée racine est quatre, tandis que
\eqref{eq:bsd-suc-normalizacion-raiz} exige un. L'exemple vérifie la
nécessité du couple couplé et de l'orientation ; il n'introduit pas d'exception
au théorème qui conserve les deux données.
\end{proof}

\begin{ablation}[Perte de pages du régulateur]
\label{abl:bsd-suc-bockstein}
La matrice \(\operatorname{diag}(T^2,T)\) montre que le premier Bockstein ne
retrouve pas la retenue d'ordre supérieur. La matrice
\(\operatorname{diag}(2,T^2)\) possède une trivialisation positive
\(\tau_2=1\), mais
\(\tau_{\mathrm{reg}}=2=\operatorname{LT}_{T}(S)\). Par conséquent, ni la
première page ni le produit des pages positives ne remplacent le régulateur
complet \eqref{eq:bsd-suc-rboc}.
\end{ablation}

\begin{proof}
Dans le premier exemple, \(d=3\), \(r_0=2\) et
\(c_I^{\deg}=1\) ; la page un ne contient pas ce dernier degré de mémoire. Dans
le second, l'unique diviseur singulier est \(T^2\), mais la direction régulière
apporte le facteur deux. L'omettre produirait un au lieu de deux. Les deux calculs
sont des cas diagonaux de \eqref{eq:bsd-suc-smith} et vérifient la nécessité de
toutes les pièces de \eqref{eq:bsd-suc-rboc}.
\end{proof}
\section*{Reconnaissance mathématique finale de Birch--Swinnerton--Dyer}

Le développement précédent s'achève avant d'employer sa formulation historique
comme critère.

Dans le domaine déclaré par le théorème
\ref{thm:bsd-suc-local-global}, la comparaison de Smith identifie la
composante finie de la droite arithmétique à la décomposition
Shafarevich--Tamagawa et démontre la finitude de
\(\operatorname{Sha}(E)\). En conséquence, son ordre dans la publication
scalaire suivante est le cardinal du groupe fini identifié, non un symbole
formel ajouté a posteriori.

La traduction scalaire de \eqref{eq:bsd-suc-linea-final} fournit, avec
\(r=\operatorname{rank}E(\mathbb Q)\),
\begin{equation}\label{eq:bsd-suc-rango-final}
 \operatorname{ord}_{s=1}L(E,s)=r
\end{equation}
et
\begin{equation}\label{eq:bsd-suc-formula-final}
 \frac{L^{(r)}(E,1)}{r!\,\Omega_E}
 =\frac{
   \lvert\operatorname{Sha}(E)\rvert\,\operatorname{Reg}(E)\prod_v c_v}
  {\lvert E(\mathbb Q)_{\mathrm{tors}}\rvert^2}.
\end{equation}
C'est la formulation conventionnelle de Birch--Swinnerton--Dyer. L'ordre,
le terme principal et les facteurs arithmétiques ne sont pas définis les uns au moyen des
autres : ce sont des publications coordonnées de l'égalité basée
\eqref{eq:bsd-suc-linea-final}, dont la racine commune a été fixée au préalable par
\eqref{eq:bsd-suc-raiz}.

